\documentclass{umbruch-position}
\usepackage{amsmath}
\usepackage{amssymb}
\usepackage{tikz}
\usetikzlibrary{arrows.meta,positioning,shapes.geometric,fit,calc}

\positionsnummer{ST-001}
\dokumenttyp{Studie}
\title{Systemtheoretische Aufarbeitung der Regressangst}
\untertitel{Eine multiperspektivische Bestandsaufnahme der deutschen\\Wirtschaftlichkeitsprüfung -- Begleitstudie zu PP-003}
\author{Um:bruch -- Denkfabrik für gesellschaftlichen Wandel}
\date{April 2026}
\versionsnummer{0.12 (Arbeitsversion)}

\zusammenfassung{%
Diese Studie nimmt das deutsche Wirtschaftlichkeitsprüfsystem nach \S\S\,106\,ff. SGB\,V multiperspektivisch in den Blick: historisch (Container-Genealogie seit 1989), rechtsdogmatisch (Selbstverständnis der Prüfverfahren, Zertifikatsfehler), verhaltensökonomisch (Ribbat 2023, Goetz 2024), gesundheitsökonomisch (Folgekostenrechnung) und international vergleichend. Zentraler neuer Befund: Das Regresssystem leidet an einem \textit{doppelten Zertifikatsfehler} --- die Auffälligkeitsprüfung misst Kostenabweichungen (enthält aber ein Schutzstufenmodell mit medizinischer Berücksichtigung), die Einzelfallprüfung misst überwiegend formale Angreifbarkeit, wobei ${\sim}$35\,\% der Prüfungen evidenzbasierte Steuerung codieren (AM-RL, G-BA-Beschlüsse). Der Zertifikatsfehler gilt in voller Schärfe nur für V2b (Formfehlerregresse). Ein zweiter Befund betrifft die Regressstatistik: Die öffentlich kommunizierte Quote von ``unter 1\,\%'' und die von Ärzten berichtete 72\,\%-Erfahrungsrate (Ribbat 2023) sind beide korrekt --- sie messen verschiedene Dinge (Definitionsdivergenz, Faktor ${\sim}$500$\times$). Die Studie ist die wissenschaftliche Begleituntersuchung zum Konzeptpapier \textit{PP-003 -- Regress-Transparenzportal}.
}
\schlagworte{Regress, Wirtschaftlichkeitsprüfung, Defensivmedizin, Versorgungsforschung, Gesundheitsökonomie, Systemanalyse}

\begin{document}
\maketitle
\newpage
\tableofcontents
\newpage


\section{Die Diagnose}

Diese Studie nimmt das deutsche Wirtschaftlichkeitsprüfsystem nach \S\S\,106\,ff. SGB\,V multiperspektivisch in den Blick. Sie ist eine Bestandsaufnahme: Was lässt sich über das System belastbar sagen, wenn man die verfügbaren empirischen, rechtsdogmatischen, historischen und vergleichenden Befunde nebeneinander legt?

Den Anlass liefert eine auffällige Diskrepanz: Die Regressdurchsetzungsquote wird öffentlich als seltenes Ereignis kommuniziert --- unter 1\,\% der Ärzte (verschiedene KV-Statistiken, Auffälligkeitsprüfung und Einzelfallprüfung zusammen; zur Differenzierung siehe Abschnitt~2.1). Dennoch berichten 72\,\% der Hausärzte, tatsächlich mindestens einmal einen Regress erlebt zu haben, und 47\,\% geben an, die Regressgefahr beschäftige ihren Praxisalltag ``stark oder sehr stark'' (Ribbat et al. 2023, n=770, TU München; V1+V2 zusammen, keine Verfahrenstrennung in der Studie). 46,7\,\% der ca. 13.900 vom KBV-Berufsmonitoring (KBV/Universität Trier 2018, dritte Erhebungswelle) befragten Medizinstudierenden nennen die Regressdrohung als Niederlassungshindernis --- obwohl sie selbst noch nie niedergelassen waren. Sind Ärzte nur besonders ängstlich --- oder hat ihre Angst einen realeren Hintergrund als die offizielle Statistik erkennen lässt? Ein Teil der Auflösung vorweg: Die ``unter 1\,\%'' sind eine Quote \textit{pro Arzt und Jahr}, dominiert durch das V1-Screening. Bei der Einzelfallprüfung (V2) ist das Risiko \textit{pro Verordnung} gering --- aber ein Arzt mit 3.000 Verordnungen pro Quartal akkumuliert über 10 Jahre 120.000 Angriffspunkte. In KV\,Nordrhein erreicht statistisch jede Praxis alle 2--3 Jahre mindestens eine V2-Prüfung. Kleines Risiko pro Fall, großes Risiko pro Arzt --- das erklärt die 72\,\%.

\subsection{Fragestellung}

Die Frage, die diese Diskrepanz aufwirft, lautet zunächst nicht, was ``richtig'' ist, sondern: \textbf{Was passiert im System, das beides zugleich produziert?} Diese Analyse versucht, Antworten zu sammeln --- aus historischen, rechtsdogmatischen, verhaltensökonomischen, gesundheitsökonomischen und international vergleichenden Quellen.

\textbf{Diese Studie ist die wissenschaftliche Begleituntersuchung} zum Konzeptpapier PP-003, das ein anonymes Regress-Transparenzportal zur empirischen Schließung einer Datenlücke vorschlägt. ST-001 dokumentiert, was sich derzeit über das System sagen lässt, und welche Datenlücken dabei sichtbar werden. PP-003 schlägt einen Weg vor, eine dieser Lücken zu schließen.

Diese Arbeitsversion ist bewusst zurückhaltend in den Schlussfolgerungen. Wo Wertungen vorgenommen werden, sind sie als solche kenntlich. Wo Befunde mehrere Lesarten zulassen, werden alle Lesarten dargestellt. Die Studie versucht nicht, eine vorgefasste These zu bestätigen, sondern sich von den verfügbaren Daten leiten zu lassen.

\subsection{Methodisches Vorgehen}

\subsubsection{Forschungsdesign}

Diese Studie entstand in einem mehrstufigen, KI-gestützten Analyseprozess im April 2026. Insgesamt wurden 31 eigenständige Recherchestränge durchgeführt, dokumentiert und in eine integrierte Argumentation überführt. Jeder Strang ist als eigenständige Recherchedatei im Arbeitsordner des Projekts dokumentiert und nachverfolgbar (Projektordner \texttt{Content/\allowbreak{}Projekte/\allowbreak{}Regress-Melder/\allowbreak{}infomaterial/}, Inventar in \texttt{WALK\_OF\_\allowbreak{}ANALYSIS.md}). Die vollständige Methodendokumentation aller Analyseschritte --- einschließlich der rekonstruierten Prompt-Kerne, Modellzuordnungen und Hypothesenrevisionen --- findet sich in Anhang~B (Abschnitt~\ref{sec:appendix-methodik}).

\subsubsection{Quellen und Suchstrategie}

Die Quellenerhebung erfolgte systematisch für den Zeitraum 2018--2026 und selektiv für ältere Quellen. Durchsuchte Datenbanken und Quellensysteme:

\begin{itemize}[nosep]
  \item \textbf{Wissenschaftliche Literatur:} PubMed, Google Scholar (Suchbegriffe: Wirtschaftlichkeitsprüfung, Regressangst, Defensivmedizin, prescribing behavior regulation, defensive medicine Germany)
  \item \textbf{Rechtsprechung:} BSG-Entscheidungsdatenbank (Senat B\,6\,KA), Landessozialgericht-Entscheidungen, dejure.org
  \item \textbf{Gesetzestexte und Materialien:} SGB\,V (laufende Fassung + historische Fassungen via BGBl), BMV-Ä, Prüfvereinbarungen ausgewählter KV-Regionen
  \item \textbf{Institutionelle Quellen:} KBV-Statistiken, GKV-Spitzenverband, WIdO-Arzneimittelreports, BRH-Sonderberichte, OECD Health at a Glance
  \item \textbf{Informationsfreiheitsanfragen:} 9 IFG-Anfragen an BMG, BAS, KBV und GPE über FragDenStaat (eingereicht 10.04.2026, Fristen laufend; dokumentiert in Anhang~C)
\end{itemize}

\subsubsection{Multi-Modell-Architektur}

Die Analyse wurde durch vier Instanzen arbeitsteilig durchgeführt:

\begin{tabular}{p{2.2cm}p{3.3cm}p{7cm}}
\textbf{Instanz} & \textbf{Modell/Version} & \textbf{Rolle} \\
\hline
Claude & Opus 4.6, 1M Kontext & Primäranalyse, Rechtsstruktur, Synthese, LaTeX-Dokumentation \\
Copilot & GPT-4o & Konzeptschärfung, Fehlkonzept-Erkennung, Kostenschätzungs-Review \\
Gemini & Deep Research & Theorierahmen (Freidson, Power, Lipsky), Quellenergänzung, Visualisierung \\
LG (Mensch) & --- & Fragestellungen, Richtungsentscheidungen, Domänenwissen, finale Freigabe \\
\end{tabular}

\vspace{0.3em}
\noindent Die Detailtabelle mit Stärken- und Schwächenprofilen jedes Modells findet sich in Anhang~B.

\subsubsection{Cross-Source-Divergenz (Multi-Agenten-Experiment)}

In einer späteren Analysephase (Phase~XIV) wurde die Multi-Modell-Architektur um eine \textbf{Multi-Agenten-Architektur} erweitert: 15 spezialisierte Agenten (6~Opus, 9~Sonnet) recherchierten parallel zu Regresshäufigkeit, Kassendaten und Evidenzsteuerung --- mit 25~Datenpunkten aus 14~unabhängigen Domains. \textbf{Zentrale methodische Lektion:} Cross-Model-Divergenz (Claude vs. Copilot vs. Gemini) kontrolliert Modell-Bias, aber \textit{nicht} Quellen-Bias. Alle Modelle schöpften aus denselben Quellen (Ärzteblatt, Virchowbund, KV~Nordrhein). Erst die systematische Multi-Quellen-Recherche brachte die SpiFa-Bayern-Daten (13.000~Regresse/Jahr) und die BMG-Zahl (47.000~V2-Verfahren) ans Licht. Zusätzlich wurden ein \textit{Naive-Suche-Test} (simulierter Journalist ohne Vorwissen) und zwei \textit{Prior-Art-Checks} (systematische Suche nach bestehender Synthese) als Pflichtschritte eingeführt (Details in Anhang~B).

\subsubsection{Ein- und Ausschlusskriterien}

\textbf{Eingeschlossen} wurden: peer-reviewed Studien, BSG-Entscheidungen, offizielle Statistiken (KBV, WIdO, Destatis), verifizierte Gesetzesmaterialien und IFG-Dokumente. Journalistische Quellen wurden nur als Ergänzung herangezogen, nicht als Primärbeleg.

\textbf{Ausgeschlossen} wurden: nicht-verifizierbare Blog-Beiträge, unattribuierte Statistiken, Gemini-Outputs ohne Claude-Double-Check (dokumentierte Halluzinationsfälle: Moosazadeh-Fehlzuordnung, französische Protestzahlen; siehe Anhang~B). Verbände-Stellungnahmen wurden als Positionspapiere gelesen, nicht als empirische Belege.

\subsubsection{Bias-Management}

Die methodische Kontrolle gegen Bestätigungsfehler erfolgte auf mehreren Ebenen:

\begin{itemize}[nosep]
  \item \textbf{Cross-Model-Divergenz:} Abweichungen zwischen Claude, Copilot und Gemini wurden systematisch als Signal genutzt. Copilot identifizierte eine Inkonsistenz in der Kostenschätzung, die zur Einführung des Dreistufenmodells führte.
  \item \textbf{Hypothesenrevisionen:} Fünf dokumentierte Fälle, in denen Ausgangshypothesen anhand der Evidenz verworfen wurden (u.\,a. Datierung der Einzelfallprüfung: GSG\,1992 $\rightarrow$ GRG\,1988; Moosazadeh-Fehlzuordnung $\rightarrow$ Zheng\,2023).
  \item \textbf{Devil's Advocate:} Ein eigener Review-Agent (Phase~10 des 13-Phasen-Review-Zyklus) versuchte systematisch, die zentrale These zu widerlegen. Die 5 identifizierten Verwundbarkeiten wurden in den Text eingearbeitet (siehe Abschnitt~\ref{sec:gegenposition}).
  \item \textbf{Unabhängiger Policy-Analyst:} Ein separater Agent ohne Kenntnis der Recherchegeschichte formulierte 7 Lösungsvorschläge, um Tunnelblick-Effekte zu umgehen.
\end{itemize}

\subsubsection{Methodische Grundsätze}

Die Studie macht folgende methodische Annahmen explizit:

\begin{itemize}[nosep]
  \item \textbf{Quellenpriorität:} Primärquellen (Gesetzestexte, BSG-Urteile, peer-reviewed Studien) vor Sekundärquellen (Stellungnahmen, Pressemitteilungen). Bei Halluzinationsrisiko: jede Quelle wurde gegen Originalpublikation geprüft.
  \item \textbf{Anonymisierung:} Alle individuellen Fälle sind anonymisiert. Keine namentliche Nennung. Keine identifizierenden Details.
  \item \textbf{Hochrechnungen mit offenen Annahmen:} Wo Schätzungen vorgenommen werden, sind die Annahmen explizit. Skepsis gegenüber Punktschätzungen, Bevorzugung von Bandbreiten.
  \item \textbf{Trennung von Befund und Bewertung:} Die normative Folgerung steht am Ende, nicht am Anfang. Wer den Befunden nicht folgt, soll an den Befunden argumentieren können.
\end{itemize}

\noindent Die Analyse gliedert sich in sieben Teile: historische Schichtung (I), rechtsdogmatische Analyse (II), Verhaltensökonomie (III), Folgekosten (IV), institutionelle Analyse (V), internationaler Vergleich (VI) und Lösungsraum (VIII). Zwischen dem internationalen Vergleich und dem Lösungsraum steht ein eigener Abschnitt zu den Grenzen der Evidenz und den stärksten Gegeneinwänden (VII).

\section{Teil I --- Historische Schichtung: Wie das System wurde, was es ist}

\subsection{Vom Solo-Container zur Containerlandschaft}

\textit{Container} meint in dieser Studie einen einzelnen Paragrafen oder eine zusammenhängende Norm, in der rechtliche Konzepte ``aufbewahrt'' werden. Der Begriff ist eine analytische Hilfe, um die Wanderung von Konzepten zwischen verschiedenen Normstandorten nachzuverfolgen.

1989 reichte ein einziger Paragraf des SGB\,V (\S\,106) für die ambulante ärztliche Wirtschaftlichkeitsprüfung. Eingeführt durch das Gesundheits-Reformgesetz (GRG) unter Norbert Blüm (CDU), in Kraft am 01.01.1989. Dieser Solo-Container enthielt \textbf{beide} Prüfverfahren von Anfang an:

\begin{itemize}[nosep]
  \item Die \textbf{Einzelfallprüfung} (V2) stammt aus dem früheren \S\,368n RVO, in der BSG-Recht\-sprechung seit 1959 (BSGE 11, 102) verankert. Sie ist nicht erst später eingeführt worden --- sie war von Anfang an dabei.
  \item Die \textbf{Auffälligkeitsprüfung} (V1) in Form der Richtgrößenprüfung wurde mit dem Gesundheitsstrukturgesetz (GSG) 1992/93 unter Horst Seehofer (CSU) als bundeseinheitliches statistisches Verfahren eingeführt. Zuvor existierte eine weniger formalisierte Durchschnittswertprüfung, die nach Fachgruppe verglich.
\end{itemize}

Was sich ändert, ist nicht das Grundprinzip (Kostenprüfung + Sanktion), sondern die Position der Verfahren innerhalb einer wachsenden Containerlandschaft.

\textbf{Was der Solo-Container 1992 außerdem enthielt:} Der originale \S\,106 SGB\,V (Fassung GRG 1989, geändert durch GSG 1992) war mehr als nur die beiden Prüfverfahren. Er regelte in einem einzigen Paragrafen: die Besetzung und Zuständigkeit der Prüfungsstelle, das Widerspruchsverfahren (Beschwerdeausschuss), die Beratungspflichten, die Sanktionsfolgen und die Kostentragung. Diese funktionalen Elemente sind 2017 auf sechs Paragrafen (\S\S\,106--106d) und die regionalen Prüfvereinbarungen verteilt worden --- bei gleichbleibendem Regelungsgehalt. Die Aufspaltung erhöhte die Regelfläche, nicht die Regelsubstanz.

2026 sind aus dem einen Container fünf ambulante (\S\S\,106, 106a, 106b, 106c, 106d), zwei stationäre (\S\S\,275, 275c) plus \S\,48 BMV-Ä, plus bundesweite Rahmenvorgaben von KBV und GKV-Spitzenverband, plus bis zu 17 regionale Prüfvereinbarungen geworden. Die geschriebene Regelfläche hat sich in 35 Jahren etwa verachtfacht. Der Regelungsgehalt ist im Kern derselbe geblieben --- Konzepte wurden aufgespalten, dupliziert, regional variiert und zwischen Normstandorten verschoben, aber der Grundmechanismus (statistische Auffälligkeit oder Kassenantrag $\to$ Prüfung $\to$ Sanktion) blieb unverändert.

\begin{figure}[htbp]
  \centering
  \includegraphics[width=0.95\textwidth]{C:/Users/User/OneDrive/.UMBRUCH/Content/Projekte/Regress-Melder/grafiken/container_genealogie_zeitachse}
  \caption{Container-Genealogie 1989--2026: Aus einem Paragrafen (\S\,106) wurden acht Container. Die große Reform 2015/2017 (GKV-VSG) ist der Wendepunkt. Quelle: eigene Darstellung.}
  \label{fig:container-genealogie}
\end{figure}

\subsection{Der Lahnsteiner Kompromiss 1992}

Im Sommer 1992 wechselte das Bundesgesundheitsministerium von Gerda Hasselfeldt (CSU) zu Horst Seehofer (CSU). Mit dem Gesundheitsstrukturgesetz (GSG) von 1992/93 (BT-Drs. 12/3608) wurde \S\,106 SGB\,V umfassend reformiert. Der Entwurf war das Ergebnis des sogenannten \textit{Lahnsteiner Kompromisses} zwischen Seehofer und dem SPD-Sozialpolitiker Rudolf Dreßler. Befürworter waren die Bundesregierung (CDU/CSU/FDP), die SPD-Opposition und alle Krankenkassen-Spitzenverbände (AOK-BV, BKK, VdAK, IKK). Kritiker (BÄK, KBV, Hartmannbund) richteten ihre Kritik primär gegen die Budgetierung und die Arzneimittelbudgets, nicht gegen die Einzelfallprüfung selbst. Die Einzelfallprüfung war damals strukturell unangefochten.

In dieser Reform wurde die Einzelfallprüfung in \S\,106 Abs.\,2 Nr.\,2 SGB\,V als bundeseinheitliche Ergänzungsmethode auf begründeten Antrag einer Krankenkasse oder einer Kassenärztlichen Vereinigung verankert. Sie war zu diesem Zeitpunkt eine Norm unter mehreren in einem einzigen Container.

\subsection{GKV-VSG 2015: Die zentrale Wendepunkt-Reform}

Der entscheidende Bruch erfolgte 2015 mit dem GKV-Versorgungsstärkungsgesetz (GKV-VSG) unter Hermann Gröhe (CDU). Die zentralen Operationen:

\begin{itemize}[nosep]
  \item \textbf{Abschaffung der bundesweit einheitlichen Richtgrößenprüfung} (in Kraft 01.01.2017). Die Richtgrößenprüfung setzte vorab definierte Euro-Beträge pro Patient und Quartal fest; Ärzte, die diese überschritten, wurden automatisch geprüft. Begründung der Abschaffung im Gesetzgebungsverfahren: mangelnde Treffsicherheit, Regionalisierung, Niederlassungs\-attraktivität.
  \item \textbf{Regionalisierung} der ambulanten Wirtschaftlichkeitsprüfung durch neue \S\S\,106a--c. Die Auffälligkeitsprüfung verbleibt in \S\,106 SGB\,V; \S\,106b regelt die Einzelfallprüfung und verankert in Abs.\,2 den Beratungsschutz für V1. Die Einzelfallprüfung wird auf die regionale Ebene der Prüfvereinbarungen ausgelagert.
  \item \textbf{Schutz ``Beratung vor Regress''} für die Auffälligkeitsprüfung im Gesetz verankert (\S\,106b Abs.\,2 SGB\,V) --- aber \textbf{für die Einzelfallprüfung explizit ausgeschlossen} (\S\,106b Abs.\,2 S.\,4 SGB\,V).
\end{itemize}

\textbf{Die KBV-Stellungnahme zum GKV-VSG 2015} (S.\,32-33) warnt explizit, dass die Verlagerung der Einzelfallprüfung in regionale Vereinbarungen ``die Gefahr einer Verschärfung der Wirtschaftlichkeitsprüfung birgt''. Das Muster, das die KBV identifizierte: Durch die Regionalisierung bekommen die Krankenkassen mehr Einfluss auf die Ausgestaltung der Prüfvereinbarungen, während die bundesgesetzliche Aufsicht sinkt. Was auf Bundesebene als Entbürokratisierung (Abschaffung der Richtgrößen) präsentiert wurde, wurde auf regionaler Ebene durch eine Ausweitung der Einzelfallprüfung kompensiert.

Das Muster, das die KBV konkret erkannte: (1)~Die Regionalisierung entzog die Ausgestaltung der Prüfverfahren dem Bundesgesetzgeber und übergab sie den regionalen Vertragspartnern (KV + Kassen). (2)~In regionalen Verhandlungen sind die Krankenkassen strukturell im Vorteil, weil sie als Zahlungspflichtige ein direktes finanzielles Interesse an scharfen Prüfverfahren haben. (3)~Was auf Bundesebene als Entbürokratisierung kommuniziert wurde (Abschaffung der Richtgrößen), konnte auf regionaler Ebene durch erweiterte Einzelfallprüfungen kompensiert werden --- ohne parlamentarische Kontrolle.

Die Kassenärzteschaft hat dieses Muster erkannt --- und politisch verloren. Der prognostizierte Effekt trat ein: In Westfalen-Lippe stiegen die Einzelfallprüfungen im ersten Jahr nach Inkrafttreten um 28\,\%, in Nordrhein erreichten sie etwa 20.000 Anträge pro Jahr. Verschiedene Erklärungen für diesen Anstieg stehen nebeneinander: Die Kassen argumentieren, dass die regionalen Prüfvereinbarungen gezielter arbeiten als die alten Richtgrößen. Die KVen vermuten, dass die Kassen die Einzelfallprüfung als Ersatz für die abgeschaffte Richtgrößenprüfung nutzen. Ob die Zunahme durch verbesserte Antragstechnik (Digitalisierung der Prüfprozesse) oder durch gezieltes Ausnutzen des Beratungsschutz-Ausschlusses getrieben wird, ist empirisch nicht geklärt.

\textbf{Zum Beratungsschutz:} Vor dem GKV-VSG 2015 gab es \textit{keinen} gesetzlich verankerten Beratungsschutz für \textit{kein} Verfahren. Das GKV-VSG führte ``Beratung vor Regress'' erstmals als gesetzliche Norm ein --- aber nur für die Auffälligkeitsprüfung (V1). Für die Einzelfallprüfung (V2) wurde der Beratungsschutz im selben Gesetz \textbf{explizit ausgeschlossen} (\S\,106b Abs.\,2 S.\,4). Das heißt: Der Gesetzgeber hat bewusst entschieden, V2 ohne pädagogischen Zwischenschritt zu belassen --- das ist kein Übersehen, sondern Design.

\subsection{Drei methodisch auffällige Wanderungen}

Die Container-Genealogie produziert drei Wanderungen, die methodisch besonders aufschlussreich sind:

\begin{itemize}[nosep]
  \item \textbf{Die \S\,106a-Umbenennungsfalle (2017):} Vor 2017 stand in \S\,106a SGB\,V die Abrechnungsprüfung. Ab 2017 steht in \S\,106a die Wirtschaftlichkeitsprüfung ärztlicher Leistungen. Die Abrechnungsprüfung wanderte gleichzeitig in den neuen \S\,106d. Wer alte Literatur liest, wird systematisch in die Irre geführt. Wer ``\S\,106a'' zitiert, muss seit 2017 dazu sagen, welche Fassung gemeint ist.
  \item \textbf{Die Auslagerungs-Kaskade der Einzelfallprüfung:} BSG-Richterrecht (1959) $\to$ Gesetzestext \S\,106 Abs.\,2 Nr.\,2 (1993) $\to$ bloße Erwähnung in \S\,106 Abs.\,1 S.\,2 plus regionale Prüfvereinbarung (2017). Die Substanz ist drei Ebenen abgewandert. Die Sichtbarkeit hat abgenommen, der Anwendungsbereich nicht.
  \item \textbf{Die parallel-gegenläufige Operation Ambulant/Stationär:} 2015 wurde die ambulante Einzelfallprüfung aus dem Gesetz fast herausgeschrieben, 2020 wurde die stationäre Einzelfallprüfung in einen eigenen \S\,275c verselbständigt. Gegenläufige Logik im selben Gesetzbuch innerhalb von fünf Jahren.
\end{itemize}

\subsection{Strategische Auseinanderreißung: Belege ja, Vorsatz schwer nachweisbar}

Die Container-Architektur weist eine strukturell intransparente Regelungsarchitektur auf, deren Entstehungsmuster sich aus dem Gesetzgebungsprozess ergibt:

\begin{itemize}[nosep]
  \item \textbf{\S\,48 BMV-Ä (Sonstiger Schaden)} ist seit Jahrzehnten außerhalb des SGB\,V kodifiziert --- nicht im Gesetz, sondern im Bundesmantelvertrag. Wer ``nur'' das Gesetz studiert, übersieht das wichtigste verordnungsbezogene Haftungsinstrument. Der Bundesmantelvertrag gilt in der öffentlichen Wahrnehmung nicht als Gesetz.
  \item \textbf{\S\,106b Abs.\,2 S.\,4 SGB\,V} schließt die Einzelfallprüfung explizit aus dem Schutz ``Beratung vor Regress'' aus. Eine bewusste asymmetrische Konstruktion: Wer von der Kasse beantragt wird, hat keinen Anspruch auf Erstgespräch.
  \item Seit 2015 entstehen statt großer Bereinigungen nur noch \textbf{Mikro-Patches} (\S\,106b Abs.\,1a, 1b, 1c --- Grippeimpfstoff, Rückruf, Versorgungslücken). Schutzklauseln werden ins SGB\,V gehoben (Sichtbarmachung), Kontrollwerkzeuge in untergesetzliche Ebenen ausgelagert (Unsichtbarmachung). Niemand korrigiert die Schichtung.
\end{itemize}

\textbf{Beobachtung:} Die Container-Landschaft bewegt sich seit 2015 nicht in Richtung Vereinfachung, sondern in Richtung weiterer Differenzierung. Welche Folgen das für die praktische Anwendbarkeit hat, lässt sich aus der reinen Containeranalyse nicht ableiten --- dafür sind die Befunde aus Teil III (Verhaltensökonomie) und Teil V (institutionelles Schweigen) heranzuziehen.

\subsection{Strukturelle Inkonsistenzen im Normgefüge: Top-10}

Eine eigene Detailanalyse hat zehn handwerklich identifizierbare Schwachstellen in der heutigen Containerlandschaft dokumentiert. Sieben davon wären mit minimalem gesetzgeberischen Aufwand zu beheben. Sie werden nicht behoben.

{\small
\begin{longtable}{p{0.4cm} p{8cm} p{4.5cm}}
\toprule
\textbf{\#} & \textbf{Inkonsistenz} & \textbf{Fixbar?} \\
\midrule
1 & \S\,106a-Umnummerierung 2017 (alte Norm = Abrechnung, neue Norm = Wirtschaftlichkeit) — Pre-2017-Literatur führt in die Irre & Nein, nur über Lesefassungen \\
2 & \S\,48 BMV-Ä unsichtbar — der wichtigste existenzvernichtende Regressgrund (BSG B\,6\,KA\,9/24\,R, 491.000\,EUR) lebt in einem Kollektivvertrag, nicht im SGB\,V & Ja, Verlinkung oder Aufnahme als \S\,106e SGB\,V \\
3 & Beratungs-Asymmetrie \S\,106b Abs.\,2 S.\,4 — ``Beratung vor Regress'' gilt explizit nicht für Einzelfallprüfungen & Ja, einzeiliger Halbsatz-Streich \\
4 & Bagatellgrenzen regional 30--150\,EUR — verfassungsrechtlich problematische Ungleichbehandlung & Ja, Mindestgrenze in \S\,106b SGB\,V \\
5 & Kein Bauplan für die Einzelfallprüfung — Verweis auf 17 regionale Prüfvereinbarungen & Nein, war 2017 politischer Wille \\
6 & Keine Aufsicht über Prüfungsstellen — keine systematische Vollzugsstatistik & Ja, Veröffentlichungspflicht in \S\,106c \\
7 & ``Unwirtschaftlich'' vs. ``unzulässig'' — keine Legaldefinition; Vollregress vs. Differenzkostenregelung als BSG-Richterrecht & Ja, Liste in \S\,106b \\
8 & ``Wirtschaftlichkeit'' in \S\,12 SGB\,V undefiniert — Tautologie statt Tatbestand & Schwer, politisch umstritten \\
9 & Sprechstundenbedarf ohne SGB-V-Container — etwa 10.000 Anträge/Jahr allein in KV Nordrhein, Rechtsgrundlage nur untergesetzlich & Ja, neuer Absatz in \S\,106b \\
10 & BMV-Ä nicht auf gesetze-im-internet.de — strukturelle Sichtbarkeitslücke am amtlichen Rechtsportal & Ja, Verlinkung durch BMJ \\
\bottomrule
\end{longtable}
}

\textbf{Beobachtung:} Sieben der zehn Inkonsistenzen sind trivial fixbar — einzeilige Gesetzesänderung, Verlinkung, Halbsatz-Streich. Sie werden seit Jahren nicht behoben. Diese Tatsache ist selbst ein Datenpunkt: \textit{Es scheint keinen Akteur zu geben, der sowohl die Motivation als auch die Befugnis zur Korrektur besitzt --- was die Persistenz dieser Schwachstellen erklären würde.}


\section{Teil II --- Die rechtsdogmatische Lücke: Verdacht ohne Substanz}

\subsection{Selbstverständnis: Die Einzelfallprüfung ist eine Verdachtsfallprüfung}

Die Einzelfallprüfung versteht sich rechtsdogmatisch eindeutig als Verdachtsfallprüfung, nicht als Screening oder Stichprobenverfahren. Drei Belege:

\begin{itemize}[nosep]
  \item \textbf{\S\,106b Abs.\,1 SGB\,V} und die Rahmenvorgaben (KBV / GKV-Spitzenverband, Stand 01.05.2020) verlangen einen ``schriftlichen und begründeten Antrag'' mit ``ausreichendem Verdacht auf Unwirtschaftlichkeit''. Die Krankenkasse muss die konkrete Verordnung, den Schaden und die Vergleichsleistung substantiiert benennen.
  \item \textbf{Die BSG-Rechtsprechung}, insbesondere B\,6\,KA\,21/19\,R vom 11.09.2019, sieht die Einzelfallprüfung als gerichteten Kontrollakt über eine konkret bezeichnete Verordnung. Die Prüfgremien haben laut BSG keinen Beurteilungsspielraum --- eine Verordnung ist entweder zulässig oder unzulässig, ohne Mittelzustand.
  \item \textbf{\S\,106b Abs.\,2 S.\,4 SGB\,V} schließt ``Beratung vor Regress'' für Einzelfallprüfungen ausdrücklich aus. Begründung des Gesetzgebers: Der Anlass sei bereits substantiiert; eine pädagogische Beratung sei deshalb entbehrlich.
\end{itemize}

Das heißt: Jeder Antrag der Krankenkasse muss --- per Gesetz --- ein substantiierter Verdachtsfall sein. Es darf keine Suchanfrage sein, keine Stichprobe, keine Routineprüfung. Die Beweislast für den Anfangsverdacht liegt bei der Kasse.

\subsubsection*{Screening versus Verdachtsprüfung — die normative Grenze}

Für die folgende Argumentation ist eine begriffliche Präzisierung wichtig. Der Unterschied zwischen einem Screening und einer Verdachtsprüfung ist nicht graduell, sondern normativ:

{\small
\begin{tabularx}{\textwidth}{X X X}
\toprule
\textbf{Kriterium} & \textbf{Screening} & \textbf{Verdachtsprüfung} \\
\midrule
Zweck & systematisches Aufspüren möglicher Probleme & gerichtete Prüfung eines konkret bezeichneten Falles \\
Beweislast & nicht erforderlich, breite Filterwirkung erwünscht & substantiierter Anfangsverdacht zwingend (\S\,106b Abs.\,1 SGB\,V) \\
Zulässige Verwerfungsquote & sehr hoch (das ist der Sinn) & nahe Null bei korrekter Vorprüfung \\
Sanktion bei Treffer & Empfehlung, Beratung, Korrektur & förmliches Verfahren mit Rechtsfolge \\
Beispiel & PKW-Verkehrskontrolle & strafrechtlicher Anfangsverdacht \\
\bottomrule
\end{tabularx}
}


\subsection{Die 99,7-\,\%-Verwerfungsrate: Was sie zeigt --- und was sie verdeckt}

Die KV Hessen wies für die \textbf{Auffälligkeitsprüfung (V1)} Arzneimittel 2015 folgende Pipeline aus: 916 Auffälligkeiten $\to$ 29 eröffnete Verfahren $\to$ 3 durchgesetzte Regresse. Endquote 0,3\,\%, verifiziert direkt im Deutschen Ärzteblatt (Artikel 113630). Vergleichbare Quoten in Thüringen, Nordrhein und Baden-Württemberg.

\textbf{Was die V1-Verwerfungsrate zeigt:} Für ein statistisches Screening ist eine hohe Verwerfungsrate \textit{erwartbar}. Die Auffälligkeitsprüfung ist ein Suchverfahren: Sie markiert statistische Abweichungen vom Fachgruppendurchschnitt, nicht individuelle Verdächte. Dass 99,7\,\% der markierten Fälle nach näherer Prüfung keinen Regress ergeben, ist die normale Funktionsweise eines Filters mit niedriger Spezifität. Aus dieser Perspektive funktioniert V1 wie vorgesehen.

\textbf{Allerdings} ändert das nichts daran, dass alle 916 markierten Ärzte die \textit{Kosten} des Screenings tragen --- Angst, Dokumentationsaufwand, Verhaltensänderung ---, auch wenn für 913 von ihnen kein Regress folgt. Das Instrument muss nicht versagen, um zu schaden.

\textbf{Was die V1-Verwerfungsrate verdeckt:} Für die \textbf{Einzelfallprüfung (V2)} existiert \textit{keine vergleichbare Pipeline-Statistik}. V2-Daten existieren (BMG: ca.\,47.000 Verfahren bundesweit 2022; SpiFa Bayern: 13.000 Regresse/Jahr), werden aber nicht als integrierte Durchsetzungsstatistik mit Verwerfungs- und Vergleichsquoten publiziert. Bekannt sind nur Volumina: Allein in der KV Nordrhein stellen Krankenkassen jährlich ca.\,20.000 Einzelfallprüfungsanträge --- bei gleichzeitig nur $\sim$2 V1-Regressen pro Jahr.

\textbf{Die naheliegende Verwunderung --- und warum sie in die Irre führt.} Auf den ersten Blick scheint die niedrige Gesamtregressquote ($<$1\,\%) zu bedeuten, dass auch bei V2 wenig passiert. Das ist ein Fehlschluss. Beide Verfahren verhalten sich \textit{genau so, wie man es von ihrem jeweiligen Design erwarten würde}:

\begin{itemize}[nosep]
  \item \textbf{V1 (Schwellenwert-Screening nach Auffälligkeiten):} Breites Screening, hoher Dropout $\to$ die meisten markierten Ärzte werden entlastet. Erwartbar.
  \item \textbf{V2 (gezieltes Verdachtsverfahren):} Die Kasse identifiziert einen konkreten Formfehler und stellt einen Antrag. Bei einem tatsächlich vorliegenden Formfehler steht das Ergebnis im Wesentlichen fest (normativer Schadensbegriff, BSG B\,6\,KA\,5/09\,R): Der Arzt \textit{wird} zahlen. Die Frage ist nur noch \textit{wie viel}.
\end{itemize}

\textbf{Warum die Durchsetzung nicht im Gericht sichtbar wird:} Beide Seiten haben rationale Gründe, das Gerichtsverfahren zu vermeiden:

\begin{itemize}[nosep]
  \item \textbf{Der Arzt:} Weiß, dass Gerichte bei festgestelltem Formfehler im Sinne der Kasse entscheiden (normativer Schadensbegriff). Ein Gerichtsverfahren kostet Zeit, Geld und psychische Belastung --- ohne realistische Erfolgsaussicht.
  \item \textbf{Die Kasse:} Hat kein Interesse an vielen Gerichtsverfahren, weil diese die Aufmerksamkeit auf die Verfahrenspraxis lenken und unangenehme Fragen aufwerfen könnten (Häufigkeit der Anträge, Verhältnismäßigkeit, Selektionskriterien). Außerdem kann die Kasse V2-Anträge taktisch einsetzen: bei Kleinigkeiten ansetzen und bei Widerstand zurückziehen --- ein Disziplinierungsinstrument, das außerhalb des Gerichts wirksamer ist als innerhalb.
\end{itemize}

\textbf{Das Ergebnis:} Geschätzt 70--80\,\% der betroffenen Ärzte lassen sich auf einen \textbf{Vergleich} ein (typisch 30--50\,\% des Regressbetrags). Der Vergleich \textit{ist} die Durchsetzung --- nur außerhalb des Gerichts. Die V2-Regressdurchsetzungsquote (gemessen am Geldfluss, nicht an Gerichtsurteilen) dürfte daher \textit{erheblich höher} liegen als die V1-Quote. Aber sie wird nirgends publiziert.

\textbf{Die eigentlichen V2-Probleme} liegen nicht in der Durchsetzungsrate (die ist strukturell hoch), sondern in vier Merkmalen des Verfahrensdesigns:

\begin{enumerate}[nosep]
  \item \textbf{Keine Sanktion bei übersensitiver Antragstellung:} Die Kasse kann beliebig viele V2-Anträge stellen, ohne Konsequenzen für Fehlanträge. Das SG Mainz 2022 (S\,3\,KA\,14/19, 260.000\,EUR) ist der einzige dokumentierte Fall, in dem eine Kasse für Rechtsmissbrauch verurteilt wurde.
  \item \textbf{Keine unabhängige richterliche Prüfung:} Da die Mehrheit der Fälle im Vergleich endet, wird keine richterliche Feststellung zum Einzelfall getroffen. Die Qualität der Kassenanträge wird nie systematisch geprüft.
  \item \textbf{Kein Beratungsschutz:} \S\,106b Abs.\,2 S.\,4 schließt Beratung vor Regress bei V2 gesetzlich aus. Der Arzt wird direkt sanktioniert.
  \item \textbf{Disziplinierungswirkung ohne Aktenzeichen:} Schon die Androhung eines V2-Antrags kann Verordnungsverhalten ändern, ohne dass je ein Verfahren stattfindet. Dieser Chilling Effect ist per Definition nicht messbar --- er hinterlässt keine Spur im System.
\end{enumerate}

Genau diese Intransparenz adressieren die eingereichten IFG-Anfragen (Anhang~C) und das Konzept eines Regress-Transparenzportals (PP-003): Nicht die V1-Verwerfungsrate ist das Problem, sondern die fehlende Messung der V2-Durchsetzungsrealität.

\medskip
\noindent\textbf{Die praktische Konsequenz in einem Satz:} Bei V2b (unzulässige Verordnung) gilt: Wird ein Formfehler festgestellt, muss der Arzt zahlen --- unabhängig davon, ob die Behandlung medizinisch korrekt war (normativer Schadensbegriff, BSG B\,6\,KA\,5/09\,R; Details in Abschnitt~\ref{sec:normativer-schaden}). Das System ``zertifiziert'' nicht medizinische Qualität, sondern formale Regelkonformität --- und sanktioniert deren Fehlen.

\subsection{Begriffsklärung: Zwei Prüfdomänen und vier Messebenen}

\subsubsection*{Drei Schichten des Regressbegriffs}

Der Begriff ``Regress'' wird von verschiedenen Akteuren unterschiedlich definiert. Diese Definitionsdivergenz erklärt, warum die öffentlich kommunizierte Quote von ``unter 1\,\%'' und die von Ärzten berichtete Erfahrungsrate von 72\,\% (Ribbat 2023) \textbf{beide korrekt} sind --- sie messen verschiedene Dinge.

\medskip
{\small
\begin{tabularx}{\textwidth}{l p{5.5cm} l l}
\toprule
\textbf{Schicht} & \textbf{Was gezählt wird} & \textbf{Ungefähre Zahl} & \textbf{Quelle} \\
\midrule
1 (eng) & Nur V1-Regresse, nur rechtskräftig durchgesetzt & ${<}$100/Jahr $\to$ 0,065\,\% & Virchowbund 2018 \\
\addlinespace
2 (amtlich) & V1 + V2, rechtskräftig durchgesetzt & 21 von 21.500 $\to$ 0,097\,\% & KV BaWü 2022 \\
\addlinespace
3 (faktisch) & Alle V2-Verfahren mit Zahlung (inkl. Vergleiche, Bagatelle) & ca.\,47.000/Jahr $\to$ ${\sim}$30\,\% & BMG GVSG 2024 \\
\bottomrule
\end{tabularx}
}

\medskip
\noindent\textbf{Der Faktor zwischen Schicht~1 und Schicht~3 beträgt ca.\,470--500$\times$.} Wer ``Regress'' als rechtskräftigen Richtgrößenregress nach Gerichtsurteil definiert (Schicht~1), kommt auf unter 0,1\,\%. Wer jede V2-Rückforderung ab 30\,EUR mitzählt (Schicht~3), kommt auf 13.000 Fälle pro Jahr allein in Bayern (SpiFa 2024). Die Diskrepanz ist kein Wahrnehmungsfehler der Ärzte, sondern ein Messproblem der Statistik. Dieser Befund wird in der Summierungsrechnung (Teil~III) und der Kommunikationsasymmetrie-Analyse empirisch untermauert.

\subsubsection*{Zwei Prüfdomänen --- funktional getrennt, nicht graduelle Abstufungen}

Das Wirtschaftlichkeitsprüfsystem besteht aus zwei Verfahren, die nicht als unterschiedlich intensive Versionen derselben Logik verstanden werden dürfen. Sie sind funktional inkommensurabel:

\medskip\noindent\textbf{Tabelle: Strukturvergleich Auffälligkeitsprüfung (V1) vs. Einzelfallprüfung (V2)}

{\small
\begin{tabularx}{\textwidth}{p{3.8cm} X X}
\toprule
\textbf{Kategorie} & \textbf{Auffälligkeitsprüfung (V1)} & \textbf{Einzelfallprüfung (V2)} \\
\midrule
\multicolumn{3}{l}{\textit{\textbf{Verfahrensstruktur}}} \\
\addlinespace[2pt]
Bezeichnung & Statistische Auffälligkeitsprüfung & Einzelfallprüfung \\
\addlinespace
Auslösetrigger & Kosten-/Mengen\-abweichung vom Fachgruppen\-durch\-schnitt & Konkreter Antrag einer Krankenkasse \\
\addlinespace
Rechts\-grundlage & \S\,106 SGB\,V & \S\,106b Abs.\,1 SGB\,V \\
\addlinespace
Auslösende Institution & Prüfungsstelle (KV\,+\,Kassen) & Krankenkasse \\
\addlinespace
Charakter & Such- und Filterverfahren & Gezieltes Verdachtsverfahren \\
\addlinespace
Verdachts\-qualität & Kein individueller Verdacht (statistisches Signal) & Individueller, sub\-stantiierter Verdacht erforderlich \\
\midrule
\multicolumn{3}{l}{\textit{\textbf{Schutz- und Präventionsstruktur}}} \\
\addlinespace[2pt]
Beratung vor Regress & Vorgeschrieben (\S\,106b Abs.\,2) & Gesetzlich ausgeschlossen (\S\,106b Abs.\,2 S.\,4) \\
\addlinespace
Präventive Gegen\-maßnahmen & Praxisbesonder\-heiten, Doku\-mentation, Anpassung & KV-Beratung, KK-Vorab\-anfrage (BSG\,B\,6\,KA\,27/12\,R), antrags\-feste Doku \\
\addlinespace
Maßnahmen bei Eintritt & Stellungnahme, Beratung, ggf.\,Anpassung & Rechtfertigung, Widerspruch, Klage \\
\midrule
\multicolumn{3}{l}{\textit{\textbf{Risikoprofil}}} \\
\addlinespace[2pt]
Häufigkeit (je Patient) & Sehr gering & Sehr gering \\
\addlinespace
Häufigkeit (Haus\-arzt, ca.\,1.200 Pat.) & Gelegentlich & Regelmäßig möglich \\
\addlinespace
Betroffenheit bei Eintritt & Teilweise & 100\,\% \\
\addlinespace
Wahrscheinlichkeit neg.\,Verfahrens\-ausgang & Niedrig & Hoch \\
\addlinespace
Wahrscheinlichkeit pos.\,Ausgang & Hoch & Mittel bis niedrig \\
\addlinespace
Folge bei neg.\,Ausgang & Regressfestsetzung & Regressfestsetzung \\
\addlinespace
Schaden\-höhe & Niedrig bis moderat & Moderat bis hoch \\
\addlinespace
Existenz\-bedrohend & Sehr gering & Niedrig, aber real \\
\bottomrule
\end{tabularx}
}

Die \textbf{Auffälligkeitsprüfung} (V1, statistisches Screeningverfahren) ist ein von der Prüfungsstelle initiiertes Verfahren zur Identifikation wirtschaftlicher Abweichungen im Vergleich zur jeweiligen Fachgruppe. Sie erfolgt ohne individuellen Verdacht. Charakteristisch ist eine bewusst niedrige Trefferquote: Die hohe Sensitivität des Filters ist systemimmanent und gewollt, da das Ergebnis lediglich ein vorläufiges Abweichungssignal darstellt. Die Auffälligkeitsprüfung trifft weder eine Aussage zur medizinischen Angemessenheit noch zur individuellen Verantwortlichkeit.

Die \textbf{Einzelfallprüfung} (V2, antragsgebundene Verdachtsprüfung) ist eine fallbezogene Prüfung, die gemäß \S\,106b Abs.\,1 SGB\,V ausschließlich auf schriftlichen und begründeten Antrag einer Krankenkasse eingeleitet wird. Sie ist von Gesetzes wegen bereits als Ankläger, Richter und Begünstigter in einer Struktur angelegt (\S\,106b Abs.\,2 S.\,4: kein Beratungsschutz; Kasse = Antragstellerin + Empfängerin des Regressbetrags). Wird ein Formfehler bestätigt, folgt unmittelbar eine sanktionsrechtliche Entscheidung --- der Bescheid ist praktisch das Urteil (siehe normativer Schadensbegriff, Abschnitt~\ref{sec:normativer-schaden}).

\textbf{Die empirische Konsequenz:} Die KV Hessen (916 Auffälligkeiten $\to$ 3 Regresse, 0,3\,\%) demonstriert normales Screening-Verhalten der Auffälligkeitsprüfung (V1) --- viel filtern, wenig finden ist für einen breiten Filter korrekt und methodisch unauffällig. Für die Einzelfallprüfung (V2) existieren keine vergleichbaren Veröffentlichungen. Die 99,7\,\%-Dropout-Rate bezieht sich ausschließlich auf V1.

\textbf{Warum zwei Verfahren, die Ähnliches zu messen scheinen?} V1 und V2 teilen dasselbe Ergebnis (Regressfestsetzung), prüfen aber grundlegend Verschiedenes: V1 fragt ``Liegt der Arzt statistisch außerhalb der Norm?'' (Fachgruppen\-durch\-schnitt), V2 fragt ``Ist \textit{diese konkrete Verordnung} formal angreifbar?'' (Einzelfall). Historisch erklärt sich die Koexistenz aus der Container-Genealogie (Teil~I): Die Einzelfallprüfung existiert seit 1959 (BSG-Richterrecht), die statistischen Prüfverfahren kamen später als ergänzendes Breitband-Instrument hinzu. Die Abschaffung der Richtgrößen 2017 hat V1 geschwächt, V2 aber unberührt gelassen --- beide Verfahren bestehen nebeneinander, weil kein Reformgesetz sie integriert hat.

\subsubsection*{Vier Messebenen --- nicht ein ``Regressquote''-Begriff}

Der Begriff ``Regressquote'' wird in der öffentlichen Diskussion und in Teilen dieser Studie für mehrere grundlegend verschiedene Messgeößen verwendet, die nicht miteinander vergleichbar sind:

{\small
\begin{tabularx}{\textwidth}{p{2.8cm} X X p{3.5cm}}
\toprule
\textbf{Begriff} & \textbf{Definition} & \textbf{Was gemessen wird} & \textbf{V1/V2} \\
\midrule
Prüfquote & Anteil der Ärzt:innen oder Fälle, die in ein formelles Prüfverfahren gezogen werden & Verfahrens\-beteiligung & V1: publiziert (regional); V2: k.D. \\
\addlinespace
Regress\-verfahren\-quote & Anteil der geprüften Fälle, bei denen eine Regress\-forderung gestellt wird & Antrags\-intensität & V1: $\sim$3\,\% (KV\,Hessen); V2: k.D. \\
\addlinespace
Regress\-durch\-setzungs\-quote & Anteil der Regress\-verfahren, bei denen ein Regress rechts\-kräftig festgesetzt wird & Verfahrens\-erfolg & V1+V2 zusammen: $<$1\,\%; getrennt: k.D. \\
\addlinespace
Regress\-volumen & Tatsächlich gezahlte Regress\-beträge in Euro & Geldfluss & V1+V2 zusammen; getrennt: k.D. \\
\bottomrule
\end{tabularx}
}

Die in dieser Studie zitierten Zahlen messen verschiedene Ebenen:

\begin{itemize}[nosep]
  \item ``Unter 1\,\% der Ärzte'' (verschiedene KV-Statistiken): \textbf{Regressdurchsetzungsquote auf Ärztebene, über beide Verfahren (V1+V2) hinweg, ohne Verfahrenstrennung}
  \item ``99,7\,\% Dropout'' (KV Hessen 2015, Auffälligkeitsprüfung): \textbf{Pipeline-Verwerfungsrate auf Verfahrensebene}
\end{itemize}

Diese zwei Zahlen sind kein Widerspruch, sondern komplementäre Messungen verschiedener Ebenen derselben Realität. Bundesweit fehlen alle vier Ebenen als einheitlich veröffentlichte Zahlen --- genau das ist die Datenbasis, die PP-003 schließen will.

\textbf{Zur Regressdurchsetzungsquote bei V2:} Für die Einzelfallprüfung existieren keine publizierten Durchsetzungsquoten. Aus der Verfahrensstruktur lässt sich jedoch ableiten, dass sie wesentlich höher liegen muss als bei V1: Fehlt ``Beratung vor Regress'' (\S\,106b Abs.\,2 S.\,4), gibt es keinen institutionalisierten Abbruchpunkt zwischen Antrag und Bescheid. Der normative Schadensbegriff (Abschnitt~\ref{sec:normativer-schaden}) macht den Regress bei festgestelltem Formfehler praktisch unvermeidlich. Die meisten Verfahren enden nicht durch Freispruch, sondern durch Vergleich --- was in der Statistik als ``kein Regress'' erscheinen kann, obwohl der Arzt 30--50\,\% des Betrags gezahlt hat.

\subsubsection*{Eskalationspipeline: Wo welche Regressquote gilt}

Beide Prüfverfahren lassen sich als Eskalationstreppe lesen. Die zitierte \textbf{,,unter 1\,\% der Ärzte``-Zahl} bezieht sich auf Stufe~9 von Verfahren~1~--- und gilt nicht für das Einzelfallverfahren, für das keine vergleichbaren Bundesstatistiken existieren. Tabellen von unten lesen: Schritt~0 = Verfahrensauslöser, Schritt~9 = höchste Eskalationsstufe (rechtskräftiger Regress).

\smallskip\noindent \textbf{Struktureller Unterschied:} Verfahren~1 ist \textit{probabilistisch}~--- viele Ärzte erreichen frühe Stufen, sehr wenige gelangen bis Schritt~n. Verfahren~2 ist nahezu \textit{deterministisch}~--- wer Schritt~0 erreicht, durchläuft die Pipeline fast vollständig. \textbf{Betroffenheit beginnt bei Schritt~1, nicht erst beim Geldfluss.}

\medskip
\noindent\textbf{Verfahren 1: Auffälligkeitsprüfung} (\S\,106 SGB\,V)
\hfill{\small$\uparrow$ Eskalation}

{\footnotesize
\begin{tabularx}{\textwidth}{p{3.5cm} p{1.9cm} p{2.3cm} p{0.6cm} p{2.1cm} p{2.1cm}}
\toprule
\textbf{Schritt / Methode} & \textbf{Umgangs-sprache} & \textbf{Verfahrens-Begriff} & \textbf{Reg.} & \textbf{P\,ges.} & \textbf{P\,$|$\,S0} \\
\midrule
Schritt 9: Regress rechts\-kräftig --- selten & ,,Der Regress ist durch`` & Regress rechts\-kräftig & 1 & k.D.* KV\,BaWü: 0,097\,\% (V1+V2) & \textbf{$\sim$0,3\,\%} KV\,Hessen: 3\,/\,916 \\
\addlinespace
Schritt 8: Klage --- Sozial\-gericht & ,,Ich gehe vor Gericht`` & Regress streitig & 1 & k.D. & k.D. \\
\addlinespace
Schritt 7: Widerspruchs\-entscheidung & ,,Entscheidung gefallen`` & Regress aufgehoben od.\,bestätigt & 1 & k.D. & k.D. \\
\addlinespace
Schritt 6: Widerspruch & ,,Ich wehre mich`` & Regress ange\-fochten & 1 & k.D. & k.D. \\
\addlinespace
Schritt 5: Regress\-fest\-setzung --- Ausnahme\-fall & ,,Ich habe einen Regress`` & Regress fest\-gesetzt (n.\,rechts\-kr.) & 1 & $\sim$0,1\,\% & $\sim$0,3\,\% \\
\addlinespace
Schritt 4: Entscheidung --- \textbf{Screening\-abschluss (Stopper für $>$99\,\%)} & ,,Die Prüfung ist durch`` & Kein Regress (Regelfall) & 0 & k.D. & \textbf{hoch} ($>$99\,\%) \\
\addlinespace
Schritt 3: Beratung vor Regress --- Schutz\-stufe & ,,Ich werde beraten`` & Regress ausge\-schlossen & 0 & k.D. & mittel \\
\addlinespace
Schritt 2: Anhörung + Praxis\-besonder\-heiten & ,,Ich erkläre mich`` & Kein Regress & 0 & k.D. & $\sim$100\,\% \\
\addlinespace
Schritt 1: Prüf\-verfahren ein\-geleitet --- Screening & ,,Die prüfen mich`` & Prüf\-verfahren (kein Verdacht) & 0 & $\sim$10--20\,\% [$^\dagger$] & $\sim$100\,\% \\
\addlinespace
Schritt 0: Trigger: Kosten-/Mengen\-abweichung erkannt & --- & Kein Verfahren & 0 & $\sim$10--20\,\% [$^\dagger$] & 100\,\% \\
\bottomrule
\end{tabularx}

\smallskip
\noindent[*] \textbf{Quellen und Messmethodik:}\\
(1)~\textbf{KV\,BaWü 2022} (0,097\,\% = ca.\,21\,/\,21.500\,Ärzte): KV\,Baden-Württemberg, Informationsseite Regressgefahr (\url{https://www.kvbawue.de/praxis/verordnungen/arzneimittel/regressgefahr}); Messmethode: durchgesetzte Arzneimittelregresse\,/\,Vertragsärzte gesamt (V1+V2, ohne Verfahrenstrennung).\\
(2)~\textbf{KV\,Hessen 2015} (3\,/\,916 = 0,3\,\%\,$|$\,S0): \textit{Deutsches Ärzteblatt}, Artikel~113630; Messmethode: vollständige Pipeline-Zählung (916 Auffälligkeitsverfahren $\to$ 29 geprüft $\to$ 3 Regresse; P\,$|$\,S0 = Regressdurchsetzungsrate V1).\\
\textit{Beide Werte: verschiedene Populationen, Verfahren, Jahre --- kein Direktvergleich.} Die ,,unter 1\,\%``-Formel gilt für V1+V2; keine Verfahrenstrennung veröffentlicht.\\
{[$^\dagger$] Bundesweite Prüfquoten nicht veröffentlicht; Schätzung auf Basis regionaler Berichte. Veröffentlichter Datenpunkt: KV\,Hessen 2015 (Dt.~Ärzteblatt, Art.~113630): 916 ausgelöste Auffälligkeitsverfahren Arzneimittel --- Gesamtzahl Vertragsärzte KV\,Hessen für dieses Jahr nicht separat publiziert (Denominator für absolute Prüfquote damit nicht rekonstruierbar).}\\
Reg. = Regress objektiv eingetreten (0\,=\,nein; 1\,=\,ja, auch wenn noch nicht rechts\-kräftig).\\
P\,ges. = Wahrscheinlichkeit diese Stufe zu erreichen (Bezugsgröße: alle Vertragsärzte).\\
P\,$|$\,S0 = Wahrscheinlichkeit bezogen auf Population die Schritt~0 erreicht hat. \quad
k.D. = keine publizierten Daten verfügbar.
}

\smallskip
\noindent\textit{\textbf{Kernaussage V1:} Screening produziert viele Signale und kaum Schaden. Nahezu alle Ärzte, die Schritt~1 erreichen, verlassen die Pipeline ohne Regress. \textbf{Schritt~4 (Screeningabschluss) ist für über 99\,\% der Betroffenen der Endpunkt} --- die Eskalation endet, bevor ein Regress entsteht.}

\medskip
\noindent\textbf{Verfahren 2: Einzelfallprüfung} (\S\,106b Abs.\,1 SGB\,V)
\hfill{\small$\uparrow$ Eskalation}

{\footnotesize
\begin{tabularx}{\textwidth}{p{3.5cm} p{1.9cm} p{2.3cm} p{0.6cm} p{2.1cm} p{2.1cm}}
\toprule
\textbf{Schritt / Methode} & \textbf{Umgangs-sprache} & \textbf{Verfahrens-Begriff} & \textbf{Reg.} & \textbf{P\,ges.} & \textbf{P\,$|$\,S0} \\
\midrule
Schritt 9: Regress rechts\-kräftig --- häufiger als V1 & ,,Der Regress ist durch`` & Regress rechts\-kräftig & 1 & k.D. & k.D. --- \textit{deutlich höher als V1} \\
\addlinespace
Schritt 8: Klage --- Sozial\-gericht & ,,Ich gehe vor Gericht`` & Regress streitig & 1 & k.D. & k.D. \\
\addlinespace
Schritt 7: Widerspruchs\-entscheidung & ,,Entscheidung gefallen`` & Regress aufgehoben od.\,bestätigt & 1 & k.D. & k.D. \\
\addlinespace
Schritt 6: Widerspruch & ,,Ich wehre mich`` & Regress ange\-fochten & 1 & k.D. & k.D. \\
\addlinespace
Schritt 5: Regress\-fest\-setzung --- Regelfall & ,,Ich habe einen Regress`` & Regress fest\-gesetzt (n.\,rechts\-kr.) & 1 & k.D. & \textbf{hoch} \\
\addlinespace
Schritt 4: Bewertung --- Verdacht bestätigt/verneint & ,,Die entscheiden`` & Regress möglich\,/\,entfällt & 0/1 & k.D. & mittel--hoch \\
\addlinespace
Schritt 3: Formale\,/\,doku\-mentarische Prüfung & ,,Die prüfen meine Unterlagen`` & Regress\-verfahren aktiv & 0 & k.D. & $\sim$100\,\% \\
\addlinespace
\multicolumn{6}{l}{\textit{--- Beratung vor Regress: gesetzlich nicht vorgesehen (\S\,106b Abs.\,2 S.\,4) --- Schutzstufe fehlt ---}} \\
\addlinespace
Schritt 2: Unterlagen\-anforderung --- Recht\-fertigungs\-phase & ,,Ich muss mich recht\-fertigen`` & Regress\-verfahren aktiv & 0 & k.D. & $\sim$100\,\% \\
\addlinespace
Schritt 1: Prüf\-verfahren ein\-geleitet --- Verdachts\-verfahren & ,,Die prüfen mich gezielt`` & Regress\-verfahren ein\-geleitet & 0 & k.D. [$^{\ddagger}$] & $\sim$100\,\% \\
\addlinespace
Schritt 0: Trigger: Konkreter Kassen\-antrag --- sub\-stantiierter Verdacht & --- & Kein Verfahren & 0 & k.D. [$^{\ddagger}$] & 100\,\% \\
\bottomrule
\end{tabularx}

\smallskip
\noindent[$^{\ddagger}$] KV\,Nordrhein: ca.\,20.000 Einzelfallprüfungs-Anträge\,/\,Jahr bei ca.\,13.500 Praxen (Quelle: interne KV-Berichte 2016--2018; zitiert in der Berichter\-stattung zur GKV-VSG-Umsetzung; keine extern verifizierbare Bundesstatistik veröffentlicht; Verfahrens\-abschluss\-quoten nicht publiziert).\\
Reg. = Regress objektiv eingetreten (0\,=\,nein; 0/1\,=\,verfahrensabhängig; 1\,=\,ja).\\
P\,ges. = alle Vertragsärzte als Bezugsgröße. \quad
P\,$|$\,S0 = bedingte Wahrscheinlichkeit gegeben Schritt~0. \quad
k.D. = keine publizierten Daten.
}

\smallskip
\noindent\textit{\textbf{Kernaussage V2:} Wenn Auslösung --- dann nahezu vollständige Eskalation. Wer Schritt~0 erreicht, bekommt in der überwiegenden Mehrzahl der Fälle einen Regress. Die fehlende Schutzstufe ,,Beratung vor Regress`` (markierte Zeile) ist kein redaktionelles Versehen, sondern explizites gesetzliches Design (\S\,106b\,Abs.\,2\,S.\,4 SGB\,V). Diese Kernaussage gilt für V2b (unzulässige/formfehlerhafte Verordnungen). Bei V2a (unwirtschaftliche Verordnungen) bestehen Verteidigungsmöglichkeiten: Der Arzt kann die medizinische Begründung vorbringen, und die Kostendifferenzmethode begrenzt den Regress auf die Differenz zum wirtschaftlichen Niveau.}

% Pipeline-V2-Grafik entfernt (v0.11, LG: passt nicht mehr zum Textstil)

\subsection{Termination without Determination: Die unsichtbare Ökonomie der Vergleiche}

Die 99,7-\,\%-Verwerfungsrate der Auffälligkeitsprüfung (V1, KV Hessen) ist aussagekräftiger als bisher dargestellt --- aber in eine unerwartete Richtung. Die bisherige Analyse hat gefragt: \textit{Warum scheitern so viele Verfahren?} Die entscheidende Folgefrage lautet: \textit{Wie enden sie --- und welche epistemische Qualität haben diese Enden?}

Regressverfahren können auf mindestens fünf Wegen enden, die sich in ihrer epistemischen Qualität fundamental unterscheiden:

{\small
\begin{tabularx}{\textwidth}{X X X}
\toprule
\textbf{Beendigungsweg} & \textbf{Rechtlicher Weg} & \textbf{Feststellung der Berechtigung?} \\
\midrule
Prüfungsstelle lehnt Antrag ab & Ablehnungsbescheid & Ja --- durch Prüfungsstelle \\
Informeller Vergleich\,/\,Zahlung & \S\S\,53--62 SGB\,X, Kassenrücknahme & Nein --- kein Bescheid, kein Urteil \\
Kasse zieht Antrag zurück & Einstellung ohne Sachentscheidung & Nein \\
Verjährung & Fristablauf & Nein \\
Gerichtliches Urteil & Urteil, bestandskräftig & Ja --- höchste epistemische Qualität \\
\bottomrule
\end{tabularx}
}

Nur zwei der fünf Wege produzieren eine inhaltliche Feststellung. Da fast alle Verfahren vor Erreichen des Gerichts enden, bleibt die Systemqualität dauerhaft ungemessen.

Das erzeugt ein \textbf{strukturelles Anreizproblem}: Kassen, die einen unbegründeten Antrag gestellt haben, können ihn vor Gericht zurückziehen, ohne dass die Unbegründetheit je festgestellt wird. Ärzte, die aus Angst vor Aufwand oder Reputationsschaden einer Vergleichszahlung zustimmen, ``anerkennen'' damit keine Schuld --- aber das System wertet die Zahlung faktisch als Bestätigung. Die Verwerfungsquote misst damit nicht die Qualität der Anträge, sondern die \textbf{Gegenwehrstärke der Betroffenen}.

Das Muster ist strukturell identisch mit dem deutschen Abmahnwesen: Viele Anträge, die meisten zahlen aus Angst vor Kosten und Aufwand, kaum einer klagt. Die Yield kommt nicht aus gerichtlichen Erfolgen, sondern aus der Asymmetrie von Angstkosten und Abwehrkosten. Das Gericht wäre der einzige Ort, an dem unabhängig festgestellt werden könnte, ob ein Verfahren berechtigt war. Weil fast kein Verfahren je ein Gericht erreicht, bleibt der Qualitätssensor dauerhaft ohne Messwerte.

\subsection{Die Asymmetrie zwischen Auffälligkeits- und Einzelfallprüfung}

Die beiden Prüfverfahren widersprechen sich strukturell:

{\small
\begin{tabularx}{\textwidth}{X X X X}
\toprule
\textbf{Prüfungsart} & \textbf{Initiator} & \textbf{Beratung vor Regress} & \textbf{Schutz} \\
\midrule
Auffälligkeitsprüfung & Prüfungsstelle (neutral) & Ja (\S\,106b Abs.\,2) & Gehoben \\
Einzelfallprüfung & Krankenkasse (Antragstellerin) & \textbf{Nein} (\S\,106b Abs.\,2 S.\,4) & Abgesenkt \\
\bottomrule
\end{tabularx}
}

Während die Auffälligkeitsprüfung über die Jahre Schutzschichten dazubekommen hat (Beratung vor Regress 2012, Praxisbesonderheiten 2011, Richtgrößen abgeschafft 2015), blieb die Einzelfallprüfung explizit ohne. Aus Sicht des Anwendungsbereichs ist sie damit das verbleibende Verfahren mit den geringsten Schutzklauseln.

Ob das eine bewusste Schutzlücke oder eine zufällige Reformresidue ist, lässt sich aus den Gesetzgebungsmaterialien nicht eindeutig klären. Die KBV-Stellungnahme zum GKV-VSG 2015 hat die Verlagerung als ``Gefahr einer Verschärfung'' kritisiert (siehe oben), wurde aber im Verfahren nicht gefolgt. Methodisch ist daher festzuhalten: Das Muster ist vorhanden, der Vorsatz nicht beweisbar.

\subsubsection{Ausschlussfristen: Wie weit zurück geprüft werden darf}
\label{sec:ausschlussfristen}

Ein weiterer Asymmetrie-Aspekt betrifft die zeitliche Reichweite der Prüfverfahren. Die Fristen wurden mehrfach reformiert. Wichtig: Es handelt sich um \textbf{Ausschlussfristen} (wirken von Amts wegen, das Recht erlischt automatisch), nicht um Verjährungsfristen (müssten vom Arzt einredeweise geltend gemacht werden). Für Ärzte ist das die günstigere Variante.

{\small
\begin{tabularx}{\textwidth}{l X X}
\toprule
 & \textbf{V1 Auffälligkeitsprüfung} & \textbf{V2 Einzelfallprüfung} \\
\midrule
\multicolumn{3}{l}{\textit{\textbf{Aktuelles Recht (ab TSVG 11.05.2019 / GVWG 20.07.2021)}}} \\
\addlinespace[2pt]
Festsetzungsfrist & \textbf{2 Jahre} ab Jahresende bzw. Honorarbescheid & 18 Mon. Antragsfrist + 12 Mon. Festsetzung = \textbf{2,5 Jahre} \\
\addlinespace
Rechtsgrundlage & \S\,106 Abs.\,3 S.\,3 SGB\,V & \S\,106 Abs.\,3 S.\,4--5 SGB\,V \\
\midrule
\multicolumn{3}{l}{\textit{\textbf{Altes Recht (vor 11.05.2019 --- gilt für Altfälle!)}}} \\
\addlinespace[2pt]
Ausschlussfrist & \textbf{4 Jahre} (BSG-Richterrecht) & \textbf{4 Jahre} (BSG-Richterrecht) \\
\addlinespace
Rechtsgrundlage & BSG 6\,RKa\,40/94 (20.09.1995) & BSG B\,6\,KA\,13/18\,R (11.09.2019) \\
\bottomrule
\end{tabularx}
}

\medskip

\textbf{Hintergrund:} Die vierjährige Ausschlussfrist wurde 1995 vom BSG richterrechtlich entwickelt (6\,RKa\,40/94), in Anlehnung an \S\,45 SGB\,I und andere sozialrechtliche Vier-Jahres-Fristen. Der Gesetzgeber hat sie mit dem TSVG (2019) und GVWG (2021) auf 2--2,5 Jahre verkürzt. Begründung im Gesetzgebungsverfahren: ``Zeiträume von mehr als zwei Jahren zwischen dem geprüften Verordnungszeitraum und dem Abschluss der Prüfungen [sind] für die Betroffenen unzumutbar.''

\textbf{Praktische Einschränkung:} Die Ausschlussfrist begrenzt nur, \textit{bis wann ein Erstbescheid ergehen muss}. Wurde das Verfahren innerhalb der Frist eingeleitet, gibt es \textbf{keine zeitliche Obergrenze} für den weiteren Instanzenzug. Das BSG hat Verfahren entschieden, die 8--11 Jahre nach dem Verordnungszeitraum bei ihm ankamen (z.\,B. B\,6\,KA\,6/21\,R, 2022: Verordnungen aus 2011). Für den betroffenen Arzt kann die Belastung durch ein laufendes Verfahren also weit über die Ausschlussfrist hinausreichen.

\subsection{Der Zertifikatsfehler: Was das System vorgibt zu messen --- und was es tatsächlich misst}

\subsubsection{Zwei inkommensurable Messdomänen}

Beide Prüfverfahren produzieren Ergebnisse, die \textbf{überwiegend nicht messen, was \S\,12 SGB\,V als Schutzziel benennt}. Sie messen Proxys --- aus verschiedenen Domänen, mit verschiedenen Methoden und fundamental verschiedener epistemischer Qualität. \textit{Einschränkung:} Ein Teil der V2-Prüfungen (geschätzt ${\sim}$35\,\%, Kategorie~A) codiert durchaus medizinische Evidenz als Formregeln --- AM-RL-Anlagen, G-BA-Nutzenbewertungen und Biosimilar-Pflichten sind inhaltlich evidenzbasiert, auch wenn sie formal geprüft werden (vgl. Abschnitt~\ref{sec:evidenz-formale-verkleidung}).

\textit{Begriffserklärung ``Zertifikat'':} In der folgenden Tabelle bezeichnet \textbf{Zertifikat} eine belastbare, inhaltliche Aussage über die medizinische Angemessenheit einer Versorgungsleistung im Sinne des \S\,12 Abs.\,1 SGB\,V. Ein Zertifikat liegt nur vor, wenn das Prüfergebnis tatsächlich misst, ob die Versorgung ``ausreichend, zweckmäßig und wirtschaftlich'' war --- nicht lediglich, ob eine formale oder statistische Auffälligkeit vorliegt.

{\small
\begin{tabularx}{\textwidth}{p{3cm} X X}
\toprule
\textbf{Instanz} & \textbf{Selbstanspruch} & \textbf{Tatsächlich gemessen\,/\,Zertifikat?} \\
\midrule
Auffälligkeitsprüfung (\S\,106 SGB\,V) & Wirtschaftliche Versorgung (Screening) & Kosten-/Mengenabweichung vom Fachgruppen\-durch\-schnitt $\to$ statistisch-ökonomisches Abweichungssignal \textbf{(kein medizin. Zertifikat)} \\
\addlinespace
Einzelfallprüfung & Medizinische Unwirt\-schaft\-lich\-keit im Einzelfall & Formale\,/\,doku\-mentarische Angreifbarkeit $\to$ formales Sanktionsresultat \textbf{(kein medizin. Zertifikat)} \\
\addlinespace
Praxisbesonderheiten & Korrektur legitimer Abweichung & Verfahrensabhängige Anerkennung; anmeldepflichtig, regional intransparent $\to$ Validität exogen (kein Zertifikat) \\
\addlinespace
Ärztliche Dokumentation & Medizinische Begründung & Selbstbeschreibung unter systemischem Druck $\to$ Dokumentation \textbf{(kein Zertifikat)} \\
\addlinespace
Medizinischer Dienst (MD) & Leistungs-\,/\,Abrechnungsprüfung & Regel- und zweckgebundene Anspruchsbewertung \textbf{(kein Qualitätszertifikat)} \\
\addlinespace
Amtsärztliche Gutachten & Status-\,/\,Verwaltungsfragen & Zustandsbewertung (Dienstfähigkeit etc.) \textbf{(kein Versorgungszertifikat)} \\
\addlinespace
Forschung\,/\,Register\,/\,QS (IQTIG etc.) & Qualitätsverbesserung & Aggregierte Qualitätssignale; nicht individualisierbar, nicht sanktionslogisch \textbf{(kein Zertifikat)} \\
\addlinespace
Schlichtungsstellen (Ärztekammern) & Behandlungsstandard & Patientenbezogene medizinische Angemessenheit $\to$ \textbf{medizinisches Zertifikat} (selten, ex post, systemisch entkoppelt) \\
\addlinespace
Zivilrechtliche Arzthaftung & Behandlungsfehler & Medizinischer Standard\,+\,Kausalität $\to$ \textbf{starkes medizinisches Zertifikat} (ex post, reaktiv) \\
\midrule
\textbf{Kein Verfahren} & \textbf{\S\,12 SGB\,V: bedarfsgerechte, zweckmäßige Versorgung} & \textbf{wird im Regress- und Prüfsystem systematisch nicht gemessen} \\
\bottomrule
\end{tabularx}
}

Medizinische Angemessenheit --- das Schutzziel des \S\,12 Abs.\,1 SGB\,V (``Die Leistungen müssen ausreichend, zweckmäßig und wirtschaftlich sein; sie dürfen das Maß des Notwendigen nicht überschreiten.'') --- wird im deutschen Vertragsarztrecht nur außerhalb des Regress- und Prüfsystems zertifiziert: in Schlichtungsverfahren der Ärztekammern und in zivilrechtlichen Arzthaftungsprozessen. Diese Verfahren sind reaktiv, selten, ex post und systemisch vollständig vom Wirtschaftlichkeitsprüfsystem entkoppelt.

Abbildungen~\ref{fig:pipeline-auff} und~\ref{fig:pipeline-einzel} zeigen beide Prüfprozesse im Detail. An keiner Stufe beider Verfahren entsteht eine medizinisch verwertbare Aussage. Auch der vollständige Rechtsweg (Widerspruch, Klage) korrigiert ausschließlich formale Entscheidungen --- er ersetzt keine medizinische Bewertung der Versorgung.

\begin{figure}[p]
\centering
\resizebox{\textwidth}{!}{%
\begin{tikzpicture}[
  proc/.style={draw, rounded corners=3pt, fill=gray!5, align=left,
               font=\small, inner sep=5pt, text width=11cm},
  trigger/.style={draw, rounded corners=3pt, fill=umbruch-red!10, align=left,
                  font=\small\bfseries, inner sep=5pt, text width=11cm},
  endok/.style={draw, rounded corners=3pt, fill=green!8, align=left,
                font=\small, inner sep=4pt, text width=5cm},
  warnbox/.style={draw, rounded corners=3pt, fill=red!8, align=left,
                  font=\small, inner sep=5pt, text width=11cm},
  dec/.style={draw, rounded corners=2pt, fill=yellow!8, align=center,
              font=\small\bfseries, inner sep=4pt, text width=6cm},
  arr/.style={->, thick},
  lbl/.style={font=\small}
]

\node[trigger] (root) at (0,0)
  {\textbf{Ärztliche Versorgung (Regelbetrieb)}\\
   Verantwortlich: Vertragsarzt \quad Externe Messung: keine\\
   Zertifikat: \textcolor{red}{$\times$} \textit{(Selbstzertifizierung durch Dokumentation)}};

\node[proc] (auff) at (0,-2.8)
  {\textbf{Auffälligkeitsprüfung}\\
   Initiator: Prüfungsstelle \quad Verfahren: Statistisches Screening\\
   Soll-Messung: Wirtschaftlichkeit (\S\,12 SGB\,V)\\
   Ist-Messung: Kosten-/Mengenabweichung vom Fachgruppen\-durch\-schnitt\\
   Outcome: statistisches Abweichungssignal \quad Zertifikat: \textcolor{red}{$\times$}};
\draw[arr] (root.south) -- (auff.north)
  node[lbl, right, midway] {stat. Abweichung erkannt};

\node[proc] (anhoer) at (0,-5.6)
  {Anhörung\,/\,Stellungnahme des Arztes\\
   Zweck: Erklärung der Abweichung \quad Mittel: Praxisbesonderheiten};
\draw[arr] (auff.south) -- (anhoer.north);

\node[dec] (dec1) at (0,-7.4) {Abweichung erklärbar?};
\draw[arr] (anhoer.south) -- (dec1.north);

\node[endok] (end1) at (8.5,-7.4)
  {Verfahren endet\\
   Entscheidung: Prüfungsstelle\\
   Wirkung Arzt: Entlastung\\
   Zertifikat: \textcolor{red}{$\times$}};
\draw[arr] (dec1.east) -- (end1.west) node[lbl, above, midway] {JA};

\node[proc] (berat) at (0,-9.6)
  {\textbf{Beratung vor Regress} (zwingend, \S\,106b Abs.\,2 SGB\,V)\\
   Verantwortlich: Prüfungsstelle \quad Zweck: pädagogische Korrektur\\
   Neue Messung: keine \quad Zertifikat: \textcolor{red}{$\times$}};
\draw[arr] (dec1.south) -- (berat.north) node[lbl, right, midway] {NEIN};

\node[dec] (dec2) at (0,-11.8) {Anpassung erfolgt?};
\draw[arr] (berat.south) -- (dec2.north);

\node[endok] (end2) at (8.5,-11.8)
  {Verfahren endet\\
   Wirkung Arzt: Anpassung\\
   Wirkung Patient: indirekt\\
   Zertifikat: \textcolor{red}{$\times$}};
\draw[arr] (dec2.east) -- (end2.west) node[lbl, above, midway] {JA};

\node[warnbox] (ausgang) at (0,-14.2)
  {\textbf{Nachteiliger Verfahrensausgang}\\
   \textbf{Soll:} Übergang in Einzelfallprüfung\\
   \textbf{Ist:} häufig direkte Regressfestsetzung\\
   $\Rightarrow$ Übergang ohne medizinische Re-Zertifizierung \quad
   Zertifikat: \textcolor{red}{$\times$}};
\draw[arr] (dec2.south) -- (ausgang.north) node[lbl, right, midway] {NEIN};

\end{tikzpicture}%
}
\caption{Pipeline A: Auffälligkeitsprüfung (statistisches Screening). An jeder Stufe fehlt ein medizinisches Zertifikat. Die Anhörung mit Praxisbesonderheiten ist der einzige Punkt, an dem der Arzt die Messvalidität aktiv beeinflussen kann.}
\label{fig:pipeline-auff}
\end{figure}

\begin{figure}[p]
\centering
\resizebox{\textwidth}{!}{%
\begin{tikzpicture}[
  proc/.style={draw, rounded corners=3pt, fill=gray!5, align=left,
               font=\small, inner sep=5pt, text width=11cm},
  trigger/.style={draw, rounded corners=3pt, fill=umbruch-red!10, align=left,
                  font=\small\bfseries, inner sep=5pt, text width=11cm},
  procrw/.style={draw, rounded corners=3pt, fill=blue!5, align=left,
                 font=\small, inner sep=5pt, text width=11cm},
  endok/.style={draw, rounded corners=3pt, fill=green!8, align=left,
                font=\small, inner sep=4pt, text width=5cm},
  endwarn/.style={draw, rounded corners=3pt, fill=red!8, align=left,
                  font=\small, inner sep=4pt, text width=5cm},
  dec/.style={draw, rounded corners=2pt, fill=yellow!8, align=center,
              font=\small\bfseries, inner sep=4pt, text width=6cm},
  arr/.style={->, thick},
  lbl/.style={font=\small}
]

\node[trigger] (trigger) at (0,0)
  {\textbf{Konkreter Verdacht (Krankenkasse)}\\
   Initiator: Krankenkasse \quad Rechtsgrundlage: \S\,106b Abs.\,1 SGB\,V\\
   Voraussetzung: schriftlicher, begründeter Antrag mit ausreichendem Verdacht};

\node[proc] (einzel) at (0,-3.0)
  {\textbf{Einzelfallprüfung}\\
   Verantwortlich: Prüfgremium \quad Verfahren: Verdachtsprüfung\\
   \textbf{V2a} (Unwirtschaftlichkeit): Wirtschaftlichkeit der Verordnung\\
   \quad Medizinische Verteidigung möglich, Kostendifferenzmethode\\
   \textbf{V2b} (Formfehler/Unzulässigkeit): Formale Angreifbarkeit\\
   \quad Medizinisch irrelevant, normativer Schadensbegriff\\
   Beratung vor Regress: \textcolor{red}{$\times$} (\S\,106b Abs.\,2 S.\,4 SGB\,V)\\
   Zertifikat: \textcolor{red}{$\times$}};
\draw[arr] (trigger.south) -- (einzel.north);

\node[proc] (pruef) at (0,-6.3)
  {Prüfung der Unterlagen\\
   Kriterien: formale Richtlinientreue, Dokumentationslage, Plausibilität\\
   \textit{(keine originäre medizinische Bewertung)}};
\draw[arr] (einzel.south) -- (pruef.north);

\node[dec] (dec1) at (0,-8.2) {Verdacht unbegründet?};
\draw[arr] (pruef.south) -- (dec1.north);

\node[endok] (end1) at (8.5,-8.2)
  {Verfahren endet\\
   Entscheidung: Prüfgremium\\
   Wirkung Arzt: Entlastung\\
   Zertifikat: \textcolor{red}{$\times$}};
\draw[arr] (dec1.east) -- (end1.west) node[lbl, above, midway] {JA};

\node[proc] (regress) at (0,-10.5)
  {\textbf{Regressfestsetzung}\\
   Wirkung Arzt: finanzielle Sanktion\\
   Wirkung Patient: indirekt (Defensivmedizin)\\
   Zertifikat: \textcolor{red}{$\times$} \textit{(formales Sanktionsresultat)}};
\draw[arr] (dec1.south) -- (regress.north) node[lbl, right, midway] {NEIN};

\node[procrw] (widersp) at (0,-13.2)
  {\textbf{Widerspruchsverfahren} (Widerspruchsausschuss)\\
   Prüfmaßstab: Rechtmäßigkeit, Verfahrensfehler, Ermessensfehler\\
   \textit{Keine originäre medizinische Neubewertung}\\
   Zertifikat: \textcolor{red}{$\times$}};
\draw[arr] (regress.south) -- (widersp.north)
  node[lbl, right, midway] {Widerspruch möglich (fristgebunden)};

\node[dec] (dec2) at (0,-15.4) {Widerspruch erfolgreich?};
\draw[arr] (widersp.south) -- (dec2.north);

\node[endok] (end2) at (8.5,-15.4)
  {Regress aufgehoben\\
   Ergebnis: formale Entlastung\\
   Zertifikat: \textcolor{red}{$\times$}};
\draw[arr] (dec2.east) -- (end2.west) node[lbl, above, midway] {JA};

\node[procrw] (klage) at (0,-17.7)
  {\textbf{Klage} (Sozialgericht)\\
   Prüfmaßstab: Rechtmäßigkeit, Beweislast\\
   \textit{Medizinische Bewertung nur mittelbar}\\
   Zertifikat: \textcolor{red}{$\times$}};
\draw[arr] (dec2.south) -- (klage.north) node[lbl, right, midway] {NEIN};

\node[dec] (dec3) at (0,-19.8) {Klage erfolgreich?};
\draw[arr] (klage.south) -- (dec3.north);

\node[endok] (end3) at (8.5,-19.8)
  {Regress aufgehoben\\
   Zertifikat: \textcolor{red}{$\times$}};
\draw[arr] (dec3.east) -- (end3.west) node[lbl, above, midway] {JA};

\node[endwarn] (rkraeft) at (0,-21.8)
  {\textbf{Regress rechtskräftig}\\
   Zertifikat: \textcolor{red}{$\times$}\\
   \textit{Auch hier: kein medizinisches Zertifikat}};
\draw[arr] (dec3.south) -- (rkraeft.north) node[lbl, right, midway] {NEIN};

\end{tikzpicture}%
}
\caption{Pipeline B: Einzelfallprüfung (Verdachtsprüfung). Auch der vollständige Rechtsweg (Widerspruch~$\to$ Klage) erzeugt kein medizinisches Zertifikat --- er korrigiert ausschließlich formale Entscheidungen.}
\label{fig:pipeline-einzel}
\end{figure}

\subsubsection{Warum so wenige Verfahren vor Gericht landen}

An keiner Stelle des Prüfprozesses --- weder bei der Auffälligkeitsprüfung noch bei der Einzelfallprüfung --- entsteht eine medizinisch verwertbare Aussage über die Versorgungsqualität. Auch der Beschwerdeausschuss prüft nur die Rechtmäßigkeit des Bescheids, nicht die medizinische Angemessenheit der Behandlung. Erst das Sozialgericht wäre die erste unabhängige Instanz. Doch fast alle Verfahren enden vorher --- durch Vergleich, Rücknahme oder Verjährung.

\textbf{Warum?} Die Antwort unterscheidet sich fundamental nach Verfahren:

\textbf{Bei V1 (Auffälligkeitsprüfung):} Die meisten Verfahren enden bereits bei Schritt~4 (Screening\-abschluss) --- der Arzt erklärt die Abweichung, die Prüfungsstelle akzeptiert. Eine Klage wäre widersinnig, weil das Verfahren erfolgreich für den Arzt endet.

\textbf{Bei V2 (Einzelfallprüfung):} Das Ergebnis steht nach Verfahrenseinleitung im Wesentlichen fest (normativer Schadensbegriff, siehe Abschnitt~\ref{sec:normativer-schaden}). Drei Mechanismen entwerten den Rechtsweg für den Arzt: (1)~Das Gericht würdigt bei Formfehlern die medizinische Korrektheit der Behandlung nicht --- selbst exzellente Medizin schützt nicht vor dem Regress. (2)~Kein Beratungsschutz (\S\,106b Abs.\,2 S.\,4): Es gibt keinen pädagogischen Zwischenschritt, der das Verfahren noch abwenden könnte. (3)~Der normative Schadensbegriff definiert den Schaden als ab Verordnung eingetreten --- die nachträgliche Heilung des Formfehlers ist prozessual ausgeschlossen.

Die Kasse ihrerseits bevorzugt einen Vergleich (30--50\,\% des Betrags, sofort, ohne öffentliche Aufmerksamkeit) gegenüber einem Prozess (100\,\% des Betrags, aber mit Risiko, Aufwand und Sichtbarkeit). Beide Seiten haben rationale Gründe, das Gericht zu meiden. Die wenigen Fälle, die Gericht erreichen, sind nicht repräsentativ --- sie sind die Spitze einer Selektion, bei der beide Seiten schwache Fälle vorher aussortiert haben.

\textbf{Machtungleichgewicht:} Für die Kasse ist ein unbegründeter Antrag praktisch folgenlos: \S\,106b SGB\,V enthält keine Sanktionen für unbegründete Anträge, keine Kostentragungspflicht bei Falsch-Positiv und keine Mengenbegrenzung. Die Prüfungskosten werden paritätisch umgelegt (\S\,106 Abs.\,4 SGB\,V). Für den Arzt ist jeder Antrag existenziell belastend. Dieses asymmetrische Risikoprofil begründet ein strukturelles Missbrauchspotenzial: Die Kasse kann es sich leisten, breit Anträge zu stellen und auf Vergleiche zu setzen, ohne ein Gerichtsverfahren zu riskieren. Die einzige dokumentierte Ausnahme: Das Sozialgericht Mainz wies 2022 einen Regress von 260.000\,EUR als \textit{rechtsmissbräuchlich} zurück (SG Mainz, S\,3\,KA\,14/19, 07.12.2022; fehlende Gegenzeichnung bei Chemotherapie-Verordnungen, medizinisch korrekt, kein persönlicher Vorteil des Arztes) --- ein Einzelfall, keine gesetzliche Schranke.

\subsubsection{Der doppelte Zertifikatsfehler}

\noindent\textit{Wichtige Differenzierung vorweg:} Der doppelte Zertifikatsfehler gilt in voller Schärfe nur für \textbf{V2b} (Formfehler/unzulässige Verordnungen). V1 hat ein Schutzstufenmodell (Anhörung, medizinische Berücksichtigung, Beratung vor Regress, Karenzzeit, Deckel 25.000\,EUR). V2a lässt medizinische Verteidigung und Kostendifferenzmethode zu. Die folgende Analyse fokussiert auf die \textit{strukturelle Anlage} beider Verfahren --- die Schutzlücken konzentrieren sich auf V2b.

\medskip
\textbf{Erster Zertifikatsfehler:} Die Auffälligkeitsprüfung misst Kostenabweichung, nicht medizinische Qualität. Der hohe Dropout ist für ein Screening-Verfahren technisch korrekt --- ein breiter Filter mit niedriger Spezifität ist funktional erwartet. Das ändert nichts daran, dass alle 916 markierten Ärzte die Kosten des Screenings tragen (Angst, Aufwand, Verhaltensänderung), auch wenn für 913 von ihnen kein Regress folgt. Das Instrument muss nicht versagen, um zu schaden.

\textbf{Zweiter Zertifikatsfehler:} Die Einzelfallprüfung beansprucht, medizinische Unwirtschaftlichkeit im Einzelfall zu prüfen (\S\,106b Abs.\,1 SGB\,V). Sie misst faktisch überwiegend formale und dokumentarische Kriterien. Die implizite Annahme, Formfehler könnten als Proxy für medizinische Fehlversorgung dienen, ist epistemisch nicht haltbar:

\begin{itemize}[nosep]
  \item Formfehler können bei exzellenter Medizin auftreten.
  \item Medizinische Fehlversorgung kann bei formgerechter Dokumentation vollständig unsichtbar bleiben.
  \item Formkonformität ist kein Proxy für medizinische Qualität.
\end{itemize}

Das Einzelfallprüfverfahren misst nicht medizinische Fehler, sondern formale Angreifbarkeit --- und behandelt diese fälschlich als medizinische Evidenz. Dieser Befund gilt in voller Schärfe für V2b (unzulässige Verordnungen). Bei V2a (unwirtschaftliche Verordnungen) kann der Arzt die Therapieentscheidung medizinisch begründen; die Kostendifferenzmethode begrenzt den Regress. Dieses Proxy-Versagen ist kein Einzelfall, sondern eine bekannte systemische Pathologie: \textit{Goodhart's Law} (``When a measure becomes a target, it ceases to be a good measure'') und \textit{Campbell's Law} (``Je stärker ein quantitativer Indikator für soziale Entscheidungsfindung genutzt wird, desto mehr ist er Korruptionsdruck ausgesetzt'') beschreiben exakt diesen Mechanismus. Die Prüfquote wird zum Ziel, die Formkonformität zum Messobjekt --- und beide verlieren dadurch ihre diagnostische Aussagekraft über die tatsächliche Versorgungsqualität.

\medskip
\noindent\textit{Kanonischer Befund:}

\begin{quote}
Das Regresssystem weist einen doppelten Zertifikatsfehler auf. Der gesetzliche Auftrag der Wirtschaftlichkeitsprüfung (\S\S\,106\,ff. SGB\,V) ist die Ausgabenkontrolle, nicht die Qualitätszertifizierung --- und diesen Auftrag erfüllt das System. Der Fehlschluss liegt in der öffentlichen und politischen Interpretation: Die niedrige Regressquote wird als Beleg für ein ``funktionierendes System'' kommuniziert, wobei ``funktionierend'' implizit medizinische Versorgungsqualität meint. Tatsächlich misst die Auffälligkeitsprüfung Kostenabweichungen und die Einzelfallprüfung formale Angreifbarkeit. Medizinische Angemessenheit im Sinne des \S\,12 SGB\,V wird systematisch nur außerhalb des Regresssystems zertifiziert. Der doppelte Zertifikatsfehler beschreibt diese Lücke zwischen Systemmessung und öffentlicher Bedeutungszuschreibung.

\smallskip\noindent\textit{Einschränkung:} V1 berücksichtigt medizinische Besonderheiten im Schutzstufenmodell (Anhörung, Praxisbesonderheiten, Beratung). V2a lässt medizinische Verteidigung zu. Der Befund gilt uneingeschränkt für V2b.
\end{quote}

\noindent\textit{Hinweis zur Gegenposition:} Ob der Zertifikatsfehler ein konstruiertes Kriterium ist --- ob also die Studie dem System ein Qualitätsversprechen unterstellt, das es nie gegeben hat ---, wird in Abschnitt~\ref{sec:gegenposition} (Einwand~2) eingehend diskutiert. Der Einwand ist dogmatisch zutreffend und wird dort eingeräumt.

\textbf{Die naheliegende Frage: Warum gibt es keine Versorgungsprüfung?} Wenn das Wirtschaftlichkeitsprüfsystem medizinische Versorgungsqualität nicht misst --- warum existiert kein paralleles System, das genau das tut? Die Antwort ist mehrstufig: (1)~Die stationäre Qualitätssicherung (IQTIG, \S\,136 SGB\,V) arbeitet mit aggregierten Indikatoren und ist nicht auf den ambulanten Bereich übertragbar. (2)~Der Medizinische Dienst (MD) prüft Leistungs- und Abrechnungsangemessenheit, aber nicht die therapeutische Entscheidung selbst. (3)~Eine individuelle Versorgungsprüfung wäre methodisch anspruchsvoll: Sie müsste den Behandlungsstandard zum Zeitpunkt der Verordnung ex post bewerten --- das leisten nur die Schlichtungsstellen der Ärztekammern (selten, reaktiv, freiwillig) und die Zivilgerichte (Arzthaftung). Keines dieser Instrumente ist systemisch an das Regressverfahren angeschlossen. Die Regelungslücke ist keine technische Unmöglichkeit, sondern ein Ergebnis historisch getrennter Verwaltungskulturen: Die Wirtschaftlichkeitsprüfung gehört zur Selbstverwaltung (KV + Kassen), die Qualitätssicherung zum G-BA, die Haftung zum Zivilrecht.

\subsubsection{Formfehler als strukturelle Strategie}

Aus dem Zertifikatsfehler der Einzelfallprüfung ergibt sich ein struktureller Strategievorteil für die antragstellende Seite. Formfehler sind ein besonders effizientes Angriffsziel:

\begin{itemize}[nosep]
  \item \textbf{Billig zu prüfen:} Kein medizinisches Gutachten nötig, keine Fachexpertise erforderlich.
  \item \textbf{Eindeutig:} Unterschrift fehlt oder Stempel statt Unterschrift --- kein Ermessensspielraum.
  \item \textbf{Angstwirksam:} Jeder Arzt kann potentiell Formfehler produzieren; jeder Arzt ist exponiert.
  \item \textbf{Nicht medizinisch streitbar:} Kein Gegengutachten kann die formale Angreifbarkeit aufheben, unabhängig davon wie gut die Medizin war.
\end{itemize}

Das erzeugt eine inverse Selektionslogik: Das System selektiert \textbf{gegen Dokumentationsschwäche, nicht gegen Versorgungsdefizite}. Ein Arzt, der systematisch schlechte Medizin macht aber sauber dokumentiert, ist im System vollständig unsichtbar. Ein Arzt, der exzellente Medizin macht aber Formfehler produziert, ist maximal exponiert.

\subsubsection{Der normative Schadensbegriff: Warum Formfehler ab Verfahrenseinleitung nicht mehr heilbar sind}
\label{sec:normativer-schaden}

Der normative Schadensbegriff des BSG (B\,6\,KA\,5/09\,R, 05.05.2010) verbindet den Zertifikatsfehler mit der Determinismus-Eigenschaft der Einzelfallprüfung. Seine Logik:

\begin{enumerate}[nosep]
  \item \textbf{Vor Verfahrenseinleitung:} Formfehler können korrigiert werden (fehlende Unterschrift nachtragen, Angaben ergänzen). Kein Schaden.
  \item \textbf{Ab Verfahrenseinleitung:} Der Schaden der Kasse gilt als normativ eingetreten --- im Moment der Bezahlung der formal fehlerhaften Verordnung. Ob das Medikament medizinisch indiziert war, ist \textbf{irrelevant}. ``Ersparte Aufwendungen'' werden nicht gegengerechnet.
  \item \textbf{Gerichtliche Prüfung:} Das Sozialgericht prüft ausschließlich: Lag der Formfehler zum Zeitpunkt der Verordnung vor? Ja oder Nein.
  \item \textbf{Ergebnis:} Ab Schritt~2 steht das Verfahrensergebnis im Wesentlichen fest.
\end{enumerate}

Das BSG hat dies 2025 (B\,6\,KA\,9/24\,R) bei 491.000\,EUR Stempel-Regress bestätigt: Der Arzt argumentierte, alle Verordnungen seien medizinisch indiziert. Das Gericht erklärte dieses Argument für irrelevant. 2024 schloss das BSG die Differenzkostenregelung (\S\,106b Abs.\,2a) für ``unzulässige'' Verordnungen (= Formfehler) ausdrücklich aus (B\,6\,KA\,5/23\,R; B\,6\,KA\,10/23\,R).

\textbf{Strukturelle Konsequenz:} Der normative Schadensbegriff erzeugt für die Einzelfallprüfung eine Anreizstruktur, die einem Vergleichssystem ähnelt: Geschätzt 70--80\,\% der betroffenen Ärzte schließen Vergleiche (typisch 30--50\,\% des Regressbetrags; Schätzung aus Praxisberichten von Regressberatern und ärztlichen Fachverbänden, keine publizierte Studie), weil die Erfolgsaussichten vor Gericht bei formalen Fehlern nahe Null liegen. Ärzte, die dennoch den Rechtsweg beschreiten, tun dies häufig in der Erwartung, das Gericht würde die medizinische Korrektheit der Behandlung würdigen. Diese Erwartung wird durch den normativen Schadensbegriff systematisch enttäuscht.

\textbf{Ausnahmen:} In seltenen Fällen haben Gerichte trotz Formfehler zugunsten des Arztes entschieden: Das Sozialgericht Mainz wies 2022 einen Regress von 260.000\,EUR als \textit{rechtsmissbräuchlich} zurück (S\,3\,KA\,14/19, 07.12.2022; Treu und Glauben, \S\,242 BGB analog --- kein realer Schaden, kein persönlicher Vorteil des Arztes). Das LSG Baden-Württemberg entschied 2025 (L\,5\,KA\,1505/23), dass eine Einzelfallprüfung nicht für Verordnungen durchgeführt werden darf, die bereits in einer Richtgrößenprüfung unbeanstandet geblieben sind (Doppelprüfungsverbot). Diese Ausnahmen sind jedoch nicht verallgemeinerbar und belegen gerade durch ihre Seltenheit, wie wenig Schutz der Rechtsweg typischerweise bietet.

\subsubsection{Praxisbesonderheiten: Der Validitäts-Reparaturmechanismus und seine Grenzen}

Die Auffälligkeitsprüfung enthält einen eingebauten Mechanismus zur Verbesserung ihrer medizinischen Validität: \textbf{Praxisbesonderheiten}. Ein Arzt, dessen Kostenprofil statistisch auffällig ist, kann geltend machen, dass seine Patientenstruktur die Abweichung erklärt (höherer Anteil schwer kranker Patienten, spezielles Fachgebiet etc.). In der Praxis scheitert dieser Korrekturmechanismus an drei strukturellen Hürden:

\begin{enumerate}[nosep]
  \item \textbf{Anmeldepflicht:} Praxisbesonderheiten müssen aktiv geltend gemacht werden --- sie werden nicht von Amts wegen berücksichtigt.
  \item \textbf{Regionale Intransparenz:} Anerkennungskriterien variieren erheblich zwischen KVen, ohne bundesweit einheitliche Maßstäbe.
  \item \textbf{Restriktive Anerkennung:} In der Praxis werden Praxisbesonderheiten eng ausgelegt; der Nachweis ist aufwendig und unsicher.
\end{enumerate}

Die Validität des statistischen Signals hängt damit nicht vom Messinstrument ab, sondern davon, ob der Arzt aktiv und erfolgreich interveniert. Wer das nicht tut --- aus Unkenntnis, Aufwandsscheu oder aus Angst, durch Widerspruch zusätzliche Aufmerksamkeit auf sich zu ziehen --- erhält ein epistemisch minderwertigeres Verfahrensergebnis. Das ist ein weiterer Mechanismus der Regressangst, der bisher wenig diskutiert wird: nicht die Angst vor dem Regress selbst, sondern die Angst davor, durch Gegenwehr zusätzliche Sichtbarkeit zu erzeugen.

\subsubsection{Das System kennt seinen Fehler: Vier informelle Ausgleichs\-mechanismen}

Das Wirtschaftlichkeitsprüfsystem hat --- trotz seiner beschriebenen Messdefizite --- informelle Korrekturkanäle entwickelt. Ihre Existenz ist analytisch bedeutsam: Sie zeigen, dass das System selbst weiß, dass seine formalen Messinstrumente unvollständig sind, und sich Ausgleichswege geschaffen hat, ohne diese institutionell zu integrieren.

{\small
\begin{tabularx}{\textwidth}{l l X l}
\toprule
\textbf{Zeitpunkt} & \textbf{Typ} & \textbf{Instrument} & \textbf{Rechtsgrundlage} \\
\midrule
Vor Verordnung & fachlich-epistemisch & KV-pharmakologische Beratung (z.\,B. KVBW, Kooperation Univ.-Kliniken) & KV-intern \\
\addlinespace
Vor Verordnung & administrativ-faktisch & Ärztliche Anfrage zur Verordnungs\-fähigkeit an Krankenkasse & BSG B\,6\,KA\,27/12\,R \\
\addlinespace
Nach Auffälligkeit & statistisch-validierend & Praxisbesonderheiten & \S\,106 Abs.\,5a SGB\,V \\
\addlinespace
Nach Regresseinleitung & prozessual-schützend & Beratung vor Regress & \S\,106b Abs.\,2 SGB\,V \\
\bottomrule
\end{tabularx}
}

Alle vier Instrumente teilen eine strukturelle Schwäche: Sie sind freiwillig, unverbindlich, regional heterogen und in keiner Bundesstatistik erfasst. Niemand weiß, wie oft sie genutzt werden~--- genau das ist die Datenlücke, die ein Transparenzportal (PP-003) schließen würde.

\textbf{Analytischer Kernsatz:} Die Existenz dieser informellen Ausgleichsmechanismen belegt, dass der Zertifikatsfehler nicht nur theoretisch bekannt ist, sondern praktisch kompensiert wird~--- aber bislang ausschließlich individuell, nicht systemisch. Wer Zugang zu diesen Wegen hat und sie kennt, ist strukturell besser geschützt. Wer sie nicht kennt, trägt das volle systemische Risiko.

\textbf{Einschränkung bei Instrument 2 (KK-Anfrage):} Eine positive Stellungnahme der Krankenkasse bindet die Kasse als \textit{Antragstellerin} in einer Einzelfallprüfung~--- sie kann nicht selbst Regress beantragen und zuvor die Verordnungsfähigkeit bestätigt haben (venire contra factum proprium). Sie bindet jedoch nicht die Prüfungsstelle bei einer von Amts wegen initiierten Auffälligkeitsprüfung. Der Schutz ist stark gegen KK-initiierte Einzelfallregresse~--- kein Schutz gegen Prüfungsstellen-initiierte Auffälligkeitsprüfungen.

\textbf{Analytische Pointe:} Die KK-Stellungnahme ist strukturell das Instrument, das einem medizinischen Zertifikat am nächsten kommt --- eine inhaltliche Bewertung der Verordnungsfähigkeit \textit{vor} dem Schadenseintritt. Dass sie existiert, aber freiwillig, informell und nicht flächendeckend ist, zeigt die Lücke im System: Der Ansatz für eine präventive Qualitätsbewertung ist vorhanden, aber nicht institutionalisiert. Ein systemisch verankertes Pendant --- etwa eine verbindliche Vorabklärung bei kostenintensiven oder off-label-Verordnungen --- könnte den Zertifikatsfehler partiell schließen.

\subsection{Evidenz in formaler Verkleidung: Die Doppelfunktion des Prüfsystems}
\label{sec:evidenz-formale-verkleidung}

Die bisherige Analyse behandelte ``Formalität'' als einheitlich medizinisch irrelevant. Das war zu pauschal. Eine genauere Betrachtung der AOK-BaWü-Prüfthemen 2025 und der AM-RL-Struktur zeigt \textbf{zwei Ebenen der Formalität}:

\medskip
{\small
\begin{tabularx}{\textwidth}{p{2.5cm} X X}
\toprule
& \textbf{Ebene 1: Echte Formalität} & \textbf{Ebene 2: Codierte Evidenz} \\
\midrule
Was geprüft wird & Stempel statt Unterschrift, fehlende ICD-Kodierung, Aut-idem-Kreuz & AM-RL Anlage~III, Biosimilar-Pflicht, Off-Label-Verbot, Mengengrenzen \\
\addlinespace
Medizinischer Gehalt & Keiner --- die Medizin stimmt, nur die Dokumentation nicht & Hoch --- G-BA hat Gleichwertigkeit/Evidenz bewertet und in Formregel übersetzt \\
\addlinespace
Zertifikatsfehler & Gilt in voller Schärfe & Gilt \textit{nicht} --- die ``Form'' \textit{ist} die Evidenz \\
\bottomrule
\end{tabularx}
}

\medskip
\noindent\textbf{Die Doppelfunktion:} Schätzt man die Verteilung der V2-Prüfanlässe anhand der AOK-BaWü-Prüfthemenliste 2025 und der SpiFa-Bayern-Daten, ergibt sich ein Dreiteilungsbild:

\medskip
{\small
\begin{tabularx}{\textwidth}{l l p{3cm} X}
\toprule
\textbf{Kategorie} & \textbf{Anteil} & \textbf{Charakter} & \textbf{Beispiele} \\
\midrule
A: Rationale Evidenzsteuerung & ${\sim}$35\,\% & Medizinischer Inhalt als Formregel codiert & Biosimilar-Pflicht, Off-Label-Kontrolle (GLP-1-RA), Opioide Mengengrenzen, RSV-Prophylaxe Altersgrenzen \\
\addlinespace
B: Reine Formalkontrolle (V2b) & ${\sim}$40\,\% & Kein medizinischer Inhalt & Stempel statt Unterschrift, fehlende ICD-Kodierung, Dokumentationslücken \\
\addlinespace
C: Bürokratie-Overhead & ${\sim}$25\,\% & Kleinstbeträge ${<}$300\,EUR & Bagatellrückforderungen, Kodierungsfehler \\
\bottomrule
\end{tabularx}
}

\medskip
\noindent\textbf{Zwei parallele Evidenzkanäle:} Das Gesundheitssystem verfügt über zwei Steuerungsmechanismen für medizinische Evidenz --- einen edukativen und einen punitiven:

\medskip
{\small
\begin{tabularx}{\textwidth}{p{3cm} X X}
\toprule
& \textbf{Fortbildungspflicht (CME, \S\,95d SGB\,V)} & \textbf{Einzelfallprüfung (V2)} \\
\midrule
Timing & \textit{Vor} dem Fehler (präventiv) & \textit{Nach} dem Fehler (punitiv) \\
Signal & ``Biosimilar ist gleichwertig'' (Lernsignal) & ``Sie haben kein Biosimilar verordnet, zahlen Sie'' (Strafsignal) \\
Wahl des Arztes & Arzt wählt Fortbildungsthemen selbst & Kasse wählt Prüfgegenstand \\
Sanktion & 10--25\,\% Honorarkürzung (indirekt) & Voller Regressbetrag (direkt) \\
Pflicht & 250 CME-Punkte / 5 Jahre & Keine --- nur reaktiv \\
\bottomrule
\end{tabularx}
}

\medskip
\noindent\textbf{Konsequenz für die zentrale These:} Der Zertifikatsfehler gilt in voller Schärfe nur für Ebene~1 (echte Formalität, Kategorie~B ${\sim}$40\,\% + Kategorie~C ${\sim}$25\,\% Bagatelle = zusammen ${\sim}$65\,\% der Prüfungen). Bei Ebene~2 (${\sim}$35\,\%, Kategorie~A) steckt die medizinische Qualität \textit{implizit} in der Formregel --- der G-BA hat die Evidenz bewertet und in eine Vorschrift übersetzt. Das System \textit{enthält} medizinische Evidenz, \textit{kommuniziert} sie aber nicht als Evidenz, sondern als Strafe. Auch bei Kategorie~A (rationale Steuerung) fehlt der Beratungsschutz (\S\,106b Abs.\,2 S.\,4): Der Arzt erhält \textit{kein} Lernsignal vor dem Regress, sondern \textit{nur} ein Strafsignal danach. Das ist wie Verkehrserziehung ausschließlich durch Bußgelder ohne Fahrschule. Die Forderung dieser Studie bleibt daher bestehen --- aber sie verschiebt sich: Nicht Abschaffung der Evidenzsteuerung, sondern Ergänzung des punitiven Kanals um einen edukativen (vgl. Einwand~5 in Abschnitt~\ref{sec:gegenposition}).

\subsection{Die operative Synthese: Drei Scanner, keine Qualitätsprüfung}
\label{sec:drei-scanner}

Fasst man die bisherigen Befunde zusammen --- Zertifikatsfehler, normativer Schadensbegriff, Formfehler-Strategie, informelle Ausgleichsmechanismen ---, ergibt sich ein Gesamtbild, das über die einzelnen Befunde hinausgeht. Das deutsche Wirtschaftlichkeitsprüfsystem betreibt drei Screening-Mechanismen. Keiner davon prüft medizinische Versorgungsqualität:

{\small
\begin{tabularx}{\textwidth}{l X X l}
\toprule
\textbf{Scanner} & \textbf{Was er sucht} & \textbf{Was er nicht prüft} & \textbf{Rechtsgrundlage} \\
\midrule
Auffälligkeitsprüfung (V1) & Statistische Abweichung vom Fachgruppendurchschnitt (Kosten\-volumen) & Ob die Versorgung medizinisch angemessen war & \S\,106 SGB\,V \\
\addlinespace
Einzelfallprüfung (V2a) & Prüft Wirtschaftlichkeit einzelner Verordnungen (medizinische Begründung möglich) & Ob die Therapieentscheidung medizinisch richtig war & \S\,106b SGB\,V \\
\addlinespace
Einzelfallprüfung (V2b) & Scannt formale Korrektheit (Unterschrift, Kodierung, Indikation). Prüft nicht Therapiequalität. & Ob die Therapieentscheidung medizinisch richtig war & \S\,106b SGB\,V \\
\addlinespace
Retaxprüfung & Formale Korrektheit der Apothekenabrechnung & Ob der Patient das richtige Medikament bekommen hat & \S\,129 SGB\,V \\
\bottomrule
\end{tabularx}
}

\medskip

Keines dieser Verfahren prüft, ob der Patient ausreichend versorgt wurde (\S\,12 SGB\,V), ob eine Unterlassung stattfand, ob die Therapieentscheidung leitliniengemäß war, oder ob der Arzt aus Angst \textit{weniger} verordnet hat als medizinisch geboten. Der Versorgungsauftrag des \S\,70 SGB\,V --- ``eine bedarfsgerechte und gleichmäßige [...] Versorgung der Versicherten zu gewährleisten'' --- hat kein eigenes Messinstrument im Prüfsystem.

\subsubsection{Das System als Formfehler-Erkennung}

Die operative Realität der Einzelfallprüfung legt eine Interpretation nahe, die über die bisherige Diskussion hinausgeht: Die Kassen suchen nicht primär nach medizinischer Unwirtschaftlichkeit, sondern nach \textbf{Formfehlern}. Der Formfehler \textit{ist} der Verdacht. Das System ist dabei extrem effizient in dem, was es operativ tut: Formfehler finden und regressieren. Nur eben nicht in dem, wofür es normativ da sein soll.

Die Kausalkette:

\begin{enumerate}[nosep]
  \item Kasse erhält Verordnungsdaten (\S\,300 SGB\,V, Apothekenrechenzentren).
  \item Prüfsoftware gleicht Verordnungen ab (PZN$\leftrightarrow$ICD, Formcheck, Dosierung).
  \item Flag: Formfehler identifiziert (fehlende Unterschrift, falscher Code etc.).
  \item Sachbearbeiter sichtet Flag, bestätigt: ``begründeter Verdacht''.
  \item Formaler Antrag an Prüfungsstelle (\S\,106b Abs.\,1).
  \item Normativer Schadensbegriff greift (BSG B\,6\,KA\,5/09\,R).
  \item Formfehler ab Verfahrenseinleitung nicht mehr heilbar.
  \item Gericht prüft: Formfehler ja/nein. Medizinische Indikation irrelevant.
  \item 70--80\,\% der Ärzte schließen Vergleich (30--50\,\% des Betrags).
  \item Regress fließt an die Krankenkasse zurück.
\end{enumerate}

Unter Lesart C erklärt sich die hohe Verwerfungsquote anders: Die 97\,\% Einstellungen sind Fälle \textit{ohne} verwertbaren Formfehler. Die 3\,\% mit Regress sind Fälle \textit{mit} Formfehler. Das System optimiert auf Einhaltung formaler Kriterien --- wie ein Rechtschreibprüfprogramm. Es tut genau das, wofür es operativ gebaut ist.

\textbf{Methodische Einschränkung:} Dies ist eine Arbeitshypothese, keine bewiesene Tatsache. Die Frage, ob Kassen tatsächlich automatisiert nach Formfehlern screenen, ist empirisch nicht abschließend geklärt (siehe Abschnitt~\ref{sec:screening-rechtsfrage}). Die technischen Voraussetzungen dafür sind vorhanden; die operative Nutzung ist plausibel, aber nicht durch öffentliche Quellen belegt.

\subsubsection{Hypothese: Digitalisierung als Treiber der Einzelfallprüfungszunahme}

Die Zunahme der Einzelfallprüfungsanträge nach 2017 (Westfalen-Lippe +28\,\%, Nordrhein $\sim$20.000/Jahr) fällt zeitlich mit der zunehmenden Digitalisierung der Kassenverwaltung zusammen. Dies legt eine testbare Hypothese nahe: \textit{Je stärker eine Krankenkasse digitalisiert ist, desto mehr Einzelfallprüfungsanträge stellt sie} --- weil automatisierte Prüfsoftware die Identifikation von Formfehler-Kandidaten erleichtert und die Grenzkosten eines Antrags gegen Null senkt.

\textbf{Operationalisierung (Vorschlag):}
\begin{enumerate}[nosep]
  \item \textbf{Digitalisierungsgrad der Kassen:} Proxy über E-Rezept-Adoptionsrate, ePA-Nutzungsquote oder IT-Ausgaben pro Versichertem (KJ1-Statistik des BAS).
  \item \textbf{Einzelfallprüfungsvolumen pro Kasse:} Daten aus IFG-Anfragen oder KV-Statistiken, sofern verfügbar.
  \item \textbf{Korrelationsanalyse:} Korreliert der Digitalisierungsgrad mit dem Antragsvolumen? Kontrollvariablen: Kassengröße, regionale Prüfvereinbarung, Versichertenstruktur.
\end{enumerate}

Diese Analyse konnte im Rahmen der vorliegenden Studie nicht durchgeführt werden (fehlende Daten zu kassenspezifischen Antragsvolumina). Sie ist als Forschungsauftrag für die Zeit nach Eingang der IFG-Antworten vorgemerkt.

\subsubsection{Informationsasymmetrie: Der Prüfkatalog als Einbahnstraße}
\label{sec:pruefkatalog}

Zur operativen Synthese gehört ein Befund über Informationszugänge, der bisher nicht systematisch dargestellt wurde:

{\small
\begin{tabularx}{\textwidth}{l X l l}
\toprule
\textbf{Akteur} & \textbf{Zugang zum Prüfkatalog} & \textbf{Echtzeit-Prüfung} & \textbf{Prüfsoftware} \\
\midrule
Krankenkasse & Eigener Katalog (mit Anbietern wie scanacs entwickelt) & Ja (automatisiert) & Ja \\
\addlinespace
Apotheke & Prüfmodul im Warenwirtschaftssystem (scanacs, ABDA-Datenbank) & Ja (vor Abgabe) & Ja \\
\addlinespace
Arzt & Nur allgemeine Verordnungshinweise der KV & Nein & Nein \\
\bottomrule
\end{tabularx}
}

\medskip

Der Verordnende --- der Arzt --- ist der einzige Akteur im System, der die konkreten operativen Prüfkriterien nicht kennt, obwohl er derjenige ist, der die Verordnung ausstellt und für Formfehler haftet. Apotheken und Krankenkassen verfügen über Echtzeit-Prüfsoftware, die Verordnungen vor Abgabe bzw. nach Einreichung automatisiert gegen Prüfkataloge abgleicht. Ärzte haben keinen Zugang zu diesen Katalogen.

Das Bestimmtheitsgebot (Art.\,20 Abs.\,3 GG) verlangt, dass der Normadressat wissen kann, welches Verhalten von ihm erwartet wird. Die Prüfvereinbarungen nach \S\,106b Abs.\,6 SGB\,V sind öffentlich --- die operativen Prüfkataloge, nach denen konkret gescannt wird, sind es nicht. Ob diese Diskrepanz rechtsstaatlich problematisch ist, ist ungeklärt.

\subsubsection{Die offene Rechtsfrage: Verdacht oder Screening?}
\label{sec:screening-rechtsfrage}

\S\,106b Abs.\,1 SGB\,V definiert die Einzelfallprüfung als antragsgebundenes Verdachtsverfahren. Das setzt voraus, dass der Verdacht \textbf{zufällig} entsteht --- ein Sachbearbeiter stößt bei der regulären Arbeit auf eine auffällige Verordnung und formuliert daraus einen begründeten Antrag.

Systematisches, flächendeckendes, automatisiertes Screening aller Verordnungen mit dem Ziel, Formfehler-Kandidaten zu identifizieren, sieht das Gesetz nicht vor. Ein solches Vorgehen wäre funktional eine Auffälligkeitsprüfung (V1) --- eingekleidet als Einzelfallprüfung (V2).

{\small
\begin{tabularx}{\textwidth}{l X X}
\toprule
\textbf{Merkmal} & \textbf{Zufälliger Verdacht (gesetzeskonform)} & \textbf{Systematisches Screening (fragwürdig)} \\
\midrule
Auslöser & Sachbearbeiter bemerkt Auffälligkeit bei Routinearbeit & Software scannt alle Verordnungen \\
\addlinespace
Umfang & Einzelne Fälle & Flächendeckend, automatisiert \\
\addlinespace
Verdachtsqualität & Konkret, fallbezogen, menschlich bewertet & Maschinell generiert, menschlich nur bestätigt \\
\addlinespace
Analogie & Polizist beobachtet zufällig Verstoß & Flächendeckende Videoüberwachung \\
\bottomrule
\end{tabularx}
}

\medskip

Ob Kassen tatsächlich systematisch screenen, ist nicht belegt. Belegt ist, dass die technischen Mittel dafür existieren (Prüfsoftware, PZN$\leftrightarrow$ICD-Abgleich, \S\,300 SGB\,V-Datenlieferungen) und dass das Volumen der Einzelfallprüfungsanträge in einigen KV-Regionen in Größenordnungen liegt, die mit zufälligen Verdachtsmomenten schwer erklärbar sind (KV Nordrhein: ca. 20.000 Anträge/Jahr). Die Frage, ob ein automatisiert identifizierter Formfehler-Kandidat einen ``begründeten Antrag'' im Sinne des \S\,106b Abs.\,1 begründen kann --- oder ob die Automatisierung des Screening-Schritts die Verfahrensgrenze zwischen V1 und V2 verwischt ---, ist rechtlich ungeklärt.

\begin{sloppypar}
\textbf{Verfassungsrechtliche Einordnung:} Das BVerfG hat 2023 im ``Hessendata''-Urteil (1\,BvR\,1547/19) entschieden, dass jede automatisierte Datenanalyse ein \textbf{eigenständiger Eingriff} in das Recht auf informationelle Selbstbestimmung ist und strengen Verhältnismäßigkeitsanforderungen unterliegt. Die Übertragbarkeit auf die GKV-Prüfpraxis ist argumentativ gut vertretbar: Wenn die Polizei nicht unbegrenzt automatisiert Daten scannen darf, kann eine Krankenkasse nicht ohne weiteres sämtliche Verordnungsdaten automatisiert durchrastern, um daraus Regressanträge zu generieren. Eine gerichtliche Klärung dieser Frage existiert bisher nicht --- es handelt sich um eine echte Regelungslücke.
\end{sloppypar}

\textbf{Forschungsbedarf:} Eine IFG-Anfrage an die Krankenkassen, ob und in welchem Umfang automatisierte Software zur Identifikation von Einzelfallprüfungs-Kandidaten eingesetzt wird, könnte diesen Befund empirisch untermauern oder widerlegen. Die operativen Prüfkataloge selbst müssten nach \S\S\,1, 2, 70 SGB\,V dogmatisch offenlegungspflichtig sein --- ein BSG-Urteil hierzu existiert jedoch nicht.

\subsection{Der innere Konflikt: Drei Codes auf einer Operationseinheit}

Die Asymmetrie zwischen den Prüfverfahren ist nur eine äußere Schicht des Problems. Die innere Schicht liegt eine Ebene tiefer: in der Konstruktion des Wirtschaftlichkeitsgebots selbst.

\subsubsection{\S\,12 SGB\,V als Doppelanforderung}

\S\,12 Abs.\,1 SGB\,V verlangt: ``Die Leistungen müssen \textit{ausreichend, zweckmäßig und wirtschaftlich} sein; sie dürfen das Maß des Notwendigen nicht überschreiten.'' Die Norm kombiniert in einem einzigen Satz vier qualitative Anforderungen ohne Vorrangregel:

\begin{itemize}[nosep]
  \item ``ausreichend'' (Versorgungsanspruch des Patienten, \S\S\,27, 31 SGB\,V)
  \item ``zweckmäßig'' (Indikationsangemessenheit, leitliniengemäß, \S\,92 SGB\,V)
  \item ``wirtschaftlich'' (Kostenoptimum, sanktioniert durch \S\,106b SGB\,V)
  \item ``nicht das Maß des Notwendigen überschreitend''
\end{itemize}

Solange es um eine \textit{Leistungsdefinition} geht, ist die Kombination tragbar — sie beschreibt einen Ziel-Korridor. Sobald sie zur \textit{Sanktionsnorm} wird (über \S\,106b SGB\,V mit existenzbedrohenden Regressbeträgen), wird sie zum Problem: Der Arzt soll sämtliche Anforderungen simultan erfüllen — und trägt die Verantwortung wenn sie nicht alle gleichzeitig erfüllbar sind.

\subsubsection{Drei Funktionssystem-Codes auf einer ärztlichen Verordnung}

Wer das systemtheoretisch betrachtet (in der Tradition Niklas Luhmanns), erkennt: Auf der Operationseinheit ``ärztliche Verordnung'' treffen drei Funktionssystem-Codes zusammen, die nicht miteinander kompatibel sind:

{\small
\begin{tabularx}{\textwidth}{X X X}
\toprule
\textbf{Code} & \textbf{Funktionssystem} & \textbf{Anforderung an die Verordnung} \\
\midrule
krank / gesund & Medizin & medizinisch indiziert, leitliniengemäß \\
Zahlung / Nichtzahlung & Wirtschaft (GKV) & wirtschaftlich, kostenoptimiert \\
Recht / Unrecht & Recht (Sozialrecht) & substantiierungsfähig, formfehlerfrei, dokumentiert \\
\bottomrule
\end{tabularx}
}

Die drei Codes sind in modernen Gesellschaften normalerweise systemtheoretisch entkoppelt: Der Arzt entscheidet medizinisch, das Wirtschaftssystem kalkuliert separat, das Rechtssystem urteilt auf einer dritten Ebene. Im deutschen Vertragsarztrecht treffen alle drei Codes auf demselben Objekt zusammen — und auf derselben Person. Der Arzt muss \textit{simultan} medizinisch, ökonomisch und juristisch korrekt handeln. Ist eine der drei Anforderungen mit den anderen nicht vereinbar, trägt er das Risiko allein.

In Luhmanns Sprache: Was hier vorliegt, ist eine Kollision funktional differenzierter Codes auf einer Operationseinheit, ohne dass die Differenzierung selbst an der Operationseinheit aufgehoben wäre. Das ist mit den Mitteln moderner Gesellschaftsstrukturen nicht auflösbar — es muss durch institutionelle Konstruktion umgangen werden. Im deutschen Fall geschieht das nicht; die Konstruktion belastet stattdessen den einzelnen Arzt.

\subsubsection{Das KBV-Statement von 2016}

Die KBV hat den Doppelbind öffentlich benannt. Regina Feldmann (damals KBV-Vorstand) im Jahr 2016:

\begin{quote}
\textit{``Der Arzt soll verantwortlich sein für die Therapie, für Dosierung und Wirkstoffauswahl. Aber nicht für den Preis.''}
\end{quote}

Das ist genau die Forderung nach Code-Trennung: Medizinischer Code beim Arzt, ökonomischer Code beim System. Sie wurde von der Gesundheitspolitik nicht aufgegriffen.

\subsubsection{Der BVerfG-Notausgang von 2005}

Das Bundesverfassungsgericht hat den Konflikt erkannt — aber nur für einen schmalen Anwendungsfall. Im sogenannten \textit{Nikolaus-Beschluss} (Az. 1 BvR 347/98 vom 06.12.2005) entschied der Erste Senat: Bei \textbf{lebensbedrohlichen oder regelmäßig tödlichen Erkrankungen} und Fehlen einer alternativen Behandlung besteht ein Anspruch auf nicht regulär erstattete Behandlung — auch dann, wenn das Wirtschaftlichkeitsgebot dagegen sprechen würde. Diese Rechtsprechung wurde 2012 in \S\,2 Abs.\,1a SGB\,V kodifiziert.

Die Normalversorgung blieb davon unberührt. Wer chronisch krank ist, ohne lebensbedrohlich zu sein — Diabetiker, Schmerzpatienten, MS-Erkrankte, viele andere —, fällt nicht unter den Notausgang. Für sie greift weiterhin der volle Doppelbind, und der Arzt trägt das Risiko.

Sämtliche Reformen seit 1989 betrafen das \textit{Verfahren} der Prüfung (\S\,106 → \S\,106b, Beratung vor Regress, Bagatellgrenzen, Differenzkostenregelung 2024) — \textbf{nicht die Grundnorm \S\,12 SGB\,V.} Der innere Code-Konflikt blieb unangetastet.

\subsection{Der theoretische Rahmen: Managerialismus und Audit Society}

Die bisherigen Beobachtungen lassen sich in zwei etablierte Theorielinien der Sozial- und Verwaltungswissenschaft einbetten, die in der deutschen gesundheitspolitischen Diskussion erstaunlich wenig rezipiert sind. Beide liefern keine neuen Fakten, aber einen Rahmen, in dem die Einzelbefunde aufhören, als ``deutsche Eigentümlichkeit'' zu erscheinen.

\subsubsection{Freidson: Die ``dritte Logik'' und ihre Erosion}

Der amerikanische Medizinsoziologe Eliot Freidson (1923--2005) hat in drei Standardwerken --- \textit{Profession of Medicine} (1970), \textit{Professional Dominance} (1970) und \textit{Professionalism: The Third Logic} (Polity, Cambridge 2001) --- das Konzept der professionellen Dominanz entwickelt. Seine Kernthese: Echte Professionen unterscheiden sich von gewöhnlichen Berufen durch drei Merkmale --- hohes Maß an Autonomie, Selbstregulierung, und das Monopol auf die Definition der eigenen Arbeitsinhalte. Diese Autonomie sei gesellschaftlich legitimiert, weil Professionen über ein esoterisches Fachwissen verfügen, das Laien (auch staatliche Administratoren) nicht adäquat beurteilen können.

Freidson benennt drei konkurrierende Logiken der Arbeitsorganisation:

\begin{itemize}[nosep]
  \item der \textbf{freie Markt} (Consumerism), der über Preissignale steuert,
  \item die \textbf{bürokratische Kontrolle} (Managerialismus), die über formale Regeln und Rechenschaftspflichten steuert, und
  \item der \textbf{Professionalismus} --- die ``dritte Logik'' --- der über ethische Bindung und kollegiale Qualitätskontrolle steuert.
\end{itemize}

Die zentrale historische These Freidsons: In westlichen Gesundheitssystemen wird der Professionalismus seit den 1980er-Jahren systematisch durch den Managerialismus ersetzt. Das \textit{fundamentale Misstrauen} des Managerialismus lässt ärztliche Entscheidungen nicht mehr per se als gemeinwohlorientiert gelten; sie müssen externen, quantifizierbaren Leistungsindikatoren unterworfen werden.

Die in Teil II beschriebene rechtsdogmatische Lücke ist aus dieser Sicht keine Verirrung des deutschen Sozialrechts, sondern eine konsequente Anwendung des managerialen Prinzips: Der Arzt muss seine Entscheidungen formal auditierbar machen, sonst gilt seine Entscheidung vor dem System nicht. Die Drei-Code-Kollision des \S\,12 SGB\,V ist der deutsche Ausdruck eines globalen Paradigmenwechsels in der Steuerung von Gesundheitssystemen.

\subsubsection{Power: Die ``Audit Society'' und ihre Rituale der Verifikation}

Der britische Soziologe und Wirtschaftswissenschaftler Michael Power, Professor of Accounting an der London School of Economics, hat 1997 in \textit{The Audit Society: Rituals of Verification} (Oxford University Press) eine weitere, präzise zugeschnittene Theorieline vorgelegt. Powers Kernthese: Die massive Ausweitung von Prüfungs-, Zertifizierungs- und Rechenschaftspflichten seit den 1990er-Jahren führt nicht notwendig zu mehr Qualität, Effizienz oder Sicherheit --- sondern zu einer \textit{Illusion von Kontrolle}. In einer ``Audit Society'' richten Organisationen ihre Aufmerksamkeit und ihre Ressourcen weg vom eigentlichen Kernzweck hin zur Produktion formal prüfbarer Dokumente und Prozesse (des sogenannten ``audit trail'').

Drei Thesen Powers sind für den deutschen Fall direkt anwendbar:

\begin{itemize}[nosep]
  \item \textbf{Rituale der Verifikation:} Audits dienen primär der Legitimationsbeschaffung von Akteuren, nicht der Qualitätssicherung. Ihre Form ist wichtiger als ihr Inhalt.
  \item \textbf{Performativität:} Audits sind keine neutralen Messinstrumente, sie verändern das Verhalten der Geprüften aktiv --- sie schaffen die Realität, die sie zu messen vorgeben.
  \item \textbf{Entkopplung von Prozess und Ergebnis:} Was auditierbar ist, wird wichtig. Was nicht auditierbar ist, verliert an Gewicht --- unabhängig von seinem tatsächlichen medizinischen Wert.
\end{itemize}

Der in Teil II dokumentierte 491.000-EUR-Formalregress wegen eines Namensstempels (BSG B\,6\,KA\,9/24\,R) ist aus dieser Sicht kein bizarres Fehlurteil, sondern das paradigmatische Endstadium der Audit Society: Der tatsächliche medizinische Schaden ist null, aber in der Logik der ``Auditability'' kollabiert der institutionelle Wert der Handlung vollständig, wenn das Ritual der handschriftlichen Unterschrift fehlt. Und die in Teil III dokumentierte Defensivmedizin ist nicht bloß eine affektive Reaktion auf ein ``Angstsystem'' --- sie ist die zutiefst logische, strukturell erzwungene Adaption an eine Prüfgesellschaft, die Vertrauen in professionelle Expertise durch den Zwang zur dokumentierten Auditierbarkeit ersetzt hat.

\subsubsection{Noordegraaf: Das Paradox des konnektiven Professionalismus}

Mirko Noordegraaf (2015, 2020) hat Freidsons defensive Konzeption des Professionalismus weiterentwickelt. Er unterscheidet \textit{protective professionalism} (Verteidigung beruflicher Exklusivität und Autonomie gegen externe Eingriffe) von \textit{connective professionalism} (kooperative Integration von Stakeholder-Anforderungen --- wie Transparenz, Effizienz, Kostendruck --- in die eigene professionelle Praxis). Noordegraafs These: In modernen, komplexen Dienstleistungssystemen ist der rein defensive Modus nicht mehr tragfähig; Professionen müssen hybride Identitäten entwickeln.

Für das deutsche Wirtschaftlichkeitsprüfsystem ergibt sich daraus ein \textbf{theoretisches Paradox}: Konnektiver Professionalismus setzt voraus, dass das institutionelle Gegenüber Kooperationsanreize bietet. Ein System, das --- wie die Regelkreisanalyse (Abschnitt~3.8) zeigt --- ausschließlich Minus-Signale sendet und Formfehler deterministisch sanktioniert, entzieht dem Arzt systematisch die Möglichkeit zur vertrauensvollen Konnektivität. Die empirische Frage, ob und wie hybride Professionalität aufrechterhalten werden kann, wenn der Versuch zur Kooperation von der regulierenden Instanz durch formalistische Bestrafung beantwortet wird, ist eine offene Forschungslücke.

\subsubsection{Foucault: Die Wirtschaftlichkeitsprüfung als Panoptikum}

Die Verhaltenswirkung des Prüfsystems lässt sich mit Michel Foucaults Panoptismus-Konzept (\textit{Surveiller et punir}, 1975, dt. \textit{Überwachen und Strafen}) präzise beschreiben. Foucaults Kernthese: Ein Überwachungssystem entfaltet seine größte Disziplinierungswirkung nicht durch tatsächliche Bestrafung, sondern durch die \textit{permanente Möglichkeit} der Beobachtung. Der Überwachte muss nicht tatsächlich beobachtet werden --- er muss nur glauben, dass er beobachtet werden \textit{könnte}.

Das deutsche Wirtschaftlichkeitsprüfsystem erfüllt die drei Kernmerkmale des Panoptikums:

\begin{enumerate}[nosep]
  \item \textbf{Asymmetrische Sichtbarkeit:} Der Arzt weiß nicht, ob und wann seine Verordnungen geprüft werden. Die Kasse sieht seine Daten in Echtzeit (Verordnungsdaten-Informationssystem VODATI, KV-Abrechnungsdaten). Der Arzt sieht die Antragsstrategien der Kasse nicht.
  \item \textbf{Internalisierung der Kontrolle:} 72\,\% der Hausärzte (Ribbat 2023) ändern ihr Verordnungsverhalten \textit{präventiv} --- nicht weil sie geprüft werden, sondern weil sie geprüft werden \textit{könnten}. Die Disziplinierung wirkt ohne Ausübung.
  \item \textbf{Ökonomie der Macht:} Die öffentlich sichtbare V1-Regressquote ($<$1\,\%) ist kein Systemversagen, sondern maximale Effizienz: Das Screening allein erzeugt flächendeckende Verhaltensänderung, ohne viele Regresse durchsetzen zu müssen. Bei V2 ist die tatsächliche Durchsetzungsrate höher (Vergleiche), aber unsichtbar --- der Disziplinierungseffekt wirkt auf beiden Ebenen.
\end{enumerate}

\noindent Die Panoptikum-Perspektive ergänzt Power (Audit-Rituale als Ersatz für Vertrauen) und Freidson (Managerialismus verdrängt Professionalismus) um die Machtdimension: Das Prüfsystem ist nicht nur ein schlechtes Messinstrument --- es ist ein hochwirksames Disziplinierungsinstrument, das gerade \textit{weil} es nicht präzise misst, seine größte Wirkung entfaltet. Präzise Messung würde die meisten Ärzte entlasten. Diffuse Überwachung diszipliniert alle.

Freidson, Power und Foucault erklären, warum das System so aussieht und so wirkt, wie es wirkt. Sie erklären noch nicht, wie Ärzte darauf reagieren. Das ist die Frage von Teil III --- und dort wird eine weitere Theorielinie nötig, die den Arzt nicht als passives Opfer, sondern als rationalen Akteur unter unauflösbaren Zwangslagen versteht: Michael Lipskys ``Street-Level Bureaucracy''.


\section{Teil III --- Die Verhaltensökonomie: Was Ärzte tatsächlich tun}

\subsection{Ribbat 2023 --- die derzeit umfangreichste deutsche Befragung}

Die Frage, wie sich die Regressangst auf das ärztliche Verhalten auswirkt, ist nicht im Bereich der Anekdoten anzusiedeln. Sie ist empirisch beantwortet --- in einer bundesweiten Studie, die merkwürdigerweise kaum jemand kennt:

\textbf{Ribbat L, Linde K, Schneider A, Riedl B (2023): Effects of the Risk of Recourse Claims on the Work of General Practitioners and Orthopedists --- A Nationwide Survey. Gesundheitswesen 85(2): 111-118. DOI: 10.1055/a-1594-2527, PMID: 34587633.}

\textbf{Methodik:} Bundesweite Zufallsstichprobe von je 1000 Hausärztinnen und 1000 Orthopäd:innen. Rücklauf 41\,\% / 39\,\% (n=770). Vierseitiger Fragebogen. Lehrstuhl für Allgemeinmedizin TU München.

\textbf{Verifizierte Hauptergebnisse} (alle Werte V1+V2 zusammen, keine Verfahrenstrennung in der Studie)\textbf{:}

\begin{itemize}[nosep]
  \item \textbf{47\,\% Hausärzte / 55\,\% Orthopäden:} Regressgefahr beschäftigt im Praxisalltag ``stark oder sehr stark''.
  \item \textbf{37\,\% / 47\,\%:} ärztliches Handeln ``stark beeinflusst''.
  \item \textbf{51\,\% Hausärzte / 25\,\% Orthopäden:} überweisen ``zumindest gelegentlich trotz Indikation'' zu Spezialisten.
  \item \textbf{72\,\% / 78\,\%:} emotionale Belastung durch Regresse ``stark oder sehr stark''.
\end{itemize}

\textbf{Interpretation:} Knapp die Hälfte der Hausärzte ändert ihr Verordnungsverhalten aus Regressangst, mehr als die Hälfte schiebt Patienten weiter, die sie selbst behandeln könnten. Direkte deutsche Evidenz, keine Hochrechnung nötig.

\noindent\textit{Hinweis zur Gegenposition:} Die methodischen Grenzen dieser Befunde --- insbesondere die Frage, ob selbstberichtetes Verhalten auf tatsächliche Unterversorgung schließen lässt --- werden in Abschnitt~\ref{sec:gegenposition} (Einwand~1) diskutiert. Zur Operationalisierung und Falsifizierbarkeit siehe den folgenden Abschnitt.

\textbf{Datum am Rande --- die fehlende Resonanz:} Eine eigene Recherche zur Medienresonanz dieser Studie zeigt: Sie wurde von keinem deutschen Fachmedium prominent aufgegriffen --- weder vom Deutschen Ärzteblatt noch von der Ärzte Zeitung, weder von Medical Tribune noch von Coliquio. Auch in den Publikumsmedien (Zeit, Spiegel, FAZ) ist keine Berichterstattung nachweisbar. Bundestagsdebatten zitieren sie nicht. Die Verbände (KBV, NAV-Virchow, Hartmannbund) beklagen ``Regressangst'' chronologisch, stützen sich dabei aber auf eigene Verbandsumfragen und nennen Ribbat nicht als Goldstandard-Beleg. Es liegt kein Boykott vor und keine Verschwörung --- es ist das ``normale Schicksal'' einer versorgungswissenschaftlichen Studie ohne PR-Maschine. Genau diese Stille ist methodisch der zweite Befund: Die deutsche Versorgungsforschung produziert Befunde, die in der gesundheitspolitischen Debatte nicht ankommen. Das ist nicht das Versagen der Forschung, sondern eine Eigenschaft der Aufmerksamkeitsökonomie des deutschen Gesundheitssystems.

\subsection{Ribbat-Operationalisierung: Von der Befragung zum Verhaltensmechanismus}

\textbf{Einzelstudien-Risiko.} Ribbat et al.\ (2023) ist die einzige große deutsche Befragungsstudie, die spezifisch die Auswirkungen der Regressangst auf das ärztliche Verordnungsverhalten misst. Eine unabhängige Replikation steht aus. Diese Studie zitiert Ribbat als den \textit{besten verfügbaren Datenpunkt}, nicht als abschließenden Beweis. Die methodische Einschränkung --- dass Ribbat \textit{subjektive Wahrnehmung und selbstberichtetes Verhalten} misst, nicht objektive Versorgungsergebnisse --- wird in Abschnitt~\ref{sec:gegenposition} (Einwand~1) eingehend diskutiert.

\textbf{Mapping zum Theorierahmen.} Die Ribbat-Befunde lassen sich auf die in dieser Studie verwendeten theoretischen Rahmenwerke abbilden. Diese Zuordnung macht die empirische Verankerung der Theorie explizit:

\begin{tabular}{p{4cm}p{3.5cm}p{5cm}}
\textbf{Ribbat-Befund} & \textbf{Theorie} & \textbf{Vorhergesagter Mechanismus} \\
\hline
72\,\% haben Regress erlebt & Spieltheorie (Teil~III) & Hohes $P_{\text{erfahren}}$ $\rightarrow$ überhöhtes $P_{\text{perceived}}$ $\rightarrow$ defensives Nash-Gleichgewicht \\
51\,\% überweisen trotz Kompetenz & Lipsky: Coping-Strategie & Weiterverweisung als Risikoverlagerung (``Outsourcing'' des Regressrisikos) \\
85\,\% berichten Verordnungsverzicht & Lipsky + Herd/Moynihan & Rationierung (Lipsky) + Administrative Burden als Verhaltenssteuerung (Herd/Moynihan) \\
47\,\% tägliche Beschäftigung & Foucault: Panoptikum & Internalisierte Überwachung ohne aktive Beobachtung; die Kontrolle wirkt auch ohne ausgelöste Prüfung \\
\end{tabular}

\vspace{0.3em}
\noindent Dieses Mapping zeigt: Die vier zentralen Ribbat-Befunde sind nicht isolierte Beobachtungen, sondern lassen sich als verschiedene Manifestationen derselben Systemwirkung lesen --- einer punitiven Steuerungslogik, die rationale Coping-Strategien erzwingt.

\textbf{Was würde die These entkräften?} Zwei empirische Befunde würden die Kausalkette erheblich schwächen: (1)~Eine Replikation von Ribbat mit weniger als 20\,\% berichtetem Verhaltenseffekt würde die Größenordnung der subjektiven Belastung infrage stellen. (2)~Eine Korrelationsanalyse zwischen Prüfintensität (Antragshäufigkeit der Kassen je KV-Region) und Verordnungsvolumen (DDD pro Kopf), die keinen Zusammenhang zeigt, würde die These einer systemischen Verhaltenssteuerung schwächen. Beide Datensätze sind derzeit nicht öffentlich verfügbar --- die hierzu eingereichten IFG-Anfragen (Anhang~C) zielen auf genau diese Lücke.

\subsection{Die internationale Bestätigung: 75,8\,\% Defensivmedizin global}

Die Ribbat-Befunde stehen nicht allein. Eine globale Meta-Analyse (\textbf{Zheng et al. 2023}, Int J Qual Health Care, PMID: 38060672) beziffert die globale Prävalenz von Defensivmedizin auf 75,8\,\%. Die deutsche Spezialstudie \textbf{Goetz et al. 2024} (BMC Primary Care, DOI: 10.1186/s12875-024-02267-x) zeigt: 27\,\% der deutschen Hausärzte mit ``starker Rechtsangst'', 54\,\% wöchentliche unnötige Labortests, 38\,\% Selbstschutz als Motiv.

Die Wirkmechanik ist also weder selten noch deutsch. Sie ist global, sie ist bekannt --- und sie ist im deutschen Fall durch ein spezifisches Sanktionssystem verschärft, das in dieser Form nur Deutschland und die Schweiz kennen (siehe Teil VI).

\subsection{Spieltheoretische Modellierung}

Wer das Regress-System spieltheoretisch betrachtet, findet drei interessante Strukturen: ein dominantes Spiel mit asymmetrischer Information, ein Nash-Gleichgewicht das alle Beteiligten schlechter stellt, und ein Gefangenendilemma in der Ärzteschaft. Die folgenden Abschnitte skizzieren die Grundstruktur; eine formale Durchmodellierung mit Auszahlungsmatrizen, Sensitivitätsanalysen und verfahrensspezifischen Parametern ist als eigenständige Folgeanalyse vorgesehen.

\subsubsection{Das Verordnungsspiel: Bayessches Spiel mit asymmetrischer Information}

Das zentrale Spiel wird zwischen Arzt und Krankenkasse millionenfach pro Quartal gespielt. Es ist ein \textbf{Bayessches Spiel}, weil beide Seiten unter Unsicherheit über die Wahrscheinlichkeit eines Regresses entscheiden müssen — und weil diese Wahrscheinlichkeit \textit{asymmetrisch verteilt} ist: Die Kasse kennt ihre eigenen Antragsstrategien, die internen Erfolgsquoten und teilweise die vertraulichen Erstattungsbeträge. Der Arzt kennt nichts davon. Er schätzt das Regressrisiko über Anekdoten, Medienberichte und Berufsbiografie.

\textbf{Grundgesamtheiten:} Die Risikobetrachtung hängt fundamental von der gewählten Bezugsgröße ab. Bei V2 (Einzelfallprüfung) ist die Wahrscheinlichkeit, dass \textit{eine bestimmte Verordnung} geprüft wird, verschwindend gering (ca.\,800\,Mio. GKV-Verordnungen pro Jahr, tausende V2-Anträge). Für den einzelnen \textit{Arzt} mit durchschnittlich ca.\,1.380 Patienten pro Quartal (KBV Gesundheitsdaten) und mehreren tausend Verordnungen pro Jahr summiert sich das Risiko:

\textbf{Die Summierungsrechnung (KV Nordrhein als Beispiel):}

{\small
\begin{tabularx}{\textwidth}{p{6cm} X}
\toprule
\textbf{Kenngröße} & \textbf{Wert} \\
\midrule
Praxen in KV\,Nordrhein & ca.\,13.500 \\
Vertragsärzte & ca.\,22.500 ($\sim$1,67 pro Praxis) \\
V2-Anträge pro Jahr & ca.\,20.000 (davon $\sim$10.000 Sprechstundenbedarf) \\
V2-Anträge pro Arzt und Jahr & $\sim$0,9 (bei Gleichverteilung$^\dagger$) \\
\addlinespace
V2-Durchsetzungsquote & \textbf{$\sim$95\,\%} führen zu Regresszahlung (Dt.\,Ärzteblatt 2020) \\
$\to$ Zahlungsvorgänge pro Arzt/Jahr & $\sim$0,85 (bei Gleichverteilung$^\dagger$) \\
\addlinespace
\textbf{Jährliche Betroffenheitsquote} & \textbf{$\sim$57\,\%} der Ärzte mit mind. 1 Zahlungsvorgang$^{\dagger\dagger}$ \\
$\to$ In 3 Jahren kumuliert & $\to$ $\sim$92\,\% mindestens einmal betroffen \\
\addlinespace
\textbf{Ribbat 2023: 72\,\% Regresserfahrung} & \textbf{Rechnerisch plausibel} \\
\bottomrule
\end{tabularx}

$^\dagger$\,Annahme: Die 20.000 V2-Anträge verteilen sich gleichmäßig auf alle 22.500 Ärzte. Die tatsächliche Verteilung ist nicht publiziert --- genau das ist Teil der diagnostizierten Intransparenz. Bei Konzentration auf bestimmte Fachgruppen oder Praxen wäre die individuelle Betroffenheitsquote für diese höher, für andere niedriger.\\
$^{\dagger\dagger}$\,Poisson-Modell: Bei $\lambda = 0,85$ Ereignissen pro Jahr ist die Wahrscheinlichkeit für $k=0$ (kein Vorgang): $P(0) = e^{-0,85} \approx 0,43$. Also: $1 - 0,43 = 0,57$ = 57\,\% der Ärzte mit mindestens einem Vorgang pro Jahr. Über 3 Jahre: $1 - e^{-3 \times 0,85} \approx 0,92$.
}

\noindent\textbf{Kleines Risiko pro Verordnung, großes Risiko pro Arzt.} Bei Gleichverteilung der V2-Anträge haben \textbf{57\,\% der Ärzte pro Jahr} mindestens einen Zahlungsvorgang (Regress oder Vergleich). Über drei Jahre kumuliert: 92\,\%. Die ``unter 1\,\%''-Zahl misst etwas völlig anderes --- nur V1-Regresse, nur rechtskräftig. Die V2-Vergleichsrealität taucht in keiner öffentlichen Statistik auf.

\textbf{Die Kommunikationsasymmetrie:} Beide Seiten sagen die Wahrheit --- über verschiedene Messgrößen. Die ``unter 1\,\%''-Zahl (Virchowbund: 0,065\,\%, basierend auf $<$100 rechtskräftigen Regressen bundesweit 2018 bei 155.000 Vertragsärzten) misst \textit{ausschließlich V1-Regresse} --- der Virchowbund selbst sagt, statistische Auffälligkeitsprüfungen ``spielen kaum noch eine Rolle''. Die KV\,BaWü-Zahl (0,097\,\%) misst V1+V2, aber nur \textit{rechtskräftig durchgesetzte} Regresse. In beiden Fällen fehlen die \textbf{V2-Vergleiche} --- geschätzt 70--80\,\% der betroffenen Ärzte zahlen, ohne dass ein ``Regress'' in der Statistik erscheint.

Der Arzt misst seine \textit{Berufslaufbahn}: Wie oft war ich in 20 Jahren mit einem V2-Verfahren konfrontiert? Bei $\sim$0,9 Anträgen pro Arzt und Jahr (KV Nordrhein) und 95\,\% Durchsetzungsquote: mehrfach. Es ist kein Widerspruch --- es ist eine systematische Perspektivdivergenz, bei der die beruhigende Zahl immer die öffentlich kommunizierte ist und die belastende Zahl in keiner Statistik auftaucht.

\textbf{Bundesweit:} 2022 gab es ca.\,47.000 Wirtschaftlichkeitsprüfungen bei niedergelassenen Ärzten (BMG GVSG-Referentenentwurf 2024). Bei ca.\,155.000 Vertragsärzten entspricht das einer Rate von ca.\,30\,\% pro Jahr --- wobei der Großteil Kleinbetragsprüfungen betrifft. \textbf{Regionale Bestätigung:} Der SpiFa meldet für Bayern 13.332 festgesetzte Regresse (2024) bei $\sim$24.000 Vertragsärzten --- Durchschnittsbetrag 334\,EUR, $\sim$70\,\% davon Bagatellregresse unter 300\,EUR (SpiFa 2024). Die bayerische Rate (max. 55\,\% bei Einzelzählung, $\sim$28\,\% bei Zwei-pro-Arzt-Bündelung) konvergiert mit dem Bundesschnitt und dem Poisson-Modell.

Die empirisch belegbaren Werte:

\begin{itemize}[nosep]
  \item \textbf{Offizielle Regressdurchsetzungsquote} (P\textsubscript{offiziell}, V1+V2 zusammen): unter 1\,\%, je nach KV-Bezirk 0,065--0,8\,\%
  \item \textbf{Tatsächliche Regresserfahrung} (P\textsubscript{erfahren}, Ribbat 2023): \textbf{72\,\% der Hausärzte} haben mindestens einmal einen Regress erlebt (V1+V2 zusammen; keine Verfahrenstrennung)
  \item \textbf{Wahrgenommenes Risiko} (P\textsubscript{perceived}): rund 15\,\% (KBV-Berufsmonitoring zeigt 50\,\% der Studierenden mit Niederlassungsangst)
\end{itemize}

\textbf{Die Diskrepanz liegt nicht zwischen Wahrnehmung und Realität, sondern zwischen offizieller Statistik und tatsächlicher Erfahrung.} Der frühere Wahrnehmungsfaktor von ca.\,50:1 (P\textsubscript{perceived}\,/\,P\textsubscript{offiziell}) suggerierte, Ärzte überschätzten ihr Risiko dramatisch. Ribbats 72\,\% zeigen: Die ärztliche Wahrnehmung ist näher an der Realität als die offizielle Statistik. Die Erklärung: Die offizielle ``unter 1\,\%''-Zahl misst nur \textit{rechtskräftig durchgesetzte} Regresse (Schritt~9), nicht die Gesamtheit der Verfahren, Anhörungen, Vergleiche und psychischen Belastungen. Beide Größen --- die offizielle Quote und die Ribbat-Erfahrungsrate --- sind V1+V2 zusammen; eine verfahrensgetrennte Messung existiert nicht. Die Angst ist berechtigt --- sie bezieht sich auf die Verfahrenserfahrung, nicht auf den Geldfluss.

Mit der Verlustaversions-Konstanten von Kahneman/Tversky ($\lambda \approx 2{,}25$) ergibt sich zusätzlich ein verstärkter gefühlter Erwartungswert. Der Arzt entscheidet nicht unter falschen Annahmen, sondern unter Berücksichtigung einer Verfahrensrealität, die in der offiziellen Statistik systematisch untererfasst ist.

\subsubsection{Defensive Unter-Verordnung als Nash-Gleichgewicht}

Bei P\textsubscript{perceived} $\approx$ 0,15 wird ``nicht verordnen'' für den Arzt zur dominanten Strategie. Bei P\textsubscript{real} $\approx$ 0,003 wäre ``verordnen'' dominant. Die Schwelle, an der die Strategie kippt, ist mathematisch berechenbar — und sie liegt deutlich unter der wahrgenommenen Risikoschwelle.

Das Ergebnis ist ein \textbf{Pareto-suboptimales Nash-Gleichgewicht}:

\begin{itemize}[nosep]
  \item Arzt verliert (ethischer Stress, defensive Mehrarbeit, dokumentationslastige Prozesse)
  \item Patient verliert (Unterversorgung, vermeidbare Folgeschäden)
  \item Krankenkasse gewinnt nichts (das nicht verordnete Medikament wird ohnehin nicht verordnet, sie zahlt nichts ein, aber sie zahlt später die Hospitalisierung)
  \item System verliert die in Teil IV bezifferten Folgekosten
\end{itemize}

Niemand will dieses Gleichgewicht. Alle spielen es individuell rational.

\subsubsection{Das Solidaritäts-Gefangenendilemma}

Das zweite zentrale Spiel ist nicht das Verordnungsspiel selbst, sondern das Spiel \textit{daneben}: das Solidaritäts-Spiel innerhalb der Ärzteschaft. Zwei Ärzte könnten gemeinsam öffentlich für leitliniengerechte Verordnung kämpfen, koordinierte Beschwerden einreichen, gemeinsam Pressearbeit machen. Jeder einzelne hat aber den Anreiz, still zu bleiben und sich selbst zu schützen.

\begin{itemize}[nosep]
  \item Wer öffentlich kämpft, wird zur Zielscheibe künftiger Prüfverfahren
  \item Wer schweigt, profitiert vom Erfolg anderer (Trittbrettfahrer in der Solidaritäts-Allmende)
\end{itemize}

Das Ergebnis: Berufsverbände sind politisch schwach, individuelle Erfahrungen bleiben anekdotisch, statt aggregiert zu werden. Genau diese Aggregation wäre die Bedingung für eine kollektive Verteidigungsfront — und sie findet aus rationalen individuellen Gründen nicht statt.

\subsubsection{Multi-Principal-Agent-Problem}

Auf einer dritten Ebene liegt eine klassische Principal-Agent-Konfiguration mit drei überlagerten Schichten:

\begin{itemize}[nosep]
  \item Schicht 1: \textbf{Patient $\to$ Arzt} — der Patient erwartet Versorgung
  \item Schicht 2: \textbf{Krankenkasse $\to$ Arzt} — die Kasse erwartet wirtschaftliches Verordnen
  \item Schicht 3: \textbf{Versicherter $\to$ Krankenkasse} — der Versicherte erwartet, dass die Kasse seine Versorgung organisiert
\end{itemize}

Das systemische Problem: \textbf{Schicht 2 dominiert das Verhalten des Arztes, Schicht 3 ist faktisch inaktiv.} Der Versicherte hat keinen direkten Einfluss darauf, wie aggressiv ``seine'' Kasse Regresse beantragt. Der Patient erlebt nur die Konsequenz (Unterversorgung), nicht die Ursache (Antragsdruck der Kasse). Die Kette ist gebrochen.

\subsubsection{Die spieltheoretisch wirksamste Intervention}

Der mathematische Befund liefert eine klare Empfehlung: \textbf{Reduktion der asymmetrischen Information ist die wirkungsmächtigste Einzelintervention.} Sobald Ärzte ihr reales Regressrisiko statt ihres empfundenen Risikos kalibrieren können, verschiebt sich das Verordnungsspiel automatisch vom defensiven zum kooperativen Gleichgewicht. Das ist eine echte Pareto-Verbesserung: Patient gewinnt, Arzt gewinnt, System gewinnt; die Kasse verliert nur ihre Abschreckungs-Rente, nicht ihre tatsächlichen Regress-Einnahmen.

\textit{Information schlägt Reform — weil Reform der Information folgt, nicht umgekehrt.}

Genau auf diese Funktion zielt das Regress-Transparenzportal in PP-003. \textbf{Eine verfahrensgetrennte Erhebung (V1 vs. V2) wäre dabei das wirksamste Korrektiv für das bisherige Narrativ:} Sollte PP-003 zeigen, dass die V2-Regressdurchsetzungsquote deutlich über der offiziellen V1+V2-Mischquote liegt, würde das die Risikokalibrierung der Ärzte rückwirkend als realistischer erweisen --- und die Dringlichkeit gesetzlicher Korrekturen erhöhen. Die fehlende Verfahrenstrennung in den offiziellen Statistiken ist daher nicht nur ein methodisches Problem, sondern ein politisches: Sie verhindert die Identifikation des asymmetrischeren Verfahrens (V2) als eigenständiges Regulierungsproblem.

\subsubsection{Regelkreisanalyse: Nur Minus-Signale}

Aus steuerungstechnischer Perspektive weist das Prüfsystem einen fundamentalen Konstruktionsfehler auf: Es enthält ausschließlich negative Rückkopplungssignale (Sanktionen bei Verstößen), aber keine positiven (Belohnung oder Bestätigung bei korrektem Verhalten).

{\small
\begin{tabularx}{\textwidth}{l X X}
\toprule
\textbf{Signaltyp} & \textbf{Beispiel im Regresssystem} & \textbf{Beispiel in einem balancierten System} \\
\midrule
Minus-Signal (Sanktion) & Regress, Auffälligkeitsmeldung, Vergleichszahlung & Konventionalstrafe \\
Plus-Signal (Bestätigung) & \textit{existiert nicht} & Prämie für leitliniengerechte Versorgung, positives Prüfergebnis \\
Null-Signal (neutral) & Keine Prüfung = kein Feedback & Keine Prüfung = kein Feedback \\
\bottomrule
\end{tabularx}
}

\noindent Ein Regelkreis, der nur Minus-Signale sendet, erzeugt systemtheoretisch vorhersagbar ein Über\-vorsichtiges Gleich\-gewicht: Das System bestraft Abweichung, belohnt aber Konformität nicht. Der Arzt erhält nie die Rückmeldung ``Ihre Verordnungspraxis ist medizinisch angemessen'' --- er erfährt nur ``Sie wurden nicht bestraft''. Die Abwesenheit von Bestrafung ist kein positives Signal. Sie ist ein Null-Signal, das der Arzt nicht von Glück unterscheiden kann.

\subsection{Der Modellarzt}

Aus den verfügbaren empirischen Daten lässt sich ein konkretes Profil des regress-ängstlichen Modellarzts ableiten. Es ist keine Konstruktion, sondern die Verdichtung der dokumentierten Befunde.

\subsubsection{Profil}

Mittleres Lebensalter (ca. 50 Jahre), Hausarzt in westdeutscher Mittelstadt, seit etwa 18 Jahren in Einzelpraxis, eine MFA, rund 1.100 Scheine pro Quartal, hoher Anteil multimorbider Patienten. Hat bereits ein bis drei Auffälligkeitsmeldungen der KV erlebt.

\textbf{Wichtig:} Die statistisch belegten Prädiktoren der Regressangst sind \textit{nicht} Alter, Berufserfahrung oder Praxisstruktur (Goetz et al. 2024 zeigen explizit: nicht signifikant). Prädiktiv sind:
\begin{itemize}[nosep]
  \item Wunsch nach rechtlicher Selbstabsicherung ($\beta = 0{,}338$)
  \item Antizipierte Klagewahrscheinlichkeit ($\beta = 0{,}250$)
  \item Geschlecht (Frauen ängstlicher, $\beta = -0{,}158$)
\end{itemize}

\textbf{Die Angst ist also strukturell, nicht biografisch.} Sie entsteht aus dem Spielfeld, nicht aus den Spielern. Strukturell exponiert sind neben Hausärzten besonders Onkologen, Diabetologen, Schmerzmediziner und Neurologen.

\subsubsection{Die dokumentierten Verhaltensanpassungen}

Was tut der Modellarzt anders? Die folgenden Verhaltensweisen sind alle empirisch belegt:

\begin{itemize}[nosep]
  \item \textbf{Quartalsende-Verschiebung} teurer Verordnungen (Himmel/Schneider 2017, HCHE Research Paper No.\,2017/14, TK-Daten: bis 14\,\% weniger Termine, +19 Prozentpunkte Notdienstbelastung)
  \item \textbf{Überweisung statt Behandlung} (Ribbat 2023: 51\,\% der Hausärzte schieben weiter, obwohl sie selbst behandeln könnten)
  \item \textbf{Verordnungsverzicht} im Lauf der Karriere (Ribbat 2023: 85\,\% der Hausärzte haben mindestens einmal aus Regressangst eine Verordnung unterlassen)
  \item \textbf{Defensive Diagnostik} (Goetz 2024: 54\,\% wöchentlich unnötige Labortests)
  \item \textbf{Off-Label-Vermeidung} trotz medizinischer Indikation
  \item \textbf{Heilmittel-Drosselung} (Reha-Verzögerung)
  \item \textbf{Patientenselektion} gegen multimorbide Patienten
  \item \textbf{Verlagerung in den stationären Sektor} (treibt die deutsche ACSC-Hospitalisierungsrate)
  \item \textbf{Niederlassungsverzicht} (KBV-Berufsmonitoring: 50\,\% der Studierenden nennen Regress als Hindernis)
\end{itemize}

\subsubsection{Hochrechnung auf die Makro-Ebene}

Bei \textbf{155.678 Vertragsärzten} (KBV 2025, ohne Psychotherapeuten), 40\,\% betroffenen Ärzten (Mittelwert aus Ribbat-Befragungswerten für Hausärzte und Orthopäden), 20 unterlassenen Verordnungen pro Arzt und Jahr und konservativen Folgeannahmen ergibt sich ein direkt zurechenbares Schadensvolumen von etwa \textbf{200 Mio.\,EUR/Jahr} aus den Verhaltensanpassungen allein.

Rechnet man den plausiblen Regressangst-Anteil an den vermeidbaren ACSC-Folgekosten (7-8\,Mrd.\,EUR/Jahr WIdO/Zi) und Non-Adhärenz-Folgekosten (10-19\,Mrd.\,EUR DIMDI HTA 65) mit konservativen 5\,\% dazu, liegt die Gesamtspanne bei \textbf{0,9-1,7\,Mrd.\,EUR/Jahr}.

Unter diesen Szenario-Annahmen (mittleres Modell) wären ca. \textbf{6.230 vermeidbare Hospitalisierungen} und ca. \textbf{125 Todesfälle pro Jahr} anteilig mit Regressangst-induzierter Unterversorgung assoziierbar. \textbf{Wichtige Einschränkung:} Diese Zahlen sind ein illustratives Gedankenexperiment auf Basis der Szenario-Annahmen, kein kausalanalytisch abgesicherter Befund. Die tatsächliche Mortalitätslast ist ohne eine eigenständige epidemiologische Studie nicht bestimmbar.

Der gesamtsystemische Rahmen liegt bei rund 15.000--20.000 vermeidbaren Todesfällen pro Jahr in Deutschland (Schätzungen aus verschiedenen Quellen, u.\,a. Sepsis-Memorandum der Leopoldina/Fachgesellschaften; Quelle für v1.0 zu präzisieren), von denen nur ein Teil regress-kausal ist. Die hier vorgenommene Hochrechnung ist absichtlich konservativ.

\subsubsection{Methodische Einschränkungen und Sensitivitätsanalyse}

Die Hochrechnung beruht auf Szenario-Annahmen, nicht auf empirisch gemessenen Parametern. Insbesondere die Attributionsrate (Anteil regress-kausal an Folgekosten) und die Frequenz unterlassener Verordnungen sind Modell-Annahmen --- keine der derzeit verfügbaren Studien misst diese Größen direkt. Die Spanne der Szenarien soll die Größenordnung des möglichen Schadens illustrieren, nicht einen Punktschätzer liefern. Drei Schlüsselparameter bestimmen die Bandbreite:

{\small
\begin{tabularx}{\textwidth}{X X X X}
\toprule
\textbf{Parameter} & \textbf{Konservativ} & \textbf{Mittel} & \textbf{Pessimistisch} \\
\midrule
Anteil betroffener Ärzte (Ribbat 2023) & 27\,\% (Goetz: starke Rechtsangst) & 40\,\% (Mittelwert HA/Ortho) & 75,8\,\% (Zheng global) \\
Unterlassene Verordnungen pro Arzt/Jahr & 5 & 20 & 50 \\
Anteil regress-kausal an Folgekosten & 2\,\% & 5\,\% & 15\,\% \\
\midrule
\textbf{Direkt zurechenbar} & ca. 25\,Mio.\,EUR & ca. 200\,Mio.\,EUR & ca. 1,4\,Mrd.\,EUR \\
\textbf{Plus anteilige Folgekosten} & ca. 0,4\,Mrd.\,EUR & ca. 0,9\,Mrd.\,EUR & ca. 4,2\,Mrd.\,EUR \\
\textbf{Vermeidbare Hospitalisierungen/Jahr} & $\sim$2.500 & $\sim$6.230 & $\sim$23.600 \\
\textbf{Vermeidbare Todesfälle/Jahr} & $\sim$50 & $\sim$125 & $\sim$472 \\
\bottomrule
\end{tabularx}
}

\textbf{Lesehinweis:} Die Bandbreite zeigt, dass selbst unter sehr konservativen Annahmen erhebliche Schadensvolumen plausibel sind. Im mittleren Szenario liegen die direkten plus anteiligen Kosten bei rund \textbf{1,1\,Mrd.\,EUR/Jahr} (200\,Mio. direkt + 0,9\,Mrd. anteilig). Der in Teil\,IV eingeführte Wert von 7--17\,Mrd.\,EUR ist der \emph{Systemrahmen} der assoziierten ACSC- und Non-Adhärenz-Folgekosten vor Attribution~--- der Nenner, aus dem durch Multiplikation mit der Attributionsrate die Regress-Folgekosten berechnet werden.

Die Modellrechnung ersetzt keine kausalanalytische Studie. Sie zeigt nur, dass die Größenordnung des Schadens unter plausiblen Annahmen die Größenordnung der Einnahmen um ein Vielfaches übertrifft.

Eine systematische Prüfung durch eine zuständige Stelle (BRH, SVR, IGES, Versorgungsforschungsinstitute) hat bisher nicht stattgefunden. Auch das ist ein Befund.

\subsection{Der Arzt als Street-Level Bureaucrat}

Die bisherigen Teile dieser Studie haben das Verhalten der Ärzte primär über verhaltensökonomische Kategorien erklärt: Verlustaversion nach Kahneman/Tversky, asymmetrische Information, Nash-Gleichgewicht. Diese Erklärung ist mathematisch elegant, aber sie reicht nicht aus. Sie beschreibt, \textit{warum} der einzelne Arzt eine bestimmte Entscheidung trifft, nicht \textit{was für einen Akteur er dabei ist}. Eine präzisere Einordnung liefert die Verwaltungswissenschaft, namentlich Michael Lipskys Theorie der Street-Level Bureaucracy.

\subsubsection{Lipskys Grundmodell}

In seinem Standardwerk \textit{Street-Level Bureaucracy: Dilemmas of the Individual in Public Services} (Russell Sage Foundation, 1980; erweiterte Neuauflage 2010) definiert Lipsky die sogenannten \textbf{Street-Level Bureaucrats} (SLB) als Akteure an der äußersten Frontlinie staatlicher oder quasi-staatlicher Leistungserbringung: Lehrer, Polizisten, Sozialarbeiter, Sachbearbeiter --- und im öffentlich finanzierten Gesundheitssystem auch Ärzte. Sie teilen drei Merkmale:

\begin{itemize}[nosep]
  \item \textbf{Direkter Kontakt} zu den Bürgern/Patienten
  \item \textbf{Erheblicher Ermessensspielraum} in der Einzelfallentscheidung
  \item \textbf{Unauflösbare Zwangslagen} durch widersprüchliche oder überbordende politische und juristische Vorgaben bei begrenzten Ressourcen
\end{itemize}

Das fundamentale Dilemma der SLB ist, dass die Nachfrage nach ihren Leistungen systematisch die verfügbaren Ressourcen übersteigt und sie gleichzeitig mit inkonsistenten Vorgaben konfrontiert sind. Sie können die formalen Anforderungen ihres Systems prinzipiell nicht alle gleichzeitig erfüllen --- also entwickeln sie zwangsläufig \textbf{Coping-Mechanismen}.

Der deutsche Vertragsarzt passt präzise in dieses Modell: \S\,12 SGB\,V verlangt, dass ärztliche Leistungen \textit{ausreichend, zweckmäßig und wirtschaftlich} sein müssen, das \textit{Maß des Notwendigen nicht überschreiten} dürfen, und gleichzeitig leitliniengerecht zu erfolgen haben. Diese Anforderungen sind in der täglichen Praxis nicht simultan erfüllbar (siehe Drei-Code-Kollision in Teil II). Der Arzt muss diesen Widerspruch im Minuten-Takt operationalisieren, während im Hintergrund das Risiko der Regressierung besteht.

\subsubsection{Lipskys Coping-Typologie, angewandt auf den Modellarzt}

Die von Lipsky identifizierten Coping-Strategien mappen erstaunlich präzise auf die bereits dokumentierten Verhaltensanpassungen des Modellarzts (siehe voriger Unterabschnitt):

{\small
\begin{tabularx}{\textwidth}{X X X}
\toprule
\textbf{Coping-Mechanismus (Lipsky)} & \textbf{Manifestation im SGB\,V-Kontext} & \textbf{Empirischer Beleg} \\
\midrule
\textbf{Rationing} (implizite Rationierung) & Zurückhaltung teurer, aber indizierter Medikamente; Off-Label-Vermeidung trotz medizinischer Rechtfertigung & Ribbat 2023, 85\,\% unterlassene Verordnungen \\
\textbf{Referral / Outsourcing} (Abwälzen) & 51\,\% der Hausärzte überweisen an Spezialisten, obwohl sie selbst behandeln könnten & Ribbat 2023 \\
\textbf{Creaming / Cherry-Picking} (Selektion) & Vermeidung multimorbider Patienten, die das Budget überdurchschnittlich belasten & Verhaltensmuster im Modellarztprofil \\
\textbf{Routinization} (Standardisierung) & Strikte Adhärenz an Leitlinien, Vermeidung individualisierter Therapiepfade & Goetz 2024, defensive Laborroutine \\
\textbf{Limiting Engagement} (Rückzug) & Quartalsende-Verschiebung, reduzierte Patientenbindung, Niederlassungsverzicht & HCHE-Kifmann, KBV-Berufsmonitoring \\
\bottomrule
\end{tabularx}
}

Diese Mapping-Tabelle ist der analytisch zentrale Punkt dieses Teils. Die Defensivmedizin des deutschen Vertragsarzts ist damit nicht mehr eine affektive Randerscheinung oder ein bedauerlicher ``Betriebsunfall'' des Systems, der sich durch Aufklärung beheben ließe. \textbf{Sie ist der klassische, lehrbuchmäßige Bewältigungsmechanismus einer Frontlinien-Rolle in einem Ressourcen- und Widerspruchskonflikt.}

\subsubsection{Die politisch-analytische Konsequenz}

Wer Lipsky anwendet, kommt zu einer politisch unbequemen Schlussfolgerung, die über die Diagnose ``Angstsystem'' deutlich hinausgeht: Der Vertragsarzt ist nicht nur ein passives Opfer der GKV-Bürokratie. Er ist vielmehr \textit{das ausführende Organ einer impliziten, politisch nicht kodifizierten Ressourcensteuerung}. Indem er durch defensive Unterversorgung, Referral-Kaskaden und Routinization das System entlastet, übernimmt er eine Rationierungs-Funktion, die der Gesetzgeber selbst nicht offen kodifizieren würde --- und die in dieser Form auch nie offen beschlossen wurde.

Das ist die präzisere Formulierung der in Teil IV entwickelten Folgekosten-These: Der Arzt ``zahlt'' nicht nur mit seiner Zeit, seiner Versicherungsprämie und seinem Stresslevel für das System. Er zahlt mit seiner Berufsethik, indem er gezwungen wird, Entscheidungen zu treffen, deren Rationierungs-Wirkung das System produziert, aber deren moralische Last allein bei ihm liegt. Das ist die eigentliche Pointe der Lipsky-Analyse: \textit{Staatliche Ressourcenknappheit wird an der Street-Level nicht nur umgesetzt --- sie wird dort auf den Schultern einzelner Frontlinien-Akteure abgeladen.}


\section{Teil IV --- Die Folgekosten: Wenn das System mehr schadet als es spart}

\subsection{Die direkten Patientenschäden}

Die Verhaltensänderungen aus Teil III führen zu konkreten Patientenschäden. Die wichtigsten dokumentierten Wirkmechanismen:

\begin{itemize}[nosep]
  \item \textbf{Quartalsende-Effekt:} Die HCHE-Kifmann-Studie (Hamburg Center for Health Economics, TK-Daten 2013/14) dokumentiert bis zu 14\,\% weniger Termine am Quartalsende und einen Anstieg der Notdienst-Belastung um 19 Prozentpunkte.
  \item \textbf{Überweisungskaskade:} 51\,\% der Hausärzte überweisen trotz Indikation an Spezialisten (Ribbat 2023). Dies erzeugt Wartezeiten, Mehrkosten und Versorgungsverzögerungen, ohne dass das Risiko gesamtsystemisch gemindert wird.
  \item \textbf{Off-Label-Vermeidung in der Onkologie:} Bis zu 60\,\% der Krebspatienten in der Regelversorgung werden off-label behandelt, weil die Wissenschaft schneller ist als die Zulassung. Trotz BSG-Rechtsprechung zur Zulässigkeit unter bestimmten Bedingungen regressierten Krankenkassen Onkologen für Standardtherapien (DGHO-Stellungnahme 2006, 2024).
  \item \textbf{GLP-1 bei Diabetes mit Leberzirrhose:} Dokumentierter Fall, in dem ein Patient nach Aussetzung der leitliniengerechten Tirzepatid-Therapie auf Insulin umgestellt wurde --- ein Medikament, das die Grunderkrankung verschlimmert.
  \item \textbf{Vermeidbare Hospitalisierungen:} OECD Health at a Glance 2025 dokumentiert für Deutschland 810 vermeidbare Hospitalisierungen pro 100.000 Einwohner gegenüber einem OECD-Mittel von 473 (+70\,\%).
\end{itemize}

\subsection{Die volkswirtschaftliche Hochrechnung}

Die direkten und indirekten Folgekosten der regress-induzierten Unterversorgung lassen sich konservativ und verifizierbar abschätzen:

{\small
\begin{tabularx}{\textwidth}{X X X}
\toprule
\textbf{Kostenbestandteil} & \textbf{Wert (Bandbreite)} & \textbf{Quelle} \\
\midrule
Vermeidbare stationäre Aufnahmen & 7-8\,Mrd.\,EUR/Jahr & WIdO Krankenhaus-Report 2016, Zi 2021 \\
Non-Adhärenz / Therapieabbrüche & 10-19\,Mrd.\,EUR/Jahr & DIMDI HTA 65, ABDA \\
Vermeidbare Todesfälle (Versorgungsdefizite) & ca. 20.000/Jahr & Monitor Versorgungsforschung 02/2025 \\
US-Defensivmedizin (proportional auf DE) & 6-25\,Mrd.\,EUR/Jahr & Mello et al. 2010 (Health Affairs), eigene Hochrechnung \\
\bottomrule
\end{tabularx}
}

\textbf{Systemrahmen der assoziierten Folgekosten (Nenner, vor Attribution):} 7--17\,Mrd.\,EUR pro Jahr (ACSC-Hospitalisierungen + Non-Adhärenz-Kosten; Quelle: WIdO/Zi, DIMDI\,HTA\,65). Dies ist der Nenner der Zurechnungsrechnung~--- nicht die regressangst-attributierte Schätzung selbst. Das Dreistufenmodell: \emph{Systemrahmen} (7--17\,Mrd.) $\times$ \emph{Attributionsrate} (5\,\%--15\,\%) $=$ \emph{Regress-Folgekosten} (0,9--4,2\,Mrd./Jahr). Im mittleren Szenario (5\,\%, inkl. direkter Verhaltenskosten) ergibt sich \textbf{0,9--1,7\,Mrd.\,EUR/Jahr} (siehe Sensitivitätsanalyse oben).

\subsection{Die Verhältnisrechnung: Schaden zu Einnahme}

Wieviel nimmt das System ein, wieviel kostet es?

{\small
\begin{tabularx}{\textwidth}{X X X}
\toprule
\textbf{Position} & \textbf{Wert/Jahr} & \textbf{Quelle} \\
\midrule
Durchgesetzte Regresse (konservativ) & ca. 5\,Mio.\,EUR & KV-Statistiken, Hochrechnung \\
Durchgesetzte Regresse (großzügig) & niedrige zweistellige Mio.\,EUR & GKV-SV 2022 (MD-stat.) \\
Bürokratiekosten ambulant (gesamt) & 4,33\,Mrd.\,EUR & Statistisches Bundesamt \\
Regress-Folgekosten (5\,\% attributiert, eigene Hochrechnung)$^\star$ & 0,9--1,7\,Mrd.\,EUR & eigene Hochrechnung \\
\bottomrule
\end{tabularx}

\smallskip\noindent$^\star$ \textit{Szenario-Annahme, keine kausale Messung.} Der Kausalschluss ``Prüfdruck $\to$ Unterversorgung $\to$ Folgekosten'' ist nicht empirisch geschlossen. Die Zahl ist ein illustratives Modell auf Basis einer unbewiesenen Attributionsannahme (5\,\%).
}

\textbf{Szenario-Verhältnis Schaden zu Einnahme (unter den genannten Modellannahmen): ca. 100:1 bis 1.000:1.} Unter diesen Annahmen wäre selbst die großzügigste Schätzung der durchgesetzten Regresse um zwei bis drei Größenordnungen kleiner als der modellierte Schaden.

\textbf{Methodische Einschränkung:} Diese Verhältnisrechnung beruht auf Hochrechnungen mit offenen Annahmen. Die Bandbreite ist absichtlich weit, weil keine deutsche Punktschätzung der durch Regressangst kausal induzierten Folgekosten existiert. Auch die konservativste Schätzung deutet jedoch darauf hin, dass die durchgesetzten Regresse die Folgekosten nicht aufwiegen. Eine systematische Prüfung dieses Verhältnisses durch eine zuständige Stelle (BRH, SVR, IGES) hat bisher nicht stattgefunden --- siehe Teil V.

\textbf{Unvollständiger Nenner:} Die Verhältnisrechnung setzt den modellierten Schaden ins Verhältnis zu den \textit{durchgesetzten Regressen} (5--50\,Mio.\,EUR). Das unterschlägt den potenziellen Nutzen des Systems: die Abschreckungswirkung, d.\,h. Mehrausgaben, die durch die bloße Existenz des Prüfsystems vermieden werden. Diese Abschreckungswirkung ist real (Einwand~3, Abschnitt~\ref{sec:gegenposition}), aber nicht quantifiziert. Eine vollständige Kosten-Nutzen-Analyse müsste die verhinderten Mehrausgaben im Nenner berücksichtigen. Da diese Größe nicht verfügbar ist, bildet die 100:1-Zahl den \textit{worst case} für das System ab, nicht den Erwartungswert.

\textbf{Vergleichsfolie zu akzeptierten gesundheitsökonomischen Modellrechnungen:} Die hier verwendete Methodik entspricht dem Standard gesundheitsökonomischer Hochrechnungen mit transparenten Annahmen. Vergleichbare Studien:
\begin{itemize}[nosep]
  \item Mello et al. (Health Affairs 2010, ``National Costs of the Medical Liability System'') schätzen die US-Defensivmedizin-Kosten auf 45,6\,Mrd.\,USD/Jahr — mit Bandbreite und expliziten Sensitivitätsannahmen
  \item Effertz et al. (European J of Health Economics 2016) berechnen die volkswirtschaftlichen Kosten der Adipositas in Deutschland mit ähnlicher Methodik
  \item DIMDI HTA-Bericht 65 zur Non-Adhärenz arbeitet mit Bandbreiten zwischen 10 und 19\,Mrd.\,EUR/Jahr
\end{itemize}
Die hier verwendete konservative Hochrechnung ist methodisch nicht ambitionierter als diese etablierten Modelle und liegt in ihren Bandbreiten in einer vergleichbaren Größenordnung. Die Modellrechnung ist also weder spekulativ noch überambitioniert — sie wendet etablierte Methodik auf einen bisher nicht rechtsförmig untersuchten Gegenstand an.

\subsection{Vergleich mit anderen großen Kostentreibern}

Zum Vergleich mit anderen großen gesellschaftlichen Kostentreibern (Systemrahmen ACSC + Non-Adhärenz = 7--17\,Mrd.; attributierte Regress-Folgekosten konservativ = 0,9--1,7\,Mrd.; $^\star$\,=\,Szenario-Annahme):

{\small
\begin{tabularx}{\textwidth}{X X}
\toprule
\textbf{Schadensquelle} & \textbf{Volkswirtschaftliche Kosten/Jahr (DE)} \\
\midrule
Rauchen & ca. 79\,Mrd.\,EUR \\
Adipositas & ca. 63\,Mrd.\,EUR \\
Alkohol & ca. 57\,Mrd.\,EUR \\
Non-Adhärenz & 10-19\,Mrd.\,EUR \\
\textbf{Systemrahmen Regress-assoziiert (ACSC + Non-Adh., vor Attr.)} & \textbf{7--17\,Mrd.\,EUR} \\
\textbf{Davon attributierbar (5\,\%, konservativ)}$^\star$ & \textbf{0,9--1,7\,Mrd.\,EUR} \\
Krankenhauskeime & 4-12\,Mrd.\,EUR \\
Bürokratie ambulant & 2,4-4,33\,Mrd.\,EUR \\
\bottomrule
\end{tabularx}
}

\textbf{Drei Vergleichsanker} (Systemrahmen 7--17\,Mrd.; attributierte Regress-Folgekosten konservativ 0,9--1,7\,Mrd.):

\begin{itemize}[nosep]
  \item Das Schaden-zu-Einnahme-Verhältnis ist selbst im konservativen Attributionsszenario (1,1\,Mrd.) noch schlechter als bei Alkohol (Verhältnis Schaden zu Steuer: 18:1) --- und produziert keinen ``Konsumnutzen''.
  \item Der Systemrahmen der regress-assoziierten Folgekosten (7--17\,Mrd.) ist 17- bis 42-fach größer als das KBV-Bürokratie-Einsparpotenzial (400 Mio.\,EUR/Jahr). Die KV administriert das System, das ihre eigenen Mitglieder so massiv belastet. Selbst die konservative Attributionsschätzung (0,9--1,7\,Mrd.) ist noch 2- bis 4-fach größer als das Einsparpotenzial.
  \item Der Systemrahmen ist größenordnungsmäßig vergleichbar mit den Kosten aller Krankenhauskeime Deutschlands zusammen (4--12\,Mrd.) --- bei einem Bruchteil der politischen, medialen und wissenschaftlichen Aufmerksamkeit. Für Krankenhauskeime existieren das RKI-KRINKO, Hygieneverordnungen und Awareness-Kampagnen. Für Regress: Schweigen.
\end{itemize}

\subsection{Administrative Burden als ``Policymaking by Other Means''}

Die Größenordnungen aus dem vorigen Unterabschnitt --- 4,33\,Mrd.\,EUR jährliche Bürokratiekosten im ambulanten Sektor, ca. 61 Arbeitstage pro Vertragsarzt und Jahr allein für Verwaltung, rund 16.000 terminale ICD-10-Codes --- werden in der deutschen Debatte üblicherweise als Ineffizienz, schlechtes Gesetzeshandwerk oder als ``wuchernde Bürokratie'' interpretiert. Diese Lesart greift zu kurz. Die US-amerikanische Policy-Forschung liefert seit einigen Jahren einen wesentlich schärferen analytischen Rahmen: die Theorie der \textbf{Administrative Burden}.

\subsubsection{Herd \& Moynihan: Bürokratie als Politikgestaltung}

In ihrem Standardwerk \textit{Administrative Burden: Policymaking by Other Means} (Russell Sage Foundation, New York 2018 --- ausgezeichnet unter anderem mit dem Louis Brownlow Book Award 2019 und dem Herbert Simon Best Book Award 2022) definieren Pamela Herd und Donald P. Moynihan ``Administrative Burden'' \textit{explizit nicht} als zufälliges Nebenprodukt träger oder inkompetenter Verwaltung. Sie ist in ihrer Lesart \textbf{``Policymaking by Other Means''} --- Politikgestaltung mit anderen Mitteln. Wenn Regierungen oder Kostenträger Leistungen nicht offen, politisch kostspielig und gesetzlich legitimiert kürzen wollen oder können, erhöhen sie strategisch die administrativen Hürden, um die Inanspruchnahme implizit zu senken.

Herd und Moynihan zerlegen administrative Lasten in drei trennscharfe, empirisch messbare Kategorien:

{\small
\begin{tabularx}{\textwidth}{X X X}
\toprule
\textbf{Kosten-Kategorie} & \textbf{Definition} & \textbf{Übertragung auf Wirtschaftlichkeitsprüfung} \\
\midrule
\textbf{Learning Costs} (Lernkosten) & Kognitiver und zeitlicher Aufwand, um die anwendbaren Regeln überhaupt zu erfassen & Bis zu 17 regionale Prüfvereinbarungen, geheime Erstattungsbeträge, ``Container-Landschaft'' des \S\,106b SGB\,V, unklare Auslegung der Substantiierungsschwelle \\
\textbf{Compliance Costs} (Befolgungskosten) & Zeit und monetäre Ressourcen für formale, gerichtsfeste Dokumentation und Kodierung & 61 Arbeitstage/Jahr pro Arzt, 4,33\,Mrd.\,EUR Bürokratiekosten ambulant, ``ICD für den Gutachter, nicht für den Patienten'' \\
\textbf{Psychological Costs} (Psychologische Kosten) & Stress, Autonomieverlust, Frustration durch die Navigation des Systems & Ribbat 2023: 72--78\,\% ``stark emotional belastet'' durch Regresse; Goetz 2024: 27\,\% ``starke Rechtsangst''; defensive Medizin als Folge \\
\bottomrule
\end{tabularx}
}

\subsubsection{Die analytische Konsequenz: Die Bürokratie funktioniert wie vorgesehen}

Wer die Herd/Moynihan-Brille aufsetzt, liest die in Teil II dokumentierten Befunde \textbf{diametral entgegengesetzt} zur naheliegenden Lesart ``Systemversagen''. Dabei wirken V1 und V2 als komplementäre Burden-Instrumente:

\begin{itemize}[nosep]
  \item \textbf{V1 (Auffälligkeitsprüfung):} Erzeugt primär \textit{Learning Costs} und \textit{Compliance Costs}. Die 99,7\,\%-Verwerfungsquote (KV Hessen 2015) bedeutet: 913 Ärzte tragen die administrativen und psychologischen Kosten des Screenings, ohne dass ein Regress folgt. Aus Burden-Perspektive ist genau das die Funktion --- die Kosten des Verfahrens selbst, nicht sein Ergebnis, steuern das Verhalten.
  \item \textbf{V2 (Einzelfallprüfung):} Erzeugt primär \textit{Psychological Costs} durch die Kombination aus fehlender Beratungsschutzstufe, normativem Schadensbegriff und hoher Durchsetzungswahrscheinlichkeit. Hier ist nicht die Breite des Filters, sondern die Unausweichlichkeit der Sanktion das Steuerungsinstrument.
\end{itemize}

Aus Sicht der Administrative-Burden-Theorie ist die Kombination beider Verfahren \textit{kein dysfunktionales Bürokratie-Versagen}. Es wirkt wie ein hochwirksam kalibriertes Steuerungsinstrument: V1 erzeugt breite Unsicherheit, V2 erzeugt gezielte existenzielle Angst. Zusammen bringen sie Ärzte dazu, auf rechtlich und medizinisch legitime Verordnungen \textit{freiwillig} zu verzichten. Diese Wirkung ist der ``Chilling Effect''.

Die Pointe von Herd und Moynihan ist, dass dieser Effekt nicht als Nebenwirkung einer Politik einzuordnen ist, sondern als ihr \textbf{intendiertes Ergebnis unter politisch ``gesichtswahrenden'' Bedingungen}: Niemand muss offen beschließen, dass Diabetiker keine GLP-1-Agonisten bekommen sollen. Es reicht, die Hürde so hoch zu legen, dass der rationale Arzt den ``Eintrittspreis'' in Form von Dokumentationsaufwand, rechtlichem Risiko und emotionaler Belastung nicht mehr zahlt. Das Ergebnis ist eine Rationierung, die politisch nie beschlossen wurde, aber faktisch eintritt.

\textbf{Methodische Vorsicht:} Wie bei der Niskanen-Analyse in Teil V (siehe dort die Vorbemerkung ``Anreize sind nicht gleich Intentionen'') gilt auch hier: Die These der bewussten Instrumentalisierung ist theoretisch schlüssig, aber historisch schwer zu beweisen. Es lässt sich nicht zeigen, dass \textit{jemand} die deutsche Wirtschaftlichkeitsprüfung mit dem expliziten Ziel entworfen hat, Ärzte zu defensivem Verordnen zu bewegen. Was sich zeigen lässt, ist, dass das System im Ergebnis genau diese Wirkung erzielt --- und dass diese Wirkung mit den öffentlich formulierten Zielen (``Wirtschaftlichkeit'', ``Beitragssatzstabilität'') gut vereinbar ist. Die Absicht liegt nicht im Gesetzestext, sie liegt in der Akzeptanz des beobachtbaren Effekts durch die Systemarchitekten.

Das führt zur nächsten Frage: Wer profitiert eigentlich davon, dass niemand dieses System stört? Diese Frage behandelt Teil V.


\section{Teil V --- Das institutionelle Schweigen}

\subsection{Der GKV-Spitzenverband 2024: Verteidigung ohne Belege}

Im Argumentationspapier vom 26.04.2024 (im Kontext der KHVVG-Anhörung) verteidigt der GKV-Spitzenverband die Einzelfallprüfung als ``effizient'' und ``überlegen''. Die Recherche zeigt: \textbf{Es wird keine einzige empirische Effizienz-Studie vorgelegt.} Keine IQWiG-Untersuchung, keine gesundheitsökonomische Publikation, keine externe Evaluation.

Stattdessen werden drei Datensätze präsentiert, die alle das eigene Argument nicht tragen:

\begin{enumerate}
  \item \textbf{Halbierung der MD-Erstattungsbeträge 2019$\to$2022} (2,38\,Mrd. $\to$ 1,19\,Mrd.\,EUR) seit Einführung der Prüfquoten. Das beweist nur, dass weniger Prüfung zu weniger Rückforderungen führt --- keine Effizienzaussage.
  \item \textbf{47-50\,\% Beanstandungsquote} der MD-geprüften Rechnungen (Q4 2023). Heißt zugleich: jede zweite Prüfung ist unberechtigt.
  \item \textbf{85\,\% Einigungsquote im Vorverfahren} --- belegt den Druckeffekt, nicht die Prüfgüte.
\end{enumerate}

Das Kernargument ist ein Beweislast-Umkehr-Argument: ``Für die Alternative liegt weder Konzept noch Folgenabschätzung vor, also bleibt der Status quo.'' Keine Kosten-Nutzen-Analyse der Prüfverfahren selbst wird angestellt.

\textbf{Wichtige Klarstellung zum Geltungsbereich:} Die GKV-SV-Verteidigung von 2024 bezieht sich primär auf die stationäre Einzelfallprüfung (\S\,275c SGB\,V, MD-Prüfung Krankenhäuser). Die ambulante Arzneimittel-Einzelfallprüfung nach \S\,106b SGB\,V wird in derselben Argumentationslinie verteidigt, aber mit anderen Zahlen, die noch dürftiger sind.

\subsection{Wer könnte prüfen --- und tut es nicht}

Die Wirtschaftlichkeitsprüfung wird von keiner zuständigen Stelle systematisch auf ihre eigene Wirtschaftlichkeit geprüft:

{\small
\begin{tabularx}{\textwidth}{X X X}
\toprule
\textbf{Institution} & \textbf{Befugnis} & \textbf{Tatsächliches Handeln} \\
\midrule
Bundesrechnungshof (\S\,88 BHO) & ja, vollumfänglich & In den GKV-Berichten 2025: keine Sonderprüfung \\
Bundesamt für Soziale Sicherung & Aufsicht (\S\,90 SGB IV) & Keine Evaluation der Prüfverfahren \\
SVR Gesundheit & Gutachten & 2014 Richtgrößenkritik, seither Schweigen \\
Bundestags-Gesundheitsausschuss & Evaluationspflichten & Ungenutzt \\
Beschwerdeausschuss \S\,106c SGB\,V & Streitschlichtung & Keine Systemkontrolle \\
G-BA & Leistungsdefinition & Nicht zuständig \\
\bottomrule
\end{tabularx}
}

\textbf{Faktisch prüft niemand systematisch.} Es gibt einzelne Zitate aus SVR-Gutachten 2014 (Richtgrößenkritik) und vereinzelte Stellungnahmen --- aber keine durchgehende Effizienz\-aufsicht. Diese Beobachtung ist methodisch wichtig, weil sie die Beleglage für Effizienzaussagen über das Prüfsystem erklärt: Wer behauptet, das System funktioniere oder funktioniere nicht, stützt sich auf indirekte Daten. Eine systematische Effizienzaufsicht der Wirtschaftlichkeitsprüfung selbst existiert in Deutschland nicht.

\subsection{Wer profitiert eigentlich? Eine Public-Choice-Analyse}

\textbf{Methodische Vorbemerkung:} Die folgende Analyse arbeitet mit Anreizstrukturen, nicht mit Motiven. \textit{Anreize sind nicht gleich Intentionen.} Es geht nicht darum, einzelnen Akteuren böse Absichten zu unterstellen, sondern darum, zu zeigen, welche strukturellen Verstärker das System unabhängig von den Absichten seiner Beteiligten produziert. Wer in einer Bürokratie einen Anreiz zur Selbsterhaltung hat, handelt deshalb nicht unmoralisch — er handelt strukturkonform.

Wenn niemand das System auf Wirtschaftlichkeit prüft und gleichzeitig die Beleglage für seinen Nutzen leer ist — wer sorgt dann eigentlich für seinen Fortbestand? Die naheliegende Antwort wären die Krankenkassen. Die empirischen Daten zeichnen ein anderes Bild.

\subsubsection{Die Nutzen-Bilanz der Spieler}

Eine Übersicht der wesentlichen Akteure und ihres Saldos aus dem Regress-System:

{\small
\begin{tabularx}{\textwidth}{X X X}
\toprule
\textbf{Akteur} & \textbf{Was kostet das System?} & \textbf{Was bringt es?} \\
\midrule
Vertragsarzt & Angst, Defensivmedizin, Existenzrisiko & nichts Direktes \\
Patient & Unterversorgung, Folgeschäden & nichts \\
Krankenkasse & Antragskosten, Personal, Gerichtsverfahren, Umlage für GPE & 5--50\,Mio.\,EUR Regresse/Jahr (vs. 4,33\,Mrd.\,EUR ambulante Bürokratie) \\
Pharmaindustrie & Umsatzverluste durch Defensivmedizin & nichts \\
KV (als Selbstverwaltung) & Verwaltungslast & nichts Direktes \\
\textbf{Gemeinsame Prüfungseinrichtungen (GPE)} & nichts (durch Umlage finanziert) & \textbf{Existenzgrund, Stellen, Budget} \\
Beschwerdeausschüsse & nichts & Sitzungspauschalen \\
Spezialisierte Sozialrechts-Anwälte & nichts & Mandate \\
Regressschutz-Angebote (Vereine, Versicherer, Beratungen) & Helfen aktiv, politisieren, dokumentieren & Bestand an Existenz des Problems gekoppelt \textit{(Ambivalenz — siehe unten)} \\
\bottomrule
\end{tabularx}
}

Aus dieser Übersicht ergibt sich ein bemerkenswerter Befund: \textbf{Die einzigen Akteure mit positivem Saldo sind die, die das System verwalten oder von ihm leben — nicht die, die es politisch entworfen oder gerechtfertigt haben.}

\textbf{Zur Verhältnismäßigkeit der Kosten:} Die GPE-Infrastruktur (bundesweit 600--900 Stellen) wird paritätisch aus Beitragsmitteln finanziert. Ihren exakten Kosten (nicht verifizierbar, s.\,u.) stehen Regress-Einnahmen von geschätzt 5--50\,Mio.\,EUR pro Jahr gegenüber. Zum Vergleich: Die GKV-Arzneimittelausgaben lagen 2024 über 50\,Mrd.\,EUR. Selbst die obere Grenze der Regress-Einnahmen entspricht damit weniger als 0,1\,\% der kontrollierten Ausgaben --- ein Verhältnis, das die Frage nach der Wirtschaftlichkeit der Wirtschaftlichkeitsprüfung selbst aufwirft.

\textbf{Warum das System dennoch Akzeptanz genießt:} Die politische Legitimation des Prüfsystems stützt sich auf zwei Grundannahmen: (1)~den \textit{Glauben an Selbstwirksamkeit} --- die Überzeugung, dass das Prüfsystem messbar Einsparungen bewirkt (empirisch nicht belegt, siehe GKV-SV 2024) --- und (2)~den Verweis auf ``überragendes Gemeinwohl'' der Beitragssatzstabilität (\S\,71 SGB\,V). Beide Annahmen sind politisch plausibel, aber empirisch ungeprüft. Niemand hat je berechnet, ob die durch das Prüfsystem verhinderten Verordnungen einen höheren Wert haben als die durch Defensivmedizin verursachten Folgekosten.

\textbf{Gegenargument --- natürliche Begrenzungen:} Ein häufig unausgesprochenes Argument für die Notwendigkeit des Prüfsystems lautet: Ohne externe Kontrolle würden Ärzte \textit{unbegrenzt} verordnen. Dieses Argument ignoriert drei natürliche Begrenzungsmechanismen: (a)~Patienten wollen nicht unnötig medikamentös behandelt werden (Nebenwirkungsaversion, Non-Compliance bei Überversorgung); (b)~die medizinische Indikation selbst setzt Grenzen (Leitlinien, Kontraindikationen, Wechselwirkungen); (c)~die ärztliche Professionslogik (Freidson, siehe Teil VIII) enthält intrinsische Qualitätsstandards. Das Prüfsystem reguliert nicht die Medizin, sondern die Abweichung von einer statistischen Norm --- eine Funktion, die auch ohne Sanktionen durch Beratung und Information erfüllt werden könnte.

\subsubsection{Niskanens Bürokratie-Modell und die GPE}

William Niskanen hat 1971 in \textit{Bureaucracy and Representative Government} ein Modell vorgelegt, das genau solche Konstellationen beschreibt: Eine Behörde, deren Output politisch nicht direkt messbar ist, maximiert ihren eigenen Etat statt das Gemeinwohl ihres Auftrags. Niskanen identifiziert sechs Kriterien für eine selbst\-erhaltende Bürokratie. Bei den Gemeinsamen Prüfungseinrichtungen (\S\,106c SGB\,V) sind alle sechs erfüllt:

{\small
\begin{tabularx}{\textwidth}{X X}
\toprule
\textbf{Niskanen-Kriterium} & \textbf{Befund für die GPE} \\
\midrule
(a) Monofunktional & Nur Wirtschaftlichkeitsprüfung; Plausibilitäts-, Abrechnungs- und Qualitätsprüfung liegen NICHT bei der GPE \\
(b) Output-unabhängige Finanzierung & Umlage durch KV und Krankenkassen, unabhängig von Verfahrenserfolg \\
(c) Informationsasymmetrie & Keine Haushaltsveröffentlichung; der Bundesrechnungshof hat NIE eine Sonderprüfung durchgeführt \\
(d) Karrierekopplung & Teamleitung setzt Verfahrensvolumen voraus; Aufstiegspfade an Fallzahlen gebunden \\
(e) \textbf{Reformresistenz} & 2017-Abschaffung der Richtgrößenprüfung führte nur zur Umbenennung — \textit{nicht zum Personalabbau} \\
(f) Kein eigener Rechtfertigungszwang & Interessen werden indirekt über die Träger-KV vertreten \\
\bottomrule
\end{tabularx}
}

\textbf{Punkt (e) ist der entscheidende empirische Beleg.} Als der Gesetzgeber 2017 die Richtgrößenprüfung abschaffte, hätte ein Output-orientiertes System Personal abbauen müssen. Stattdessen blieb das Personal und wurde in andere Prüfaufgaben reallokiert. Die Prüfungseinrichtung Bayern beschäftigt aktuell ca. 75 Mitarbeiter (Apotheker, PTA, Juristen, BWLer, Rechtsanwaltsfachangestellte) ausschließlich für Wirtschaftlichkeitsprüfungen. Bundesweit etwa 600--900 Stellen.

\subsubsection{Kosten extern nicht verifizierbar — und das ist selbst ein Befund}

Eine eigene Detailrecherche zum Versuch, exakte Eurozahlen für die Prüfinfrastruktur (Personal, Sachkosten, Sitzungsgelder) zu ermitteln, führt zu einem bemerkenswerten Ergebnis: \textbf{Die Kosten sind aus offenen Quellen nicht verifizierbar.} Drei strukturelle Gründe:

\begin{itemize}[nosep]
  \item Die GPE sind keine staatlichen Behörden. Sie tauchen in keinem Bundes- oder Landeshaushalt mit eigenständigen Posten auf.
  \item Die Finanzierung erfolgt paritätisch durch KV und Krankenkassen-Verbände (\S\,106c Abs.\,2 SGB\,V), aber die Kostenanteile werden in den allgemeinen Verwaltungsbudgets der Träger nicht granular ausgewiesen --- sie erscheinen in den öffentlich zugänglichen Geschäftsberichten nicht als eigenständige Position ``Prüfinfrastruktur''.
  \item Der Bundesrechnungshof prüft die GPE nicht in Form tiefergehender Wirtschaftlichkeits-Sonderberichte — und zwar systematisch nicht, weil hier primär GKV-Beitragsmittel und ärztliche Honorarabzüge fließen, keine direkten Steuermittel des Bundes. Die Zuständigkeit liegt formal außerhalb des klassischen BRH-Prüfungsbereichs.
\end{itemize}

Diese Nicht-Verifizierbarkeit ist nicht ein Mangel der Recherche, sondern ein strukturelles Merkmal des Systems. Genau diese Informationsasymmetrie ist Niskanens Kriterium (c) — und sie ist der Grund, warum die institutionelle Selbsterhaltung der GPE bisher nicht extern belegt werden \textit{konnte}, obwohl sie strukturell offensichtlich ist. Die Beweislage für die Effizienz des Systems ist nicht nur leer --- sie ist strukturell nicht herstellbar, weil die nötigen Kostendaten nicht granular erfasst oder veröffentlicht werden.

\textbf{IFG-Hebel:} Eine systematische IFG-Anfrage an die Prüfungsstellen aller KV-Bezirke (oder als Sammelanfrage über den Spitzenverband) könnte diese Lücke schließen. Die Erfolgsaussichten sind unklar, weil die Gemeinsamen Prüfungseinrichtungen ihre Stellung als ``Selbstverwaltung'' tendenziell als IFG-frei interpretieren. Eine solche Anfrage wäre damit zugleich Recherche und politischer Test.

\subsubsection{Der Drehtür-Effekt}

Ein zweites Niskanen-typisches Muster: der Wechsel von Mitarbeitern der Prüfungseinrichtungen in spezialisierte Anwaltskanzleien, die auf Regressabwehr ausgerichtet sind. Dieses Muster ist strukturell beobachtbar: Ehemalige Leitungspositionen in Wirtschaftlichkeitsprüfungseinrichtungen finden sich in Fachkanzleien wieder, die genau diese Verfahren vertreten.

Dieser Personalwechsel ist rechtlich zulässig. Er erzeugt jedoch eine Anreizstruktur, in der spezialisierte Anwaltskanzleien wirtschaftlich von der Fortdauer des Prüfsystems abhängen --- auch wenn ihre Mandanten die Verfahren am liebsten gar nicht erst hätten. Die Kanzleien sind damit die zweite Gruppe mit strukturellem Interesse an der Kontinuität des Systems.

\subsubsection{Die wahren Profiteure}

Aus der empirischen Analyse ergibt sich eine überraschende Reihenfolge der strukturellen Profiteure des Systems:

\begin{enumerate}[nosep]
  \item \textbf{Spezialisierte Medizinrechts-Kanzleien} (Drehtür-Effekt belegt)
  \item \textbf{GPE-Beschäftigte} (600--900 Stellen, monofunktional)
  \item \textbf{Regressschutz-Angebote} (Vereine, Versicherer, Beratungen — siehe Anmerkung zur Ambivalenz unten)
\end{enumerate}

\textbf{Krankenkassen sind nicht auf Platz 1.} Sie kassieren wenige Millionen Euro Regress, zahlen aber die Umlage für die GPE. Ihr Nettonutzen ist nahe Null. \textit{Sie halten das System aus politisch-kultureller Legitimation am Leben, nicht aus rationaler Kostenkalkulation.}

\subsubsection*{Zur Ambivalenz der Regressschutz-Angebote}

Akteure wie Therapiefreiheit e.V., die Beratungsangebote des Virchowbundes und private Regressschutz-Versicherer nehmen in dieser Analyse eine strukturell ambivalente Rolle ein. \textbf{Einerseits} sind sie die wichtigsten Verbündeten betroffener Ärzte: Sie helfen konkret, dokumentieren Einzelfälle, politisieren das Thema, verfassen Stellungnahmen, treten als Prozessbegleiter auf und sind häufig die einzigen Stimmen, die das Regress-System öffentlich kritisieren. \textbf{Andererseits} ist ihr Bestand strukturell an die Existenz des Problems gekoppelt. Würde das Regress-System morgen abgeschafft, verlöre ihr Angebot seinen Zweck.

\textbf{Methodischer Hinweis:} Diese Beobachtung ist keine moralische Kategorie. Sie beschreibt eine Anreizstruktur, keine Intention. Anreize sind nicht gleich Motive (siehe Vorbemerkung am Anfang dieses Abschnitts). Die Akteure selbst handeln in bester Absicht und leisten für betroffene Ärzte unverzichtbare Arbeit.

\textbf{Strategische Konsequenz:} Gerade diese Akteure sind natürliche Kooperationspartner für ein Regress-Transparenzportal. Wer das Problem wirklich kennt, hat das größte Interesse daran, es messbar und sichtbar zu machen. Die Ambivalenz ist kein Grund zur Distanz, sondern eine Einladung zur gemeinsamen Veränderung des Systems.

\subsubsection{Die Forschungslücke}

Bemerkenswert ist, dass die deutsche Public-Choice-Literatur die GPE bisher nicht empirisch aufgearbeitet hat. Es existieren keine wissenschaftlichen Studien, die das Niskanen-Modell auf das deutsche Wirtschaftlichkeitsprüfungs-System anwenden. \textbf{Diese Forschungslücke ist selbst ein Datenpunkt} — und sie bestätigt die Informations\-asymmetrie-These: Wo niemand hinschaut, entstehen genau die strukturellen Eigenlogiken, die Niskanen beschreibt.


\section{Verfassungsrechtliche Angriffsvektoren}

Ein weiterer Befund der Studie lässt sich aus den vorangegangenen Teilen ableiten: Wenn das System einen rechtsdogmatischen Doppelbind enthält (Teil II), wenn es Patientenschäden in Milliardenhöhe erzeugt (Teil IV) und wenn niemand seine Effizienz prüft (Teil V), liegt die Frage nahe, ob es verfassungsrechtlich angreifbar ist. Diese Frage haben wir aus der Perspektive eines erfahrenen Verfassungsrechtlers analysiert. Das Ergebnis ist nicht eindeutig — aber es ist auch nicht negativ.

\subsection{Was nicht funktioniert: Der naheliegende Weg über Art.\,12 GG}

Der Standard-Angriff auf Regressbescheide läuft über Art.\,12 Abs.\,1 GG (Berufsfreiheit). Das BVerfG ordnet Regresse jedoch dogmatisch auf der niedrigsten Stufe der Drei-Stufen-Theorie ein und winkt die Finanzstabilität der GKV als ``überragendes Gemeinwohl'' regelmäßig durch.

Frischer Beleg: Der \textbf{Kammerbeschluss 1 BvR 985/24 vom 01.10.2025}. Eine ärztliche Verfassungsbeschwerde gegen einen Regress wurde dort ohne ausführliche Substantiierung als unzulässig abgewiesen. Karlsruhe hat das Thema also bereits auf dem Tisch — und akzeptiert die etablierte Linie.

\subsection{Die Backdoor: Art.\,20 Abs.\,3 GG, Bestimmtheitsgrundsatz und Wesentlichkeitstheorie}

Der eigentliche verfassungsrechtliche Hebel liegt eine Ebene tiefer. \S\,12 SGB\,V formuliert das Wirtschaftlichkeitsgebot tautologisch (``ausreichend, zweckmäßig, wirtschaftlich''). Für eine Leistungsnorm ist das tragbar. Für eine \textit{Sanktionsnorm mit existenzvernichtender Wirkung} — etwa die 491.000\,EUR im Stempelregress (BSG B\,6\,KA\,9/24\,R) — ist es ein Bestimmtheitsproblem.

Die Leitentscheidungen des BVerfG sind hier auf der Seite einer Kritik:

\begin{itemize}[nosep]
  \item \textbf{Kalkar I (BVerfGE 49, 89)} — Wesentliche Entscheidungen muss der Gesetzgeber selbst treffen
  \item \textbf{Numerus Clausus I (BVerfGE 33, 303)} — Wesentlichkeitstheorie auch im Vorfeld grundrechtsrelevanter Eingriffe
\end{itemize}

Beide Entscheidungen verlangen, dass der Gesetzgeber wesentliche Sanktionsbedingungen selbst regelt — und nicht an selbstverwaltungsrechtliche Prüfvereinbarungen delegiert. Die deutsche Wirtschaftlichkeitsprüfung tut genau das: Sie regelt die Substanz der Einzelfallprüfung in regionalen Prüfvereinbarungen, während das Gesetz selbst nur einen tautologischen Rahmen vorgibt.

\subsection{Die BSG-Selbstentlarvung}

Eine zusätzliche dogmatische Schwachstelle liegt in der jüngsten BSG-Rechtsprechung selbst. In den Urteilen B\,6\,KA\,5/23\,R und B\,6\,KA\,10/23\,R (beide vom 05.06.2024) hat der 6. Senat die \textbf{Differenzkostenregelung} für einen Teilbereich anerkannt: Bei Wirtschaftlichkeitsregressen wird nur der tatsächliche Mehrschaden ersetzt, nicht der volle Verordnungswert. Für \textit{Formalregresse} gilt das ausdrücklich nicht — dort bleibt der Vollregress.

In der Begründung räumt der BSG offen ein, der Vollregress diene der \textit{Verhaltenssteuerung}. Damit ist er als \textbf{Sanktion} klassifiziert, nicht als Schadensausgleich. Sanktionen unterliegen aber strengeren verfassungsrechtlichen Anforderungen — vor allem an Bestimmtheit. Das BSG hat in eigenen Worten geliefert, was eine Verfassungsbeschwerde braucht.

\subsection{Die Eigentumsdogmatik des regressierten Medikaments}

Die Eigentumsfrage des regressierten Medikaments ist verfassungsrechtlich anomal:

\begin{itemize}[nosep]
  \item Der Hersteller liefert das Medikament an den Großhandel
  \item Der Großhandel an die Apotheke
  \item Die Apotheke gibt es an den Patienten ab
  \item Der Patient verbraucht es (Untergang nach \S\,90 BGB)
  \item \textbf{Der Arzt war zu keiner Zeit Eigentümer}
  \item Eine Rückforderung vom Patienten scheitert (\S\,812 BGB: kein Leistungsverhältnis)
  \item Trotzdem haftet der Arzt mit voller Kostenhöhe
\end{itemize}

Das ist eine verfassungsrechtlich erklärungsbedürftige Konstruktion. Der Arzt zahlt für etwas, das er nie besessen hat, das nicht mehr existiert, das ihm nichts genützt hat — und das dem Patienten genützt hat, ohne dass er von ihm zurückfordern könnte. Eine Verfassungsbeschwerde mit primärem Fokus auf Art.\,14 GG (Eigentumsgarantie) hat hier ein nicht ausgeschöpftes Argumentationsfeld.

\textbf{Methodische Abgrenzung:} Die Argumentation behauptet kein Eigentum des Arztes an der ärztlichen Leistung selbst. Was sie behauptet, ist eine eigentumsrechtlich begründete Disposition: Die Regress-Forderung greift auf Vermögenswerte des Arztes zu, die durch eine Verordnung entstanden sind, an deren Eigentumssache (Medikament) der Arzt zu keinem Zeitpunkt rechtlich beteiligt war. Das ist juristisch ein zweifelhafter Eingriff in das Vermögen, nicht in das Eigentum an einer Sache. Die Linie ist daher Art.\,14 Abs.\,1 Satz 2 GG (Inhalts- und Schrankenbestimmung) — dort verlangt das BVerfG die Beachtung des Verhältnismäßigkeitsgrundsatzes (siehe BVerfGE 50, 290 — Mitbestimmungsurteil; BVerfGE 100, 226 — Denkmalschutz Rheinland-Pfalz). Das Argument ist also nicht: ``Der Arzt war Eigentümer''. Das Argument ist: ``Die Vermögensbelastung steht außer Verhältnis zum tatsächlichen Schaden, weil der Schaden \textit{strukturell null} ist.''

\subsection{Die Doppelbeschwerde-Strategie}

Aus der Persona-Analyse ergibt sich eine konkrete strategische Empfehlung: Eine \textbf{Doppelbeschwerde} mit zwei Beschwerdeführer-Profilen.

\textbf{Beschwerde 1 — der Arzt:}
\begin{itemize}[nosep]
  \item Hauptangriff: Wesentlichkeitstheorie und Bestimmtheit (Art.\,20 Abs.\,3 GG)
  \item Flanke: Art.\,14 GG (Differenzkosten als milderes Mittel)
  \item Flanke: Art.\,3 GG (Beratung-vor-Regress-Asymmetrie zwischen Auffälligkeits- und Einzelfallprüfung)
\end{itemize}

\textbf{Beschwerde 2 — der Patient:}
\begin{itemize}[nosep]
  \item Hauptangriff: Art.\,2 Abs.\,2 Satz 1 GG (staatliche Schutzpflicht für körperliche Unversehrtheit)
  \item Strategie: Fortentwicklung der Nikolaus-Rechtsprechung (1 BvR 347/98) auf schwere chronische Erkrankungen
\end{itemize}

Der strategische Wert der Doppelbeschwerde liegt darin, dass die zweite Beschwerde das BVerfG zu einer inhaltlichen Auseinandersetzung mit der Schutzpflicht zwingt. Auch eine formale Ablehnung würde rechtspolitisch wertvolle Begründung produzieren.

\subsection{Einschätzung der Erfolgsaussichten}

Die folgende Einschätzung ist eine \textit{analytische Einordnung}, keine juristische Prognose. Die tatsächlichen Erfolgsaussichten hängen von der Fallkonstellation, der Begründungsqualität und der verfassungsgerichtlichen Agenda ab:

\begin{itemize}[nosep]
  \item Vollumfängliche Stattgabe: nicht absehbar, aber angesichts der Dogmatik-Hürden gering
  \item Rechtsfortbildende Begründung (auch bei formaler Ablehnung): realistischeres Ziel, da das BVerfG sich zur Schutzpflicht positionieren müsste
\end{itemize}

Parallel zur Beschwerde gehörten EMRK-Vorbereitung (Art.\,6 und Art.\,1 ZP\,1) sowie eine politische Begleitstrategie.

\textbf{Methodische Einordnung:} Diese Studie ist keine anwaltliche Schrift. Die hier dargestellte Analyse ist eine \textit{Argumentations-Skizze}, keine fundierte Verfassungsbeschwerde-Vorbereitung. Wer das System tatsächlich verfassungsrechtlich angreifen wollte, bräuchte spezialisierten anwaltlichen Beistand und einen konkreten Fall mit dokumentierter Beschwer.


\section{Teil VI --- Internationaler Vergleich}

\subsection{Deutschland und die Schweiz als Sonderwege}

Deutschland und die Schweiz sind in Westeuropa praktisch die einzigen Länder mit flächendeckender individueller Regress-Prüfung gegen einzelne Ärzte für Verordnungen. Andere Länder steuern ärztliches Verordnungsverhalten ohne Individual-Sanktionen:

{\small
\begin{tabularx}{\textwidth}{X X}
\toprule
\textbf{Land} & \textbf{Mechanismus} \\
\midrule
UK (NHS) & QOF / CQUIN / NICE-Empfehlungen, professionelle Selbstregulation, keine individuellen Regress-Verfahren für Verordnungen \\
Niederlande & Capitation / NHG-Standards, kollektive Steuerung \\
Schweden, Norwegen, Dänemark & Zentrale Erstattungsregeln, professionelle Selbstregulation \\
Frankreich & MSO/MSAP-Verfahren (\textit{Mise Sous Objectif} / \textit{Mise Sous Accord Pr\'ealable}) mit algorithmengestützter Vorauswahl. Seit 2024 wachsender Widerstand durch MG France, FMF und Médecins Pour Demain unter der Begriffsprägung \textit{``d\'elit statistique''} --- siehe eigene Subsection unten \\
Schweiz & Art.\,56 KVG / WZW-Verfahren --- strukturell vergleichbar zu DE, aber transparenter \\
\textbf{Deutschland} & \textbf{\S\,106b SGB\,V plus 17 regionale Prüfvereinbarungen, ohne Beratungsschutz für Einzelfallprüfungen} \\
\bottomrule
\end{tabularx}
}

\subsection{Frankreich: Le ``D\'elit Statistique''}

Ein Teil der Recherche für diese Studie galt der Frage, ob Deutschland mit seinem Regress-System tatsächlich allein steht. Die kurze Antwort: Nein. Seit etwa 2023 entwickelt sich in Frankreich eine strukturell parallele Debatte, die unter dem von Ärztegewerkschaften geprägten Begriff \textbf{``d\'elit statistique''} (``statistisches Delikt'' oder ``statistisches Vergehen'') geführt wird. Sie ist das direkte Analogon zum deutschen Fall --- nur auf ein anderes Verordnungsobjekt zugeschnitten.

\subsubsection{Das Instrument: MSO und MSAP}

Die französische gesetzliche Krankenversicherung (\textit{Caisse nationale de l'Assurance Maladie}, CNAM) verfügt seit dem \textbf{D\'ecret 2011-551 vom 19. Mai 2011} (Regierung Fillon) über ein Verfahren namens \textbf{Mise Sous Objectif} (MSO) und sein schärferes Pendant \textbf{Mise Sous Accord Pr\'ealable} (MSAP). Die Mechanik ist einfach:

\begin{itemize}[nosep]
  \item Die CNAM analysiert jährlich statistisch das Verordnungsverhalten jedes Vertragsarztes, im Fokus vor allem \textit{indemnit\'es journali\`eres} (IJ --- Arbeitsunfähigkeitsbescheinigungen / Krankschreibungen).
  \item Ärzte, die laut Algorithmus ``signifikant über dem regionalen Durchschnitt'' verordnen, werden zu einem Gespräch einbestellt.
  \item Wer im Gespräch keine zufriedenstellende Erklärung liefert, wird per Verfügung unter MSO gestellt --- das bedeutet: sechs Monate lang verpflichtendes ``Reduktionsziel'' unter Androhung finanzieller Sanktionen durch Finanzstrafkommissionen (\textit{commissions des p\'enalit\'es financi\`eres}).
  \item Wer bei MSAP landet, muss jede einzelne Krankschreibung vorab von der Kasse genehmigen lassen.
\end{itemize}

\subsubsection{Größenordnung: Was die offiziellen Zahlen zeigen}

Die veröffentlichten Zahlen der CNAM und der französischen Ärztegewerkschaften sind für 2023--2024 belastbar dokumentiert (Quellen: \textit{Le Quotidien du M\'edecin}, \textit{Egora}, \textit{Previssima}, FMF-Cellule-Juridique):

\begin{itemize}[nosep]
  \item \textbf{Kampagne 2023--2024:} 745 Ärzte wurden unter MSO gestellt; 416 MSO-Verfügungen über jeweils sechs Monate (01.09.2023--29.02.2024). 310 erfüllten ihr Reduktionsziel, \textbf{106 wurden vor Finanzstrafkommissionen gebracht}.
  \item \textbf{Gesamtjahr 2024:} ca. 680.000 Kontrollen, davon 230.000 mit medizinischer Prüfung, \textbf{ca. 7.000 Sanktionen}.
  \item \textbf{2025:} Fortsetzung der Kampagne in zwei Wellen; die CNAM rechnet mit einer erhöhten Zahl von MSO-Verweigerungen und der Umstellung auf MSAP.
\end{itemize}

Diese Zahlen sind keine ``Massenproteste im Zehntausender-Bereich'' --- sie bezeichnen strukturellen, institutionell verfestigten Disziplinierungsdruck auf tausende Ärzte, mit existenzrelevanten Konsequenzen für einige Hundert pro Jahr. Sie sind damit größenordnungsmäßig durchaus vergleichbar mit der deutschen Situation.

\subsubsection{Die Begriffsprägung und der gewerkschaftliche Widerstand}

Den Begriff ``d\'elit statistique'' prägten primär zwei Akteure:

\begin{itemize}[nosep]
  \item \textbf{MG France} (der größte französische Hausarzt-Verband): ``Le d\'elit statistique est ill\'egal'' --- das statistische Delikt ist ein rechtswidriges Konstrukt. Die CNAM ``adjustiere'' ihre Statistiken und erschaffe damit eine ``pseudo-vergleichbare'' Aktivität, die die spezifische Patientenstruktur (Alter, Geschlecht, Multimorbidität, soziale Härten) einer Praxis systematisch unterkalkuliere.
  \item \textbf{F\'ed\'eration des M\'edecins de France (FMF):} Ihre juristische Zelle (\textit{cellule juridique}) hat unter dem Titel ``Histoire du d\'elit statistique'' eine rechtsdogmatische Analyse vorgelegt, die argumentiert, dass die rein statistische Begründung eines Verdachts gegen einzelne Ärzte weder mit dem Bestimmtheitsgebot noch mit dem Prinzip der individuellen Verantwortung vereinbar sei. Der Widerstand gegen MSO hat sich seit 2024 über ein breites Spektrum französischer Ärzte-Verbände ausgeweitet, einschließlich des Kollektivs \textit{Médecins Pour Demain}.
\end{itemize}

Die Kernargumentation der französischen Verbände liest sich fast wörtlich wie das, was die KBV in Deutschland sagt: Ein statistischer Ausschlag ist kein valider Beweis für unethisches oder unwirtschaftliches Verordnen. Er spiegelt in der Regel die spezifische Patientenstruktur einer lokalen Praxis wider. Die algorithmische Androhung zwinge Ärzte in ein ethisches Dilemma, in dem sie gegen ihr professionelles Gewissen kranken Patienten notwendige Krankschreibungen verweigern müssten, um die eigene Praxis vor Sanktionen zu schützen.

\subsubsection{Methodische und strategische Konsequenz}

Der Frankreich-Befund verändert die Einordnung des deutschen Sonderwegs in drei Hinsichten:

\begin{enumerate}
  \item \textbf{Der deutsche Fall ist kein isoliertes juristisches Artefakt.} Das Muster der algorithmengestützten Einzelarzt-Sanktionierung findet in mindestens zwei großen europäischen Gesundheitssystemen statt. Es ist Teil eines europäischen Trends zur \textit{mikroindividuellen Kostenkontrolle} --- nicht eine pathologische Einzelerscheinung des deutschen SGB\,V.
  \item \textbf{Das Verordnungsobjekt ist austauschbar, die Steuerungslogik nicht.} In Frankreich zielt der Algorithmus auf Krankschreibungen (IJ), in Deutschland primär auf Arzneimittel und Heilmittel. Der zugrundeliegende Mechanismus --- statistischer Vergleich + Einzelfall-Disziplinierung + finanzielle Existenzdrohung --- ist strukturell derselbe.
  \item \textbf{Der kollektive Widerstand ist in Frankreich erfolgreicher organisiert.} Die französischen Ärzte-Verbände haben ihr Solidaritäts-Gefangenendilemma (siehe Teil III) überwunden, weil sie eine politisch streikfähige Kultur haben. Die deutsche Ärzteschaft schafft diesen Übergang zur gemeinsamen öffentlichen Artikulation bisher nicht --- ein Befund, der die in Teil III beschriebene spieltheoretische Konstellation zusätzlich belegt.
\end{enumerate}

Der deutsche und der französische Fall bestätigen sich also gegenseitig: Es geht in beiden Systemen um dieselbe Frage, und in beiden Fällen wird auf dieselbe Ursache ein unterschiedlich organisierter Widerstand gebaut. Der Import des Begriffs ``d\'elit statistique'' in die deutsche Debatte ist nicht polemisch --- er ist sachlich präzise.

\subsection{Was andere Länder anders machen}

\begin{itemize}[nosep]
  \item \textbf{Positive Anreize statt Sanktionen:} QOF in UK belohnt Qualität, statt Abweichung zu bestrafen.
  \item \textbf{Capitation statt Einzelleistung:} Pauschale Vergütung pro Patient (NL, UK) reduziert den Anreiz, jede Verordnung einzeln zu kontrollieren.
  \item \textbf{Zentrale Preisbildung:} NICE in UK, Spain, Australien --- der ``richtige Preis'' wird nicht im Einzelfall ermittelt, sondern systemisch festgelegt.
  \item \textbf{Professionelle Selbstregulation:} GP-Körperschaften sanktionieren Fehlverhalten kollegial, nicht über finanzielle Drohung.
\end{itemize}

Die deutsche Praxis ist nicht alternativlos. Sie ist eine spezifische historische Wahl, die in ihrer heutigen Ausprägung international die Ausnahme ist.

\subsection{Diametrale Manifestation desselben Mechanismus}

Die internationale Defensivmedizin-Literatur belegt einen Mechanismus, der auch für Deutschland gilt --- aber sich diametral entgegengesetzt manifestiert. Diese Beobachtung ist methodisch zentral, weil sie erklärt, warum die deutschen Folgekosten schwerer messbar sind als die US-amerikanischen, ohne dass der zugrundeliegende Mechanismus ein anderer wäre.

\subsubsection{USA: Druck vom Patienten $\to$ Überbehandlung}

In den USA ist die zentrale Druckquelle das Tort Law (Kunstfehlerklage). Der Aggressor ist der Patient (und sein Anwalt). Die Reaktion der Ärzte ist \textit{Defensive Medicine} im engeren Sinne: zusätzliche Tests, zusätzliche Bildgebung, häufigere stationäre Aufnahmen, prophylaktische Eingriffe.

\textbf{Klassische empirische Belege:}
\begin{itemize}[nosep]
  \item Kessler \& McClellan (Quarterly Journal of Economics 111[2], 1996, S. 353-390): Haftungsreformen senken Gesundheitsausgaben um 5--9\,\% ohne negative Auswirkungen auf Mortalität oder Komplikationen --- ein direkter Effizienzbeleg für defensive Überbehandlung.
  \item Currie \& MacLeod (Quarterly Journal of Economics 2008): Schadensersatzdeckelungen verändern die Wahl von Behandlungswegen (z.B. C-Sections) systematisch.
  \item Mello et al. (Health Affairs 2010): Schätzen die Kosten der Defensivmedizin in den USA auf 45,6\,Mrd.\,USD/Jahr.
\end{itemize}

\subsubsection{Deutschland: Druck von der Krankenkasse $\to$ Fehlversorgung in zwei Richtungen}

In Deutschland ist die zentrale Druckquelle nicht der Patient, sondern die Krankenkasse (über das Wirtschaftlichkeitsprüfungs-System). Die Reaktion der Ärzte ist jedoch \textbf{nicht eindimensional ``weniger''}, sondern eine \textit{Fehlversorgung in zwei Richtungen}:

\begin{itemize}[nosep]
  \item \textbf{Unterversorgung bei Verordnungen:} Ärzte verschreiben weniger, vermeiden teure Medikamente, verzichten auf Off-Label-Therapien --- selbst wenn sie medizinisch indiziert wären.
  \item \textbf{Überversorgung bei Diagnostik und Überweisungen:} Ärzte veranlassen zusätzliche Bluttests, überweisen an Fachärzte, fordern bildgebende Diagnostik an --- nicht weil der Patient es braucht, sondern zur eigenen Absicherung (Goetz 2024: 54\,\% wöchentlich unnötige Labortests; Ribbat 2023: 51\,\% überweisen trotz eigener Behandlungskompetenz).
  \item \textbf{Belastung des Patienten:} Die defensive Überweisung geht zu Lasten des Patienten: zusätzliche Arzttermine, Wartezeiten, Urlaubstage, Anfahrtswege --- für Untersuchungen, deren Ergebnis im Grunde bereits feststeht. Kranke Menschen werden gezwungen, noch häufiger zum Arzt zu gehen, weil ihr Arzt sich absichern muss.
\end{itemize}

\noindent Der präzisere Begriff ist daher \textbf{Fehlversorgung}, nicht Unterversorgung: Das System lenkt ärztliche Ressourcen von der Therapie (weniger Verordnungen) in die Absicherung (mehr Diagnostik, mehr Überweisungen) um. Der Nettoverlust liegt beim Patienten.

\subsubsection{Methodische Konsequenz}

Der psychologische Wirkmechanismus (Antizipation eines Schadens, Verlustaversion nach Kahneman/Tversky, defensives Risikoverhalten) ist in beiden Systemen \textit{identisch}. Was sich unterscheidet, ist die institutionelle Anreizstruktur, die das defensive Verhalten in eine bestimmte Richtung lenkt:

{\small
\begin{tabularx}{\textwidth}{X X X}
\toprule
 & \textbf{USA (Tort Law)} & \textbf{Deutschland (SGB\,V)} \\
\midrule
Druckquelle & Patient (Klage) & Krankenkasse (Regress) \\
Sanktionsrichtung & ``zu wenig getan'' & ``zu viel/falsch verordnet'' \\
Defensive Reaktion & \textit{mehr} Behandlungen & Verordnungen $\downarrow$, Diagnostik/Überweisungen $\uparrow$ \\
Netto-Effekt & Überversorgung & \textbf{Fehlversorgung} (weniger Therapie, mehr Absicherung) \\
Empirische Sichtbarkeit & gut messbar (Kostensteigerung) & schlecht messbar (Versorgungslenkung unsichtbar) \\
Wer trägt die Last? & System (Kosten) & \textbf{Patient} (Termine, Wartezeiten, Lebenszeit) \\
Folgekosten & Überversorgung & Unterversorgung + Hospitalisierung \\
\bottomrule
\end{tabularx}
}

Die diametrale Manifestation erklärt, warum die internationale Literatur (vor allem die US-amerikanische) so umfangreich zur Defensivmedizin existiert, während die deutsche Variante kaum quantifiziert ist: \textbf{Unterversorgung ist methodisch viel schwerer zu messen als Überversorgung.} Wenn ein US-Arzt einen unnötigen Test anordnet, taucht das in der Abrechnung auf. Wenn ein deutscher Arzt eine notwendige Verordnung unterlässt, taucht es nirgends auf --- außer später in der Hospitalisierungsstatistik, und auch dort ohne Ursachenzuschreibung.

\textbf{Diese Asymmetrie ist die zentrale methodische Begründung für das Konzept eines Regress-Transparenzportals.} Solange Unterlassungen nicht sichtbar sind, kann ihre Verbreitung nicht bewiesen werden --- und solange sie nicht bewiesen ist, gibt es keinen politischen Reformdruck. Das Portal würde diese Sichtbarkeitslücke erstmals systematisch schließen.

Die Tort-Reform-Studien aus den USA können damit als \textit{methodische Validierung} des deutschen Falls gelesen werden: Sie zeigen, dass regulatorischer Druck auf Ärzte systematisch ihr Verhalten verändert. Was sie nicht zeigen, weil sie es nicht zeigen können, ist die spezifisch deutsche Ausprägung dieses Effekts.


\section{Grenzen der Evidenz und Gegenposition}
\label{sec:gegenposition}

Diese Studie argumentiert, dass das Wirtschaftlichkeitsprüfsystem systemische Fehlanreize setzt --- insbesondere die Einzelfallprüfung (V2) mit ihrem fehlenden Beratungsschutz und dem normativen Schadensbegriff, aber auch die Auffälligkeitsprüfung (V1) über den Zertifikatsfehler. Eine ehrliche Analyse muss die stärksten Gegenargumente benennen und ihnen standhalten. Die folgenden sechs Einwände sind die gewichtigsten, die gegen die zentrale These vorgebracht werden können. Wo möglich, wird angegeben, welches Verfahren (V1, V2 oder beide) der jeweilige Einwand betrifft.

\subsection{Einwand 1: Korrelation ist keine Kausalität (betrifft V1+V2)}

\textbf{Der Einwand:} Die Ribbat-Studie (2023) misst \textit{subjektive Wahrnehmung und selbstberichtetes Verhalten}, keine objektiven Versorgungsergebnisse. Dass 47\,\% der Hausärzte angeben, die Regressgefahr beeinflusse ihr Verschreibungsverhalten, beweist nicht, dass Patienten tatsächlich unterversorgt werden. Steigende GKV-Arzneimittelausgaben (über 50\,Mrd.\,EUR 2024) sprechen gegen eine systematische Unterversorgung.

\textbf{Befund:} Der Einwand trifft einen realen Punkt. Diese Studie verfügt nicht über eine kausalanalytische Studie, die Versorgungsergebnisse direkt mit Prüfintensität korreliert --- etwa einen Vergleich von KV-Regionen mit unterschiedlicher Einzelfallprüfungsfrequenz. Ein solches natürliches Experiment wäre methodisch der stärkste Beleg, existiert aber nicht.

Was die Studie zeigen kann: (a) Die verhaltensökonomische Literatur belegt, dass regulatorischer Druck ärztliches Verhalten systematisch verändert (Kessler \& McClellan 1996, Currie \& MacLeod 2008). (b) Steigende Gesamtausgaben schließen sektorale Unterversorgung nicht aus --- Biosimilar-Wachstum und Hochpreistherapien treiben das Volumen, während bei Standardverordnungen Zurückhaltung herrschen kann. (c) Die Datenlücke selbst --- dass niemand misst, ob Unterlassungen stattfinden --- ist ein Befund über das System, nicht über die Behauptung.

\textbf{Zur Operationalisierung:} Steigende GKV-Arzneimittelausgaben sind als Gegenargument gegen Unterversorgung methodisch ungeeignet. Die Ausgabensteigerung wird primär durch drei Faktoren getrieben, die mit dem Verordnungsverhalten niedergelassener Ärzte wenig zu tun haben: (1)~Hochpreisarzneimittel (Onkologika, Biologika, Gentherapeutika), deren Stückpreise in die Hunderttausende gehen; (2)~demografischer Wandel (mehr Multimorbidität, mehr Polypharmazie); (3)~Patentauslauf-Zyklen, die kurzfristig zu Generika-Einsparungen und langfristig zu neuen Hochpreistherapien führen. Ein Arzt, der aus Regressangst auf ein 30-EUR-Standardmedikament verzichtet, erscheint in der GKV-Gesamtstatistik nicht --- während ein einziges neues Krebsmedikament für 100.000\,EUR den Aggregatwert dominiert. Die Gesamtausgaben messen den Preistrend des Arzneimittelmarkts, nicht das Verordnungsverhalten der Praxis.

\subsection{Einwand 2: Der Zertifikatsfehler ist ein Strohmann (betrifft V1+V2)}

\textbf{Der Einwand:} \S\,106b prüft Wirtschaftlichkeit, nicht Qualität. Das System \textit{verspricht} kein medizinisches Zertifikat. Die Studie erfindet ein Sollkriterium und kritisiert dessen Nicht-Erfüllung.

\textbf{Befund:} Der Einwand ist dogmatisch korrekt: Das Wirtschaftlichkeitsprüfsystem ist als Ausgabensteuerungsinstrument konzipiert, nicht als Qualitätssicherung. Der Zertifikatsfehler beschreibt nicht einen Fehler \textit{im} System, sondern die Lücke zwischen Systemmessung und öffentlicher Bedeutungszuschreibung: Die niedrige Regressquote wird in der politischen Kommunikation als Beleg für ein ``funktionierendes System'' herangezogen --- aber was ``funktioniert'', ist die Ausgabenkontrolle, nicht die Versorgungsqualität.

Der Kern des Arguments bleibt: Wer die öffentlich kommunizierte Gesamtquote ($<$1\,\%, V1+V2 zusammen, dominiert durch V1) als Entwarnung kommuniziert, setzt implizit ein Qualitätsversprechen voraus, das das System nicht einlöst. Diese Interpretationslücke ist real, auch wenn das System sie nicht selbst produziert.

\textbf{Präzisierung:} Die Analyse in Teil~II zeigt, dass das Prüfsystem --- insbesondere V2 --- faktisch Formeinhaltung kontrolliert, nicht Ausgaben und nicht Versorgungsqualität. Der normative Schadensbegriff (Abschnitt~\ref{sec:normativer-schaden}) definiert den ``Schaden'' als formale Regelabweichung, unabhängig von medizinischem Nutzen oder tatsächlicher Kostenwirkung. Das Prüfsystem ist damit weder ein Kostenkontrollinstrument (es fragt nicht, ob die Ausgabe sinnvoll war) noch ein Qualitätsinstrument (es fragt nicht, ob die Behandlung angemessen war) --- es ist ein Formkontrollinstrument mit Sanktionswirkung. Diese funktionale Einordnung verschärft den Strohmann-Einwand: Der Zertifikatsfehler ist nicht nur eine Interpretationslücke, sondern eine strukturelle Fehlzuordnung des gesamten Verfahrenszwecks.

\subsection{Einwand 3: Das Präventions-Paradoxon --- der Nutzen wird nie gerechnet (betrifft primär V1)}

\textbf{Der Einwand:} Die Studie rechnet nur Kosten (Unterversorgung, Defensivmedizin, Folgeschäden), nie den Nutzen. Ohne Prüfsystem könnten die Arzneimittelausgaben unkontrolliert steigen. Die Abschreckungswirkung ist real und möglicherweise volkswirtschaftlich bedeutsam.

\textbf{Befund:} Der Einwand ist berechtigt. Diese Analyse fokussiert auf die Schadensseite und erhebt nicht den Anspruch einer vollständigen Kosten-Nutzen-Analyse. Die Abschreckungswirkung des Prüfsystems ist real und schwer messbar --- das macht sie nicht weniger real. Beide Seiten --- Nutzen wie Schaden --- argumentieren mit nicht-messbaren Größen. Dieses symmetrische Problem einzuräumen ist methodisch redlich. Ein natürliches Experiment bot sich 2017 mit der Abschaffung der Richtgrößenprüfung: Wenn der Verordnungsanstieg nach 2017 messbar war, spräche das für eine substanzielle Abschreckungswirkung. Eine Auswertung der verfügbaren Daten (WIdO, GKV-Arzneimittelindex) zeigt: Die GKV-Arzneimittelausgaben stiegen nach 2017, jedoch primär durch patent\-geschützte Hochpreisarzneimittel, nicht durch einen Mengenanstieg bei Standardverordnungen. Die Verordnungsmenge (DDD) war 2018 sogar rückläufig. Ein klarer ``Dammbruch'' nach Wegfall der Richtgrößen ist nicht nachweisbar --- was darauf hindeutet, dass die Abschreckung durch Richtgrößen ohnehin begrenzt war. Allerdings substituierten regionale Prüfverfahren (Zielwerte, Durchschnittswerte) die Richtgrößen teilweise, sodass der kausale Schluss begrenzt ist.

\subsection{Einwand 4: ``Deterministisch'' ist eine Setzung, kein Befund (betrifft V2)}

\textbf{Der Einwand:} Die Eskalationspipeline für V2 (Einzelfallprüfung) ist fast komplett ``k.D.'' (keine Daten). ``Deterministisch'' als Eigenschaft zuzuschreiben, ohne über Transitionswahrscheinlichkeiten zu verfügen, ist ein \textit{argumentum ad ignorantiam} --- aus der Abwesenheit von Daten wird ein Befund konstruiert.

\textbf{Befund:} Der Einwand ist formal gültig. ``k.D.'' heißt: keine publizierten Daten verfügbar. Die Bezeichnung ``deterministisch'' beschreibt die \textit{strukturelle Anlage} des Verfahrens (kein Beratungsschutz, normativer Schadensbegriff, kein Ermessensspielraum bei Formfehlern), nicht eine empirisch gemessene Transitionsrate. Wir formulieren deshalb: ``strukturell so angelegt, dass nach Kassenantrag eine hohe Übergangswahrscheinlichkeit anzunehmen ist --- empirisch jedoch nicht verifizierbar, weil die nötigen Verfahrensdaten nicht publiziert werden.'' Die Ehrlichkeit über die Datenlücke schwächt das Argument nicht --- sie ist selbst ein Befund über die Transparenz des Systems.

\subsection{Einwand 5: Ein Drittel des Systems ist rationale Steuerung (betrifft V2)}

\textbf{Der Einwand:} Geschätzt 35\,\% der V2-Prüfungen codieren medizinische Evidenz als Formregeln (AM-RL, G-BA-Beschlüsse, Biosimilar-Pflicht). Das System \textit{ist} eine Evidenzsteuerung --- die Studie übertreibt die Dysfunktionalität und zeichnet ein zu einseitiges Bild.

\textbf{Befund:} Der Einwand ist berechtigt und hat geschätzt 35--40\,\% Erklärungskraft (getestet als alternative Hypothese H2, siehe Anhang~B). Drei Einschränkungen:

\begin{enumerate}[nosep]
  \item \textbf{Auch bei evidenzbasierter Steuerung fehlt der Lernkanal.} Der Arzt erfährt \textit{dass} er zahlen muss, nicht \textit{warum} die Evidenzlage so entschieden hat. Der punitive Kanal (V2) läuft parallel zum edukativen Kanal (CME/Fortbildung, \S\,95d SGB\,V), ohne mit ihm verbunden zu sein. Ein Arzt, der eine AM-RL-Anlage nicht kennt, wird bestraft, statt informiert.
  \item \textbf{V2b (Kategorie~B, ${\sim}$40\,\%) bleibt vollständig dysfunktional.} Stempel-Regress 491.000\,EUR bei medizinisch korrekter Behandlung ist durch keine Evidenzsteuerungslogik rechtfertigbar.
  \item \textbf{Kein Beratungsschutz --- auch nicht bei Kategorie~A.} \S\,106b Abs.\,2 S.\,4 schließt ``Beratung vor Regress'' bei V2 gesetzlich aus, unabhängig davon ob die Prüfung evidenzbasiert oder rein formal ist.
\end{enumerate}

\noindent Die Forderung dieser Studie präzisiert sich durch den Einwand: Nicht Abschaffung der Evidenzsteuerung (Kategorie~A), sondern Ergänzung des punitiven Kanals um einen edukativen --- Beratungsschutz und Vorabinformation \textit{auch} bei V2 (vgl. Abschnitt~\ref{sec:evidenz-formale-verkleidung}).

\subsection{Einwand 6: Positive Entwicklungen ausgeblendet (betrifft V1+V2)}

\textbf{Der Einwand:} Das System reformiert sich: Abschaffung der Richtgrößenprüfung (2017), Differenzkostenregelung (BSG 2024), regionale Bagatellgrenzen. Die Studie zeichnet ein statisches Bild eines dynamischen Systems.

\textbf{Befund:} Drei wesentliche Reformen sind anzuerkennen:

\begin{itemize}[nosep]
  \item \textbf{Abschaffung Richtgrößenprüfung (2017):} Positiv, aber betrifft ausschließlich V1 (Auffälligkeitsprüfung). Die Einzelfallprüfung (V2) --- deren strukturelle Asymmetrie den Kern dieser Studie bildet --- blieb unangetastet.
  \item \textbf{Differenzkostenregelung (BSG B\,6\,KA\,5/23\,R, 2024):} Positiv, aber vom BSG explizit auf ``wirtschaftlich unzweckmäßige'' Verordnungen beschränkt. Bei ``unzulässigen'' Verordnungen (= Formfehler) wurde die Anwendung ausdrücklich ausgeschlossen --- also genau dort, wo Lesart~C den operativen Schwerpunkt verortet.
  \item \textbf{Regionale Bagatellgrenzen:} Regional unterschiedlich, nicht bundeseinheitlich. Schließen Hochbetragsregresse nicht aus.
  \item \textbf{Fristenverkürzung (TSVG 2019 / GVWG 2021):} Positiv. Die Ausschlussfrist wurde von 4 Jahren (BSG-Richterrecht) auf 2 Jahre (V1) bzw. 2,5 Jahre (V2) verkürzt (siehe Abschnitt~\ref{sec:ausschlussfristen}). Das begrenzt die zeitliche Rückreichweite neuer Prüfungen --- beseitigt aber nicht die strukturellen Probleme innerhalb laufender Verfahren.
\end{itemize}

Das System reformiert sich punktuell. Die strukturelle Asymmetrie zwischen V1 und V2 --- insbesondere der fehlende Beratungsschutz (\S\,106b Abs.\,2 S.\,4) und die Nicht-Anwendbarkeit der Differenzkostenregelung bei Formfehlern --- bleibt unangetastet. Die Reformen adressieren die Symptome, nicht den in dieser Studie beschriebenen Mechanismus.

\subsection*{Zusammenfassung: Was diese Studie behauptet --- und was nicht}

Diese Studie behauptet \textbf{nicht}, dass das Wirtschaftlichkeitsprüfsystem abgeschafft werden sollte, dass alle Regresse ungerechtfertigt sind, oder dass Ärzte keiner Kontrolle bedürfen. Sie behauptet:

\begin{enumerate}[nosep]
  \item Das System misst überwiegend nicht, was es zu messen vorgibt (Zertifikatsfehler) --- wobei ${\sim}$35\,\% der V2-Prüfungen durchaus evidenzbasierte Steuerung leisten (Abschnitt~\ref{sec:evidenz-formale-verkleidung}). Das Problem ist nicht die Evidenz, sondern der Kommunikationskanal: punitiv statt edukativ.
  \item Die Einzelfallprüfung ist strukturell asymmetrisch gegenüber der Auffälligkeitsprüfung ausgestaltet (kein Beratungsschutz, normativer Schadensbegriff).
  \item Die Datenlage über die Wirkung des Systems ist so dünn, dass weder Befürworter noch Kritiker ihre Position empirisch vollständig belegen können.
  \item Die Transparenzlücke ist nicht zufällig, sondern strukturell --- und sie zu schließen wäre der wirksamste einzelne Schritt zur Versachlichung der Debatte.
\end{enumerate}


\section{Teil VII --- Lösungsraum}

\subsection{Verweis auf PP-003: Das Regress-Transparenzportal}

Die zentrale Lösung, die diese Studie unterstützt, ist nicht in ST-001 selbst beschrieben, sondern im Konzeptpapier \textbf{PP-003 -- Regress-Transparenzportal}. Es schlägt ein anonymes, verifiziertes Meldeportal für ärztliche Regressforderungen vor, das die Datenlücke schließt, die in dieser Studie als zentrale institutionelle Lücke identifiziert wurde.

PP-003 ist ein Lösungsangebot, kein Anspruch. Die Studie ST-001 begründet, warum die Lücke existiert und welche Folgen sie hat. PP-003 zeigt einen praktikablen Weg, sie zu schließen.

\subsection{Das Transparenz-Paradoxon --- und warum PP-003 es berücksichtigt}

Eine Transparenzforderung ist in der technokratisch aufgeklärten Debatte bequem --- sie klingt selbstverständlich richtig. Die Organisationsforschung der letzten zwei Jahrzehnte hat aber gezeigt, dass Transparenz keine neutrale Eigenschaft eines Systems ist. Sie wirkt verhaltensändernd, und zwar nicht immer in die gewünschte Richtung. Wer ein Transparenzportal vorschlägt, muss diese Gegenwirkung ernst nehmen --- sonst baut er einen weiteren Schritt im bürokratischen Rüstungswettlauf, den er eigentlich beenden will. Die entscheidende Referenz zu diesem Punkt ist Ethan Bernstein.

\subsubsection{Bernsteins Befund}

In einer heute vielzitierten Feldstudie in der zweitgrößten Mobiltelefon-Fabrik der Welt (China) hat Ethan S. Bernstein von der Harvard Business School dokumentiert, was er das ``Transparency Paradox'' nennt. Veröffentlicht als \textit{``The Transparency Paradox: A Role for Privacy in Organizational Learning and Operational Control''} in \textit{Administrative Science Quarterly} 57 (2), Juni 2012, S. 181--216 (DOI: 10.1177/0001839212453028). Die Kernbefunde:

\begin{itemize}[nosep]
  \item \textbf{Entgegen der Management-Erwartung} führt mehr Sichtbarkeit nicht zu mehr Produktivität. Sie führt, unter bestimmten Bedingungen, zu \textit{weniger}.
  \item Beobachtete Akteure entwickeln \textit{Verbergungstechniken} und bauen ``Compliance Theater'' auf, das der Beobachtung genügt, aber mit dem tatsächlichen Arbeitsvorgang nur noch lose verbunden ist.
  \item Sie verfallen in einen \textit{Spotlight-Effekt}: Risiko- und Innovationsbereitschaft sinken, weil jede Abweichung vom Muster als potenzieller Fehler registriert werden könnte.
  \item Bernsteins empirisch belegter Gegenvorschlag: \textbf{``Zones of Privacy''} --- bewusst entkoppelte, beobachtungsfreie Räume, in denen Akteure experimentieren, lernen und professionelle Urteilskraft einsetzen können, ohne unmittelbar sanktionsexponiert zu sein.
\end{itemize}

\subsubsection{Die ernsthafte Kritik an PP-003}

Angewandt auf das deutsche Regress-Transparenzportal ist das eine unbequeme, aber notwendige Kritik. Wer ein anonymes Meldeportal baut, das aggregierte Regressdaten pro Kasse, Wirtschaftsraum und Wirkstoffgruppe sichtbar macht, \textit{verändert das Verhalten der Beteiligten} --- nicht nur der Ärzte, die das Portal befüllen, sondern vor allem der Kostenträger, die später darin auftauchen. Zwei pessimistische Szenarien sind anhand der Bernstein-Literatur zu erwarten:

\begin{enumerate}
  \item \textbf{Compliance Theater auf Kassenseite:} Kassen, die sich als ``Top-Ablehner'' oder ``Ärzte-Schreck'' in einem öffentlichen Dashboard wiederfinden, haben nicht unbedingt den Anreiz, milder zu prüfen. Sie haben zunächst den Anreiz, ihr Prüfverhalten \textit{informeller} und \textit{weniger sichtbar} zu machen --- etwa über Beratungsgespräche, vertrauliche Vergleiche, modifizierte Auswahl-Algorithmen, die die statistischen Spitzen glätten, ohne die Praxis zu ändern.
  \item \textbf{Spotlight-Effekt auf Arztseite:} Die öffentliche Sichtbarkeit dokumentierter Regressfälle könnte, anders als intendiert, die subjektive Bedrohungswahrnehmung verstärken, statt sie zu kalibrieren. Einzelne Fallbilder --- etwa der 491.000-EUR-Stempelregress (BSG B\,6\,KA\,9/24\,R) --- sind psychologisch hochwirksam und lassen sich im Aggregations-Dashboard nicht wegrechnen. Die Verfügbarkeitsheuristik würde die reale Wahrscheinlichkeit überformen und die Defensivmedizin eher zusätzlich befördern.
\end{enumerate}

Beide Szenarien sind nicht zwangsläufig --- aber sie sind plausibel genug, um bei der Gestaltung des Portals mitgedacht zu werden.

\subsubsection{Warum PP-003 robuster ist, als es auf den ersten Blick wirkt}

Eine detaillierte Prüfung von PP-003 zeigt, dass mehrere Gestaltungsentscheidungen das Bernsteinsche Paradoxon bereits bewusst umgehen --- wenngleich im Konzeptpapier nicht unter diesem Namen:

\begin{itemize}[nosep]
  \item \textbf{Wahlkabinen-Prinzip:} Die strikte kryptografische Trennung zwischen Identitätsprüfung (EFN-Verifikation) und Falldatenbank stellt sicher, dass die meldende Person nicht auf ihren Meldungen zurückverfolgt werden kann. Das ist \textit{präzise} das, was Bernstein eine ``Zone of Privacy'' für den Frontlinien-Akteur nennt: Der Arzt kann seine Erfahrung in den Datenpool einspeisen, ohne dass sie auf ihn zurückfällt. Die Bereitschaft, überhaupt zu melden, steht und fällt mit diesem Schutz.
  \item \textbf{k-Anonymität (k$\geq$5):} Keine Zelle im öffentlichen Dashboard wird veröffentlicht, in der weniger als fünf Einzelmeldungen zusammenlaufen. Das verhindert die Rück-Identifikation einzelner Ärzte über Kombinationen seltener Merkmale --- und begrenzt zugleich die psychologische Wirkung von Einzelfall-Spotlights.
  \item \textbf{Fokus auf Aggregate, nicht auf Einzelfälle:} Das Dashboard publiziert Verhältniszahlen und Trends, keine anekdotischen Erzählungen. Das zielt bewusst auf die Kalibrierung des Wahrscheinlichkeits-Empfindens und gegen den Spotlight-Effekt einzelner Schock-Geschichten. Einzelfall-Dokumentation bleibt Journalismus und Zivilgesellschaft überlassen, nicht dem Portal selbst.
\end{itemize}

\subsubsection{Was PP-003 allein nicht lösen kann}

Ehrlich gesagt: Ein Portal allein wird das Bernsteinsche Paradoxon auf der Kassenseite nicht beheben. Wenn Krankenkassen sich aus der öffentlichen Sichtbarkeit zurückziehen und ihre Sanktionspraxis in informellere Kanäle verlagern, verschiebt das Portal das Problem nur, ohne es zu lösen. Der strukturelle Anreiz zum Compliance Theater auf Kostenträgerseite bleibt --- es sei denn, das Portal wird durch \textbf{rechtlich bindende Berichtspflichten} und \textbf{unabhängige Aufsicht} flankiert. Genau diese Flankierung leisten die in der nächsten Subsection vorgeschlagenen minimal-invasiven Reformhebel (Halbsatz-Streich \S\,106b Abs.\,2 S.\,4, Publikationspflicht der Prüfungsstellen, Sichtbarmachung \S\,48 BMV-Ä).

\textbf{Strategische Konsequenz:} PP-003 ist ein notwendiger, aber nicht hinreichender Schritt. Seine Wirkung wird voll erst im Verbund mit strukturellen Reformen sichtbar. Das Portal erzeugt den \textit{informationellen Druck}, die Reformhebel den \textit{institutionellen}. Eines ohne das andere verpufft.

\subsection{Die juristisch belegten Top-3-Schutzmaßnahmen}

Eine eigene Detailrecherche zur Praxisrealität zeigt: Es gibt drei juristisch wirksame Schutzmaßnahmen, die in Sozialgerichtsverfahren regelmäßig den Unterschied machen:

\begin{enumerate}
  \item \textbf{Lückenlose Dokumentation zum Zeitpunkt der Verordnung.} SG Marburg 14.02.2024 (S 18 KA 96/23): Wer die Behandlungsdokumentation vorlegen kann, gewinnt. BSG B\,6\,KA\,26/13: nachgereichte Dokumentation reicht nicht; sie muss zeitnah erstellt sein.
  \item \textbf{Vorab-Klärung in den expliziten Ausnahmefällen} (Cannabis nach \S\,31 Abs.\,6 SGB\,V, häusliche Krankenpflege \S\,37, Soziotherapie \S\,37a, stationäre Reha \S\,40, bestimmte Heil-/Hilfsmittel). \textbf{Wichtig:} Bei normalen Arzneimitteln mit Indikation ist Vorab-Genehmigung gesetzlich verboten (\S\,29 BMV-Ä). Hier hilft nur defensive Dokumentation oder im Grenzfall eine BSG-konforme Bitte um Stellungnahme (B\,6\,KA\,27/12\,R) — die juristisch nicht bindet, aber Beweislast-Vorteile schafft.
  \item \textbf{Substantiierter Vortrag von Praxisbesonderheiten BEREITS im Verwaltungsverfahren} — nicht erst vor dem Sozialgericht (LSG BW 15.11.2023, SG Marburg 31.01.2024).
\end{enumerate}

Alle drei Maßnahmen sind softwaregestützt umsetzbar. Genau das ist der Anlass für die folgende Skizze.

\subsection{Open-Source-Verordnungstool: Konkrete MVP-Spezifikation}

Eine eigene Detailrecherche (siehe \texttt{RECHERCHE\_PVS\_TOOL\_MACHBARKEIT\_DETAIL.md}) zeigt: Ein fokussiertes Browser-Tool ohne Patientendaten ist in 6--9 Monaten von 1--2 Entwickler:innen machbar (Prototype-Fund-Volumen ca. 47.500\,EUR). Es ist ausdrücklich \textit{kein Praxisverwaltungssystem} (PVS), sondern ein Companion-Tool, das daneben arbeitet.

\textbf{Eingabe:} ICD-10-GM-Code + Wirkstoff/ATC.\\
\textbf{Verarbeitung:} Plausibilitätsprüfung gegen öffentliche Regelwerke (AM-RL Anlage III/V/VI, PRISCUS 2.0, GKV-SV-Praxisbesonderheiten).\\
\textbf{Ausgabe:} Ampel (grün / gelb / rot) mit Begründungstext und Quellenverweis.

\textbf{MUST-Have-Funktionen für den MVP:}

\begin{enumerate}[nosep]
  \item Echtzeit-Plausibilitätsprüfung gegen AM-RL Anlage III/V/VI
  \item Strukturierte Begründungspflicht-Templates mit Direkt-Eintrag in die Patientenakte (Audit-Trail)
  \item Container-sensitiver Vorab-Klärungs-Workflow (Cannabis vs. normale Arzneimittel — verschiedene Rechtswege)
  \item Bundesweite Praxisbesonderheiten-Erkennung mit Quartalsmarkierung
  \item Manipulationssicherer Compliance-Log (Beweiskraft im Sozialgerichtsverfahren)
\end{enumerate}

\textbf{Trägermodell-Empfehlung:} Eine gemeinnützige Trägerorganisation in Kooperation mit einem medizinischen Review-Partner (z.\,B. Therapiefreiheit e.V., DDG, BVDN). Förderprofil: Prototype Fund der Open Knowledge Foundation Deutschland oder vergleichbare Civic-Tech-Förderlinien. Keine KV als Bauträger (struktureller Interessenkonflikt).

\textbf{Begründung gegen ein vollständiges Open-Source-PVS:} Ein vollständiges KBV-zertifiziertes PVS würde 60--120 Personenjahre Initialentwicklung plus 1--2\,Mio.\,EUR/Jahr Betrieb erfordern (Hauptknackpunkte: Datenlizenzen ABDATA/MMI, KBV-Zertifizierung, gematik-Audits). Für einen Think Tank nicht machbar. Das Companion-Tool umgeht diese Hürden, indem es ausdrücklich kein PVS ist und keine Patientendaten verarbeitet.

\subsection{Politische Reform: Drei minimal-invasive Hebel}

Aus den vorangegangenen Befunden ergeben sich drei Reformschritte, die jeweils einzeilig im Gesetzestext umsetzbar sind und politisch realistisch erscheinen, weil sie den bestehenden Container nicht sprengen, sondern entgräten. Sie sind nach Aufwand-Wirkung-Verhältnis priorisiert:

\subsubsection*{Hebel 1: Halbsatz-Streich \S\,106b Abs.\,2 S.\,4 SGB\,V}

Der vierte Satz von \S\,106b Abs.\,2 SGB\,V schließt die Einzelfallprüfung explizit aus dem Schutz ``Beratung vor Regress'' aus. Eine Streichung dieses Halbsatzes wäre die kleinste mögliche Änderung mit der größten Wirkung: Damit gälte auch für Einzelfallprüfungen, dass beim erstmaligen Anlass eine Beratung statt eines Regresses erfolgt.

\textbf{Begründung:} Die Asymmetrie zwischen Auffälligkeits- und Einzelfallprüfung beim Beratungsschutz (siehe Teil II) ist die einzige rechtsdogmatisch erklärungsbedürftige Schutzlücke. Sie wurde 2017 im GKV-VSG eingefügt und nicht begründet.\\
\textbf{Aufwand:} Eine Zeile.\\
\textbf{Wirkung:} Senkt den Druck auf alle Vertragsärzte ohne die Prüfungsmöglichkeit aufzuheben.\\
\textbf{Politische Akzeptanz:} Hoch, weil bereits in der KBV-Kritik 2015 vorgeschlagen.

\subsubsection*{Hebel 2: Publikationspflicht der Prüfungsstellen (Vollzugsstatistik)}

\S\,106c SGB\,V regelt heute die Aufgaben der Prüfungsstellen, enthält aber keine Veröffentlichungspflicht für die Vollzugsstatistik. Eine Ergänzung um einen neuen Absatz ``Die Prüfungsstelle veröffentlicht jährlich eine Statistik über Verfahrenseinleitungen, -einstellungen und -gründe. Die Daten sind in maschinenlesbarer Form bereitzustellen'' wäre die zweite minimale Reform.

\textbf{Begründung:} Ohne diese Daten ist jede Effizienzdiskussion unmöglich. Der Wegfall der Informationsasymmetrie ist nach der spieltheoretischen Analyse (Teil III) die wirkungsmächtigste Einzelintervention. Sie würde auch das in PP-003 vorgeschlagene Transparenzportal teilweise überflüssig machen — was gewollt ist.\\
\textbf{Aufwand:} Ein Absatz.\\
\textbf{Wirkung:} Ermöglicht erstmals empirische Aufsicht und kalibriert die Risikowahrnehmung der Ärzte.\\
\textbf{Politische Akzeptanz:} Mittel, weil die GPE strukturell verlieren — aber durch die KV RLP bereits präzedenzlos erlaubt.

\subsubsection*{Hebel 3: Sichtbarmachung von \S\,48 BMV-Ä im SGB\,V}

\S\,48 BMV-Ä (``Sonstiger Schaden'') ist die rechtliche Grundlage des Stempelregresses (BSG B\,6\,KA\,9/24\,R, 491.000\,EUR), lebt aber außerhalb des SGB\,V im Bundesmantelvertrag. Wer ``nur'' das Gesetz liest, übersieht ihn. Eine Aufnahme des Tatbestands als neuer \S\,106e SGB\,V (oder eine ausdrückliche Verlinkung im SGB\,V) wäre der dritte minimale Schritt.

\textbf{Begründung:} Die Sichtbarmachung ist eine Bedingung dafür, dass die rechtsdogmatische Diskussion über Bestimmtheit und Verhältnismäßigkeit (siehe Verfassungsrechts-Teil) überhaupt im allgemeinen Rechtsbewusstsein geführt werden kann.\\
\textbf{Aufwand:} Ein neuer Paragraf oder eine Verlinkung.\\
\textbf{Wirkung:} Ändert die Norm nicht inhaltlich, hebt aber die Sichtbarkeit.\\
\textbf{Politische Akzeptanz:} Hoch, weil rein redaktionell — kein Akteur verliert dadurch.

\subsubsection*{Was diese drei Hebel nicht lösen}

Die Reform-Skizze adressiert nicht den \textit{strukturellen Doppelbind} aus Teil II (\S\,12 SGB\,V als Drei-Code-Kollision). Dessen Auflösung wäre eine größere Reform — sie müsste das Wirtschaftlichkeitsgebot von der Sanktionsmechanik trennen. Die natürliche Anschlussfrage führt zu einem breiteren Lösungsraum.

\subsection{Sieben weitere Interventionspfade}

Die folgende Liste wurde durch eine unabhängige Gegenlesung der Studie ermittelt (frischer Blick ohne Kenntnis der Recherche\-geschichte). Die Vorschläge leiten sich ausschließlich aus den in ST-001 dokumentierten Befunden ab. Reihenfolge nach geschätzter Wirksamkeit, nicht nach Realisierbarkeit.

\begin{enumerate}

  \item \textbf{Öffnung der operativen Prüfkataloge für Ärzte.}
  Die Studie zeigt (Teil~II), dass Kassen und Apotheken Echtzeit-Prüfsoftware besitzen, während der verordnende Arzt die konkreten Prüfkriterien nicht kennt. Das verletzt das rechtsstaatliche Bestimmtheitsgebot: Wer für Formfehler existenziell haftet, muss die Maßstäbe kennen. Eine Veröffentlichungspflicht für die operativen Prüfkataloge in den regionalen Prüfvereinbarungen würde die Informationsasymmetrie direkt beseitigen.
  \textit{Aufwand: mittel (Verwaltungsvorschrift). Realisierbarkeit: mittel --- rechtsstaatlich schwer abzulehnen.}

  \item \textbf{Sanktionsfolge für unbegründete Kassenanträge.}
  Die Studie belegt (Teil~II, Machtungleichgewicht): Für die Kasse ist ein unbegründeter Antrag folgenlos --- keine Kostentragung, keine Mengenbegrenzung. Das erzeugt das Abmahnwesen-Muster. Eine Kostenfolge bei vollständiger Verwerfung des Antrags (z.\,B. Pflicht zur Übernahme der Stellungnahmekosten des Arztes) würde die Antragsqualität strukturell erhöhen. Das Prinzip existiert in der ZPO (\S\,91) und im Abmahnrecht bereits.
  \textit{Aufwand: mittel (neuer Absatz in \S\,106b). Realisierbarkeit: niedrig--mittel.}

  \item \textbf{Positive Rückkopplungssignale im Prüfsystem.}
  Die Regelkreisanalyse (Teil~III) identifiziert: Das System sendet ausschließlich Minus-Signale. Kein Arzt erfährt je ``Ihre Verordnungspraxis ist unauffällig''. Eine institutionalisierte positive Rückmeldung --- etwa ein jährliches Unauffälligkeits-Zertifikat der KV bei regelkonformer Praxis --- würde das Panoptikum-Gleichgewicht durchbrechen und eine kalibrierbare Risiko\-ein\-schätzung ermöglichen. Internationales Vorbild: QOF im NHS.
  \textit{Aufwand: gering (KV-interne Maßnahme, kein Gesetz nötig). Realisierbarkeit: hoch.}

  \item \textbf{Kostenfolgen-Pflichtabschätzung bei Regressanträgen.}
  Die Public-Choice-Analyse (Teil~V) zeigt: Die Netto-Einnahmen aus Regressen (5--50\,Mio.\,EUR/Jahr) stehen in keinem Verhältnis zu den Systemkosten (4,33\,Mrd.\,EUR ambulante Bürokratie). Eine Pflichtabschätzung des erwarteten Regressbetrags vs. der Verfahrenskosten vor jedem V2-Antrag würde irrationale Kleinstanträge unterbinden.
  \textit{Aufwand: mittel (Verwaltungsvorschrift). Realisierbarkeit: mittel --- das Argument ``Wirtschaftlichkeit der Wirtschaftlichkeitsprüfung'' ist politisch schwer abzulehnen.}

  \item \textbf{Bundeseinheitliche Bagatellgrenze für die Einzelfallprüfung.}
  Die Studie dokumentiert (Top-10 Inkonsistenzen, Nr.\,4): Regionale Bagatellgrenzen variieren zwischen 30 und 150\,EUR --- eine verfassungsrechtlich bedenkliche Ungleichbehandlung. Eine bundesgesetzliche Mindestgrenze (z.\,B. 200\,EUR) in \S\,106b SGB\,V würde Formfehler-Regresse bei Niedrigbeträgen eliminieren, die maximal Angst und minimal Ertrag produzieren.
  \textit{Aufwand: gering (ein Absatz). Realisierbarkeit: hoch --- bereits in KBV-Forderungsliste.}

  \item \textbf{Maximale Gesamtverfahrensdauer gesetzlich begrenzen.}
  Die Studie dokumentiert (Abschnitt~\ref{sec:ausschlussfristen}): Ausschlussfristen begrenzen nur den Erstbescheid. Das BSG entschied Fälle 8--11 Jahre nach dem Verordnungszeitraum. Eine gesetzliche Gesamtverfahrensobergrenze (z.\,B. 5 Jahre ab Verordnung bis rechtskräftiger Abschluss) würde die psychologische Dauerbelastung begrenzen, die Ribbat 2023 als Haupttreiber der Defensivmedizin identifiziert.
  \textit{Aufwand: mittel (neuer Absatz in \S\,106). Realisierbarkeit: mittel --- justizpolitisch anschlussfähig an die TSVG-Fristenverkürzung 2019.}

  \item \textbf{Unabhängige Effizienzprüfung des Prüfsystems selbst.}
  Die Studie belegt (Teil~V, Niskanen-Analyse): Kein Akteur (BRH, BAS, SVR, Bundestag) prüft die Wirtschaftlichkeit der Wirtschaftlichkeitsprüfung. Die GPE-Kosten sind extern nicht verifizierbar. Ein gesetzlicher Evaluationsauftrag an den SVR Gesundheit oder das IQWiG --- alle 5 Jahre eine Kosten-Nutzen-Analyse des Prüfsystems --- würde die Niskanen-Dynamik der institutionellen Selbsterhaltung durchbrechen.
  \textit{Aufwand: gering (Evaluationsklausel in \S\,106c). Realisierbarkeit: hoch --- das Prinzip ``Evaluation öffentlicher Ausgaben'' ist politischer Konsens.}

\end{enumerate}

\subsubsection{Konzeptioneller Vorschlag: Governance-Pharmakovigilanz}

Die Pharmakovigilanz überwacht systematisch die Nebenwirkungen von Arzneimitteln nach der Zulassung. Für regulatorische Eingriffe in die Versorgungspraxis existiert kein Äquivalent. Weder der SVR Gesundheit noch das IQWiG noch der Bundesrechnungshof evaluieren systematisch, ob die Wirtschaftlichkeitsprüfung \textit{mehr Schaden anrichtet als sie verhindert}. Die vorliegende Studie argumentiert, dass ein solches Monitoring --- eine \textbf{Governance-Pharmakovigilanz} --- institutionalisiert werden sollte: eine regelmäßige, unabhängige Bewertung der Versorgungs-Nebenwirkungen regulatorischer Maßnahmen, mit denselben methodischen Standards, die für Medikamente selbstverständlich sind (Meldepflichten, Signaldetektion, regelmäßige Benefit-Risk-Assessments).

\medskip
\noindent \textbf{Verhältnis zu PP-003 und den drei Minimalhebeln:} PP-003 und die Hebel 1--3 (Abschnitt oben) sind kurzfristig umsetzbar und brauchen keine politische Mehrheit. Die sieben hier genannten Pfade setzen teilweise institutionelle Initiative voraus, ergänzen aber den Lösungsraum in Dimensionen, die PP-003 allein nicht adressiert: Machtungleichgewicht (Nr.\,2, 4), Informationsasymmetrie beim Normadressaten (Nr.\,1), Verfahrensdauer (Nr.\,6), Systemselbstevaluation (Nr.\,7) und positive Anreizsetzung (Nr.\,3).

\subsection{Kosten-Nutzen-Vergleich aller Lösungsvorschläge}

Die folgende Tabelle bewertet alle in dieser Studie entwickelten Lösungsvorschläge nach Aufwand, politischer Realisierbarkeit, erwartetem Nutzen und resultierendem Kosten-Nutzen-Verhältnis (KNV). Die Einschätzungen stützen sich auf die in den vorhergehenden Abschnitten dargestellten Begründungen.

{\small
\begin{longtable}{p{0.3cm} p{4.5cm} p{1.2cm} p{1.4cm} p{1.5cm} p{1.2cm} p{1.2cm}}
\toprule
\textbf{\#} & \textbf{Maßnahme} & \textbf{Aufw.} & \textbf{Realis.} & \textbf{Nutzen} & \textbf{KNV} & \textbf{Zeit} \\
\midrule
1 & Halbsatz-Streich \S\,106b Abs.\,2 S.\,4 (Beratungsschutz für V2) & gering & hoch & hoch & \textbf{A+} & kurz \\
\addlinespace
2 & Sichtbarmachung \S\,48 BMV-Ä im SGB\,V & gering & hoch & mittel & \textbf{A} & kurz \\
\addlinespace
3 & Bundeseinheitliche Bagatellgrenze & gering & hoch & mittel & \textbf{A} & kurz \\
\addlinespace
4 & Positive Rückkopplung (Unauffälligkeits-Zertifikat) & gering & hoch & mittel & \textbf{A} & kurz \\
\addlinespace
5 & Unabhängige Effizienzprüfung des Prüfsystems & gering & hoch & hoch & \textbf{A} & mittel \\
\addlinespace
6 & Publikationspflicht der Prüfungsstellen (Hebel~2) & gering & mittel & transform. & \textbf{A$-$} & kurz \\
\addlinespace
7 & \textbf{PP-003: Regress-Transparenzportal} & mittel & mittel & hoch & \textbf{B+} & mittel \\
\addlinespace
8 & Öffnung der operativen Prüfkataloge & mittel & mittel & hoch & B+ & mittel \\
\addlinespace
9 & Kostenfolgen-Pflichtabschätzung bei Anträgen & mittel & mittel & mittel & B & mittel \\
\addlinespace
10 & Maximale Gesamtverfahrensdauer & mittel & mittel & mittel & B & mittel \\
\addlinespace
11 & Open-Source-Verordnungstool & mittel & mittel & mittel & B & mittel \\
\addlinespace
12 & Sanktionsfolge für unbegründete Kassenanträge & mittel & niedrig & hoch & B & lang \\
\addlinespace
13 & IFG-Anfragen an Prüfungsstellen & gering & niedrig & gering--mittel & C & kurz \\
\addlinespace
14 & Governance-Pharmakovigilanz & hoch & niedrig & transform. & B$-$ & lang \\
\bottomrule
\end{longtable}
}

\noindent\textit{Aufw. = legislativer, technischer, finanzieller Aufwand. Realis. = politische Machbarkeit. KNV = Kosten-Nutzen-Verhältnis (A+ bis C). Zeit: kurz $<$1\,J., mittel 1--3\,J., lang $>$3\,J.}

\medskip
\noindent\textbf{Einordnung PP-003:} Das Regress-Transparenzportal positioniert sich im oberen Mittelfeld (Rang~7). Es ist --- wie diese Studie betont --- \textit{notwendig, aber nicht hinreichend}. Seine strategische Funktion ist die eines zivilgesellschaftlichen Druckaufbaus: Wird die Publikationspflicht (Rang~6) politisch realisiert, wird das Portal teilweise überflüssig --- was ausdrücklich das gewünschte Ergebnis wäre. PP-003 ist der Plan~B, der Plan~A erzwingen soll.

\medskip
\noindent\textbf{Strategisches Fazit:} Die höchste Hebelwirkung liegt in der Kombination der Ränge 1--6: geringe Aufwände, hohe Realisierbarkeit, kurzfristig umsetzbar. Der Halbsatz-Streich (Rang~1) ist der effizienteste Einzelhebel --- \textit{eine Zeile Gesetzestext} würde die strukturelle Asymmetrie zwischen V1 und V2 beseitigen und den wichtigsten einzelnen Treiber der Regressangst eliminieren.


\section{Quellen und Recherche-Dateien}

Diese Studie integriert Befunde aus zahlreichen parallelen Recherchesträngen. Die folgende Übersicht listet die thematischen Blöcke, zu denen eigene Arbeitsdokumente im projektinternen Infomaterial-Ordner existieren:

\begin{itemize}[nosep]
  \item Grundzahlen und Fallbeispiele (Regressquoten, KV-Statistiken, dokumentierte Einzelfälle)
  \item Pipeline-Daten und Prüfungsarten (Auffälligkeitsprüfung, Einzelfallprüfung, Abrechnungsprüfung)
  \item Formfehler-Rechtsprechung und Eigentumsdogmatik (BSG B\,6\,KA\,9/24\,R u.\,a.)
  \item Bürokratiekosten und Kodierungsbelastung in Arztpraxen
  \item Akademische Einordnung (Chilling Effect, CIRS, Pharmakovigilanz)
  \item Folgekosten und Patientenschäden (Ribbat 2023, HCHE, OECD, Verhältnisrechnung)
  \item Internationaler Vergleich und wissenschaftliche Primärliteratur
  \item Gesetzgebungshistorie der Prüfverfahren (1989--2026)
  \item Mechanik und Selbstverständnis der Einzelfallprüfung
  \item Changelog der Einzelfallprüfung (1989--2026)
  \item Argumentation des GKV-Spitzenverbandes zur Einzelfallprüfung
  \item Container-Genealogie der Wirtschaftlichkeitsprüfung
  \item Liga-Tabelle der volkswirtschaftlichen Kostentreiber
  \item Software-Markt und OpenPrescribing.net (internationaler Vergleich)
  \item PVS-Kostenanalyse und Open-Source-Realisierbarkeit
  \item Machbarkeitsanalyse des Verordnungs-Plugin-Tools
  \item Eigenanalyse zum Systemwiderspruch der Prüfverfahren (Reliabilität vs. Validität)
  \item Anonymisierte Fallstudie zu GLP-1-Verordnung bei Leberzirrhose
  \item GPE-Selbsterhaltung (Niskanen-Analyse)
  \item Strukturelle Inkonsistenzen im Normgefüge (Top-10)
  \item Container-Virus (rechtsdogmatische Zielkonflikte)
  \item Verfassungsrechtlicher Angriffsvektor (Backdoor-Analyse)
  \item Theoretische Einordnung: Managerialismus und Audit Society (Freidson 1970/2001, Power 1997)
  \item Ärzte als Street-Level Bureaucrats (Lipsky 1980)
  \item Administrative Burden als Policymaking by Other Means (Herd \& Moynihan, Russell Sage 2018)
  \item Transparency Paradox (Bernstein, ASQ 2012) als Selbstkritik des Transparenzportal-Konzepts
  \item Frankreich: D\'elit statistique, MSO/MSAP-Verfahren der CNAM (MG France, FMF, Médecins Pour Demain; 2023--2025)
\end{itemize}

\textbf{Die in v0.4 neu eingearbeiteten internationalen und theoretischen Primärquellen (alle extern verifiziert):}

\begin{itemize}[nosep]
  \item Freidson, Eliot (1970): \textit{Profession of Medicine}. New York: Dodd, Mead.
  \item Freidson, Eliot (1970): \textit{Professional Dominance}. New York: Atherton Press.
  \item Freidson, Eliot (2001): \textit{Professionalism: The Third Logic. On the Practice of Knowledge}. Cambridge: Polity Press.
  \item Power, Michael (1997): \textit{The Audit Society: Rituals of Verification}. Oxford: Oxford University Press. ISBN 0-19-828947-2.
  \item Lipsky, Michael (1980, erw. 2010): \textit{Street-Level Bureaucracy: Dilemmas of the Individual in Public Services}. New York: Russell Sage Foundation. 30th Anniversary Expanded Edition 2010, ISBN 9780871545442.
  \item Herd, Pamela \& Moynihan, Donald P. (2018): \textit{Administrative Burden: Policymaking by Other Means}. New York: Russell Sage Foundation. ISBN 9780871544445. Louis Brownlow Book Award 2019, Herbert Simon Best Book Award 2022.
  \item Bernstein, Ethan S. (2012): ``The Transparency Paradox: A Role for Privacy in Organizational Learning and Operational Control''. \textit{Administrative Science Quarterly} 57(2): 181--216. DOI: 10.1177/0001839212453028.
  \item D\'ecret n\textsuperscript{o} 2011-551 du 19 mai 2011 (Frankreich): Rechtsgrundlage für MSO-Verfahren der CNAM gegen Vertragsärzte.
  \item MG France: Stellungnahmen zum ``d\'elit statistique'' (2024/2025), z.\,B. in \textit{Le Quotidien du M\'edecin}, \textit{Egora}.
  \item F\'ed\'eration des M\'edecins de France (FMF), Cellule Juridique: ``Histoire du d\'elit statistique'' und ``La responsabilit\'e juridique de la CNAM face \`a la pression exerc\'ee sur les m\'edecins pour d\'elit statistique'' (fmfpro.org, 2024--2025).
  \item CNAM-Kontrollzahlen 2024: ca. 680.000 Kontrollen, 230.000 mit medizinischer Prüfung, 7.000 Sanktionen; MSO-Kampagne 2023--2024 mit 745 betroffenen Ärzten und 106 Finanzstrafkommissions-Verfahren (Quellen: Le Quotidien du M\'edecin, Previssima, Egora).
\end{itemize}

\subsection*{Formales Literaturverzeichnis}

Das folgende Verzeichnis listet alle 50 formalisierten Quellen dieser Studie --- peer-reviewed Artikel, Monographien, Rechtsprechung, institutionelle Quellen und Gesetzestexte. Inline-Zitationen im Text werden schrittweise auf \texttt{\textbackslash cite\{\}} umgestellt.

\nocite{*}
\bibliographystyle{unsrt}
\bibliography{ST-001_references}

\subsection*{Interne Arbeitsdokumente}


\section*{Anhang: Glossar}
\addcontentsline{toc}{section}{Anhang: Glossar}

Diese Studie verwendet einige fachliche Begriffe, die hier kurz erklärt werden:

\begin{sloppypar}
\begin{itemize}[leftmargin=1em, itemsep=0.6em]
  \item \textbf{Administrative Burden} --- Begriff der US-amerikanischen Policy-Forschung (Herd \& Moynihan, Russell Sage 2018). Bezeichnet administrative Hürden nicht als zufälliges Bürokratie-Nebenprodukt, sondern als \textit{``Policymaking by Other Means''} --- strategisch eingesetzte Lasten zur impliziten Zugangsbegrenzung. Drei Kostenkategorien: Learning Costs, Compliance Costs, Psychological Costs.
  \item \textbf{Audit Society} --- Theorie von Michael Power (Oxford UP, 1997). Beschreibt die Ausweitung von Prüfungs- und Rechenschaftspflichten auf alle gesellschaftlichen Bereiche. Prüfungen wirken performativ (sie verändern das Verhalten der Geprüften) und dienen primär der Legitimationsbeschaffung, nicht der Qualitätssicherung (``Rituale der Verifikation'').
  \item \textbf{Auffälligkeitsprüfung} --- Statistisches Prüfverfahren nach \S\,106 SGB\,V, das von der Prüfungsstelle von Amts wegen durchgeführt wird, wenn ein Arzt im Vergleich zur Fachgruppe statistisch auffällig wirtschaftet.
  \item \textbf{Container (im Sinne dieser Studie)} --- Ein einzelner Paragraph oder eine zusammenhängende Norm, in dem rechtliche Konzepte ``aufbewahrt'' werden. Der Begriff ist eine analytische Hilfe, um die Wanderung von Konzepten zwischen verschiedenen Normstandorten nachzuverfolgen, nicht ein juristischer Fachbegriff.
  \item \textbf{D\'elit statistique} --- Aus dem Französischen: ``statistisches Delikt''. Begriffsprägung der französischen Ärztegewerkschaften MG France und FMF seit 2023, gerichtet gegen die rein statistisch begründete Sanktionspraxis der CNAM (MSO/MSAP-Verfahren). Bezeichnet den Vorwurf, dass ein rein statistischer Ausschlag gegenüber dem regionalen Durchschnitt kein valider Beweis für unethisches oder unwirtschaftliches Verordnen sein kann, weil er die spezifische Patientenstruktur einer Praxis systematisch unterkalkuliert.
  \item \textbf{Differenzkostenregelung} --- Berechnungsmethode, bei der nur der tatsächliche Mehrschaden ersetzt wird, nicht der volle Verordnungswert. Vom BSG 2024 (B\,6\,KA\,5/23\,R, B\,6\,KA\,10/23\,R) für einen Teilbereich anerkannt; bei Formalregressen ausgeschlossen.
  \item \textbf{Drei-Code-Kollision} --- Eigenprägung dieser Studie. Beschreibt das systemtheoretische Phänomen, dass auf einer Operationseinheit (ärztliche Verordnung) drei verschiedene Funktionssystem-Codes (medizinisch / ökonomisch / juristisch) zusammentreffen, ohne dass die Differenzierung an dieser Operationseinheit aufgehoben wäre. Bezugspunkt: Niklas Luhmann, Theorie funktional differenzierter Gesellschaften.
  \item \textbf{Einzelfallprüfung} --- Prüfverfahren nach \S\,106b SGB\,V, das von einer Krankenkasse auf Antrag eingeleitet wird und sich auf eine konkret bezeichnete Verordnung bezieht. Der Antrag muss substantiierten Anfangsverdacht enthalten. Die Einzelfallprüfung ist explizit vom Schutz ``Beratung vor Regress'' (\S\,106b Abs.\,2 SGB\,V) ausgeschlossen.
  \item \textbf{GPE --- Gemeinsame Prüfungseinrichtung} --- Körperschaft des öffentlichen Rechts nach \S\,106c SGB\,V, paritätisch von KV und Krankenkassen-Spitzenverbänden getragen. Zuständig ausschließlich für die Wirtschaftlichkeitsprüfung. Bundesweit ca. 600--900 Stellen.
  \item \textbf{Managerialismus} --- Steuerungslogik, die professionelle Autonomie durch externe, quantifizierbare Leistungsindikatoren und formale Rechenschaftspflichten ersetzt. Nach Eliot Freidson eine der drei konkurrierenden ``Logiken'' der Arbeitsorganisation neben dem freien Markt und dem Professionalismus (``Third Logic'').
  \item \textbf{MSO / MSAP} --- Französisch: \textit{Mise Sous Objectif} (Setzung unter Reduktionsziel) und \textit{Mise Sous Accord Pr\'ealable} (Vorab-Genehmigungspflicht). Sanktionsverfahren der französischen CNAM gegen Ärzte, deren Verordnungen statistisch über dem regionalen Durchschnitt liegen. Eingeführt durch D\'ecret 2011-551 (Regierung Fillon). Das zentrale Analogon zum deutschen \S\,106b SGB\,V.
  \item \textbf{Niskanen-Modell} --- Ökonomisches Modell der Bürokratie nach William Niskanen (1971), das beschreibt, wie Behörden mit nicht direkt messbarem Output ihre eigenen Etats statt das Gemeinwohl ihres Auftrags maximieren. Sechs Identifikationskriterien.
  \item \textbf{Nikolaus-Beschluss} --- Entscheidung des BVerfG vom 06.12.2005 (1 BvR 347/98). Bei lebensbedrohlichen oder regelmäßig tödlichen Erkrankungen und Fehlen einer alternativen Behandlung besteht ein Anspruch auf nicht regulär erstattete Behandlung. Kodifiziert 2012 in \S\,2 Abs.\,1a SGB\,V. Gilt nicht für die Normalversorgung.
  \item \textbf{Professional Dominance} --- Konzept von Eliot Freidson (\textit{Profession of Medicine} 1970, \textit{Professional Dominance} 1970, \textit{Professionalism: The Third Logic} 2001). Beschreibt die gesellschaftlich legitimierte Autonomie, Selbstregulation und Definitionsmacht einer echten Profession. Wird laut Freidson seit den 1980er-Jahren durch Managerialismus systematisch erodiert.
  \item \textbf{Public Choice} --- Theorierichtung der ökonomischen Analyse politischer und administrativer Prozesse. Untersucht, wie Eigeninteressen von Akteuren in Bürokratien und politischen Institutionen kollektive Ergebnisse prägen.
  \item \textbf{Street-Level Bureaucracy (SLB)} --- Theorie von Michael Lipsky (\textit{Street-Level Bureaucracy}, Russell Sage 1980). Beschreibt Akteure an der Frontlinie staatlicher Leistungserbringung (Lehrer, Polizisten, Sozialarbeiter, Ärzte im GKV-System), die in direktem Kontakt mit Bürgern stehen, über erheblichen Ermessensspielraum verfügen und unter unauflösbaren Zwangslagen (Ressourcenknappheit + widersprüchliche Vorgaben) agieren. Sie entwickeln zwangsläufig Coping-Mechanismen: Rationing, Referral/Outsourcing, Creaming, Routinization, Limiting Engagement.
  \item \textbf{Substantiierungsschwelle} --- Beweisanforderung an einen Antrag auf Einzelfallprüfung. Die Krankenkasse muss laut \S\,106b Abs.\,1 SGB\,V und den Rahmenvorgaben (KBV/GKV-SV, Stand 01.05.2020) einen ``ausreichenden Verdacht auf Unwirtschaftlichkeit'' substantiiert vortragen.
  \item \textbf{Transparency Paradox} --- Organisationstheoretischer Befund von Ethan S. Bernstein (\textit{Administrative Science Quarterly} 57[2], 2012, S. 181--216). Mehr Sichtbarkeit und Beobachtung kann kontraproduktiv wirken: Akteure entwickeln Verbergungstechniken, bauen ``Compliance Theater'' auf und verfallen in einen Spotlight-Effekt (Risiko- und Innovationsvermeidung). Bernsteins Gegenvorschlag: ``Zones of Privacy'' als bewusst entkoppelte, beobachtungsfreie Räume zur Ermöglichung professionellen Urteils.
  \item \textbf{Verdachtsprüfung vs. Screening} --- Normative Unterscheidung. Verdachtsprüfung verlangt einen substantiierten Anfangsverdacht und führt im Trefferfall zu einer Sanktion. Screening ist ein systematisches Suchverfahren ohne Anfangsverdacht und führt im Trefferfall zu Empfehlung oder Beratung. Die Einzelfallprüfung ist nach geltendem Recht eine Verdachtsprüfung.
  \item \textbf{Wesentlichkeitstheorie} --- Verfassungsrechtliche Doktrin des BVerfG (Kalkar I, BVerfGE 49, 89; Numerus Clausus I, BVerfGE 33, 303). Wesentliche Entscheidungen --- insbesondere grundrechtsrelevante --- müssen vom Gesetzgeber selbst getroffen werden und dürfen nicht an untergesetzliche Ebenen delegiert werden.
  \item \textbf{Wirtschaftlichkeitsprüfung} --- Sammelbegriff für die Prüfung der Wirtschaftlichkeit ärztlicher Verordnungen nach \S\,106 ff. SGB\,V. Umfasst heute Auffälligkeitsprüfung, Einzelfallprüfung sowie weitere Spezialverfahren.
\end{itemize}
\end{sloppypar}


\vspace{1cm}
\begin{center}
Dieses Dokument steht unter CC BY 4.0.\\[1em]
Kontakt: hallo@um-bruch.org\\
Web: um-bruch.org
\end{center}

\newpage
\section*{Anhang B: Methodik --- Walk of Analysis}
\addcontentsline{toc}{section}{Anhang B: Methodik (Walk of Analysis)}
\label{sec:appendix-methodik}

Diese Studie entstand in einem mehrstufigen, KI-gestützten Analyseprozess zwischen dem 08. und 11. April 2026. 31 Recherche-Dateien, 3 Session-Berichte und 108 manuelle PDF-Kommentare dokumentieren den Verlauf. Die vollständige erweiterte Methodendokumentation mit allen Teilanalysen, rekonstruierten Prompt-Kernen, Modellrechnungen und Cross-Model-Divergenzen ist verfügbar unter \texttt{WALK\_OF\_ANALYSIS\_EXTENDED.md}.

\textbf{Methode der Methodenrekonstruktion:} Die Dokumentation des eigenen Analyseprozesses wurde selbst mehrstufig erstellt: (1)~Haiku-Schwarm-Extraktion der Prompt-Kerne aus 31 Recherche-Dateien, (2)~Opus-Synthese aus Session-Kontexten, (3)~menschliche Erinnerung an Richtungsentscheidungen (LG), (4)~Konsolidierung aus 3 Session-Berichten. Diese Meta-Methode --- die Rekonstruktion der Analyse durch die analysierten Werkzeuge selbst --- ist Teil der methodischen Transparenz.

\subsection*{Multi-Modell-Architektur}

{\small
\begin{tabularx}{\textwidth}{l l X}
\toprule
\textbf{Modell} & \textbf{Rolle} & \textbf{Kontrollfunktion} \\
\midrule
Claude (Opus 4.6) & Primäranalyse, Synthese (29/31 Dateien) & Kohärenzprüfung, Konsistenz-Review \\
\addlinespace
Copilot (OpenAI) & Konzeptschärfung, Core-Etablierung & Fehlkonzept-Erkennung, Begriffsklärung \\
\addlinespace
Gemini (Deep Research) & Theorierahmen, Quellenergänzung & Unabhängige Literatursuche, Visualisierung \\
\addlinespace
LG (Mensch) & Fragestellungen, Richtungsentscheidungen & Plausibilitätsprüfung, Fallkenntnis, Freigabe \\
\bottomrule
\end{tabularx}
}

\subsection*{Analyseverlauf in 9 Phasen}

{\small
\begin{tabularx}{\textwidth}{l l X X}
\toprule
\textbf{Phase} & \textbf{Dateien} & \textbf{Leitfrage} & \textbf{Schlüsselbefund} \\
\midrule
I. Empirie & 4 & Was wissen wir gesichert über das Regresssystem? & Regressquote $<$1\,\%, gefühlte Bedrohung $>$85\,\%. Diskrepanz = Ausgangspunkt \\
\addlinespace
II. Einordnung & 7 & Ist das Phänomen in der Literatur bekannt? & Goetz 2024, Zheng 2023. Pipeline: 916$\to$3 (KV Hessen). Lücke: Kein ``Governance-Pharmakovigilanz''-Konzept \\
\addlinespace
III. Kosten & 6 & Was kostet das System? Gibt es Präventionssoftware? & 0,9--1,7\,Mrd.\,EUR Folgekosten (konservativ). Keine Open-Source-Präventionssoftware \\
\addlinespace
IV. Rechtsstruktur & 5 & Wie funktioniert V2 genau? & V2 nahezu deterministisch nach Antrag. Kein Beratungsschutz. Container-Genealogie: 1 Norm $\to$ 8 Container \\
\addlinespace
V. Modellierung & 5 & Kann man das Verhalten formal modellieren? & Nash-Gleichgewicht = defensive Unterverordnung. 9 Normdefekte. \S\,12 SGB\,V unauflösbar \\
\addlinespace
VI. Verfassung & 2 & Verfassungsrechtliche Angriffsvektoren? & Art.\,2 Abs.\,2 GG (Patientenschutzpflicht) = stärkster Hebel. GPE = Niskanen-Bürokratie \\
\addlinespace
VII. Gemini & 2 & Wo sind Lücken? Externes Double-Check & GPE-Kosten nicht verifizierbar. Ribbat medial ``still'' \\
\addlinespace
VIII. Theorie & 5 Quellen & Welche Theorien fundieren die Befunde? & Freidson, Power, Lipsky, Herd/Moynihan, Bernstein. Alle bestätigen: Audit-Bürokratie funktioniert wie vorgesehen \\
\addlinespace
IX. Core & Dialog & Was ist die kanonische Wahrheit? & 7+2 CORE-Referenzpunkte (CORE~1--7 + CORE3: Schutzstufenmodell/V2a-V2b + CORE4: Definitionsdivergenz/Doppelfunktion). 4 Messebenen. V1$\neq$V2. Drei Schichten Regressbegriff \\
\bottomrule
\end{tabularx}
}

\textbf{Gesamtumfang:} 31 Recherche-Dateien, 6+ Session-Berichte, 3 Review-Sessions (13-Phasen-Review-Zyklus mit 9 parallelen Agenten), 108 manuelle PDF-Kommentare als Sichtprüfung. Geschätzte Gesamtarbeitszeit: ca. 60--80 Stunden (Mensch + KI kombiniert, April 2026).

\subsection*{Dokumentierte Korrekturen und Hypothesenrevisionen}

{\small
\begin{tabularx}{\textwidth}{X X X}
\toprule
\textbf{Korrektur} & \textbf{Anlass} & \textbf{Ergebnis} \\
\midrule
``Moosazadeh 2023'' & Literatur-Double-Check & Fehlzuordnung in Gemini-Quelle $\to$ korrigiert auf Zheng et al. 2023 \\
\addlinespace
Einzelfallprüfung ab GSG 1992 & Rechtshistorische Tiefenrecherche & Bereits GRG 1988 (Blüm) $\to$ Hypothese um 4 Jahre verschoben \\
\addlinespace
7--17 Mrd = attributierte Kosten & Copilot-Review & Systemrahmen, nicht Attribution $\to$ Dreistufenmodell eingeführt \\
\addlinespace
``125 GPs protestierten 2024'' (FR) & Faktencheck Frankreich & Gemini-Halluzination $\to$ entfernt, echte CNAM-Zahlen eingesetzt \\
\addlinespace
1\,\%-Quote = V1-only & Core-Etablierung & Korrigiert: V1+V2 zusammen, keine Verfahrenstrennung publiziert \\
\bottomrule
\end{tabularx}
}

\subsection*{Abhängigkeitsstruktur}

Die 31 Analysen bilden eine gerichtete Wissenskaskade in fünf Schichten: (1)~Empirie $\to$ (2)~Vertiefung $\to$ (3)~Synthese $\to$ (4)~Bewertung $\to$ (5)~Theorierahmen + Core. Jede Schicht baut auf den Befunden der vorherigen auf. An keiner Stelle wurde eine Schlussfolgerung ohne dokumentierte Quellengrundlage gezogen.

\textbf{Methodische Regel:} Etablierte Klassiker (Freidson, Power, Lipsky) waren bei Gemini Deep Research zuverlässig. Aktuelle politische Ereignisse mit spezifischen Akteuren (Frankreich: ``d\'elit statistique'') erforderten systematischen Double-Check --- Größenordnungen und Urheberschaften waren verzerrt.

\subsection*{Phase XIV: Multi-Agenten-Experiment (Cross-Source-Divergenz)}

In einer späteren Analysephase wurden 15 spezialisierte Agenten (6~Opus, 9~Sonnet) parallel eingesetzt, um die Quellen-Bubble der ``unter 1\,\%''-Zahl systematisch zu durchbrechen:

{\small
\begin{tabularx}{\textwidth}{l l X}
\toprule
\textbf{Agent} & \textbf{Modell} & \textbf{Aufgabe und Schlüsselbefund} \\
\midrule
Agent 1--3 & Sonnet & Datensammlung (25 Datenpunkte, 14 Domains), Impressum-Verifizierung, Tenor-Kategorisierung \\
\addlinespace
Agent A--B & Sonnet & Widerspruchs-Erklärungen (6/14 thematisieren den Widerspruch, alle als ``Wahrnehmungsfehler''); Naive Suche (``unter 1\,\%'' in ersten 5 Treffern) \\
\addlinespace
Agent C & Opus & Unabhängige V2-Zahlen-Suche $\to$ BMG GVSG: 47.000 V2-Verfahren bundesweit \\
\addlinespace
Prior-Art & Opus & Systematische Suche (20+ Anfragen) $\to$ Synthese (Definitionsdivergenz) nirgends publiziert \\
\addlinespace
EFP-1/2 & Sonnet/Opus & KVBW-Prüfthemen; Alternative Hypothese H2 $\to$ 35--40\,\% Erklärungskraft \\
\addlinespace
KK-1--4 & Sonnet & Krankenkassen-Website-Bewertung (Bias-Kontrolle: 2 Agenten, umgekehrte Reihenfolge) \\
\addlinespace
Evidenz & Opus & Suche nach These ``Evidenz in formaler Verkleidung'' $\to$ originär \\
\bottomrule
\end{tabularx}
}

\textbf{Zentrale methodische Erkenntnis:} Cross-Model-Divergenz (Claude vs. Copilot vs. Gemini) kontrolliert Modell-Bias, aber \textit{nicht} Quellen-Bias. Alle drei Modelle schöpften aus denselben Quellen (Ärzteblatt, Virchowbund). Erst Cross-Source-Divergenz (14 Domains, verschiedene Akteursgruppen) brachte die SpiFa-Bayern-Zahlen ans Licht. Für quantitative Befunde ist Multi-Quellen-Recherche zwingend.

\textbf{Zwei Prior-Art-Checks bestätigt:} Beide zentralen Synthesen (Definitionsdivergenz als Erklärung für Ribbat-Diskrepanz; Evidenz in formaler Verkleidung als Doppelfunktion des Systems) sind nach systematischer Recherche \textbf{originär} --- alle Einzelbausteine publiziert, Verbindung nirgends.

Vollständige Dokumentation: 14 RECHERCHE-Dateien in den Arbeitsmaterialien und METHODIK\_RECHERCHE\_PROTOKOLL.md.

\subsection*{Prompt-Archäologie: Quantitative Analyse der Mensch-KI-Interaktion}

Über den gesamten Analysezeitraum (07.--11.04.2026) wurden \textbf{208 Prompts} systematisch dokumentiert und klassifiziert: 127 menschliche Eingaben und 81 LLM-generierte Aufträge (an Sub-Agenten, Gemini, Copilot). Die Klassifikation erfolgte nach einer eigens entwickelten 5-Stufen-Taxonomie (Rohdaten $\to$ Themenfilter $\to$ Klassifikation $\to$ Aggregation $\to$ Statistik).

\textbf{Zentrale Befunde zur Mensch-KI-Forschungsdynamik:}

\begin{enumerate}[nosep]
  \item \textbf{Alle 6 Wendepunkte waren menschliche Startprompts.} Kein einziger Durchbruch wurde von der KI eigenständig ausgelöst. Die Wendepunkte entstanden durch kurze, oft beiläufige Fragen (Median: 10,5 Wörter --- kürzer als der Durchschnitt aller Prompts).

  \item \textbf{Bestätigungs-zu-Korrektur-Verhältnis: 1,25:1.} Bei unkritischer Mensch-KI-Interaktion wäre ein Verhältnis von 5:1 oder höher zu erwarten. LG korrigierte fast so oft wie er bestätigte --- Indikator für fehlenden Zustimmungs-Bias.

  \item \textbf{Korrekturrate stieg von 2,8\,\% (Tag~1) auf 15,8\,\% (Tag~5).} Der Mensch wurde im Verlauf kritischer, nicht zahmer. Die produktivsten Korrekturen (H104, H108, H111) lösten die CORE3- und CORE4-Paradigmenwechsel aus.

  \item \textbf{33,9\,\% der Prompts erzeugten ${\sim}$95\,\% der Kernerkenntnisse.} Startprompts (Hebelfaktor 2,9$\times$) und Korrekturen (2,7$\times$) waren die produktivsten Prompt-Typen. Bestätigungen und Nachfragen trugen kaum zu originären Befunden bei.

  \item \textbf{Proaktiv:Reaktiv = 1,74:1.} Der Mensch agierte 1,7$\times$ so oft proaktiv (neue Fragen, neue Richtungen) wie reaktiv (Bestätigung, Korrektur). Die Forschung war klar menschlich geführt, nicht KI-getrieben.

  \item \textbf{KI-Autonomie-Gradient:} Die KI war bei strukturierten Aufgaben zuverlässig (Tabellen, Recherche, Dokumentation). Bei synthetischen, interpretierenden Aufgaben (Bedeutung der 99,7\,\%, Definitonsartefakt, Bilanzbewertung) lag sie ohne menschliche Korrektur falsch oder unvollständig. Konservative Schätzung: Ohne menschliche Korrekturen wären mindestens 3 der 5 zentralen CORE-Referenzpunkte nicht in ihrer jetzigen Form entstanden.
\end{enumerate}

\noindent\textbf{Methodische Implikation:} Diese Analyse legt nahe, dass KI-gestützte Forschung nicht als ``KI forscht, Mensch prüft'' zu verstehen ist, sondern als ``Mensch fragt, KI sucht, Mensch korrigiert und synthetisiert''. Die entscheidende Wertschöpfung liegt in den Fragen und Korrekturen --- nicht in den Antworten. Die vollständige Prompt-Dokumentation (208 Prompts, 5 Aggregationsstufen) ist in den Arbeitsmaterialien verfügbar.

\newpage
\section*{Anhang C: Offene Forschungsfragen und eingereichte IFG-Anfragen}
\addcontentsline{toc}{section}{Anhang C: Offene Forschungsfragen und IFG-Anfragen}
\label{sec:appendix-ifg}

Die zentrale Empfehlung der Review-Zyklen dieser Studie lautete: ``Ein realer V2-Datensatz wäre der größte einzelne Qualitätssprung.'' Die Gesamtzahl der V2-Verfahren ist inzwischen bekannt (BMG GVSG-Referentenentwurf 2024: ca.\,47.000 bundesweit). Was fehlt, ist die Aufschlüsselung: nach V2a/V2b, nach Kassenart, nach Ausgang (Regress/Vergleich/Einstellung), nach Betragsklasse. Um die Datenlücke aktiv zu adressieren, wurden im April 2026 neun IFG-Anfragen (Informationsfreiheitsgesetz, \S\S\,1\,ff. IFG) zu zentralen offenen Forschungsfragen eingereicht. Die Anfragen sind öffentlich einsehbar auf FragDenStaat.de und auf um-bruch.org/presse/ifg-anfragen/.

\subsection*{Übersicht der Anfragen}

{\small
\begin{tabularx}{\textwidth}{l l X l}
\toprule
\textbf{\#} & \textbf{Empfänger} & \textbf{Forschungsfrage} & \textbf{Status} \\
\midrule
A1 & BMG & Bundesweite Prüfstatistiken \S\S\,106b/106d: Verfahrenszahlen, Ergebnisse, Dropout-Gründe, Volumen gefordert vs. durchgesetzt & Eingereicht \\
\addlinespace
A2 & GPE Rheinland-Pfalz & Regionale Prüfstatistiken (gleicher Fragenkatalog, regional zugespitzt) & Eingereicht \\
\addlinespace
A3 & BAS & Eurozahlen GPE: Personalkosten, Sachkosten, Verwaltungsbudgets, monofunktionale Stellen & Eingereicht \\
\addlinespace
A5 & GKV-Spitzenverband & Berechnungsgrundlage der ``500\,Mio.\,EUR Mehrausgaben''-Behauptung (Stellungnahme KHVVG, BT-Drs 20(14)220(6)) & Eingereicht \\
\addlinespace
A6 & Bundesrechnungshof & Sonderprüfungen der Wirtschaftlichkeitsprüfung seit 2005 --- falls keine: Begründung & Eingereicht \\
\addlinespace
A7 & SVR Gesundheit / BMG & Interne Berichte oder Beratungspapiere zur Effizienz der Wirtschaftlichkeitsprüfung nach 2014 & Eingereicht \\
\addlinespace
A8 & BMJ / Statistisches Bundesamt & Sozialgerichtsstatistik zu Wirtschaftlichkeitsprüfungsverfahren (Anzahl, Ausgang, Dauer) & Eingereicht \\
\addlinespace
A9 & BMG & Normwidersprüche und Regelungslücken: 10 dokumentierte strukturelle Inkonsistenzen im Normgefüge & Eingereicht \\
\addlinespace
A11 & KBV & Anwendung \S\,48 BMV-Ä (``Sonstiger Schaden'') in den letzten fünf Jahren & Eingereicht \\
\bottomrule
\end{tabularx}
}

\subsection*{Was die Antworten für diese Studie bedeuten würden}

\begin{itemize}[nosep]
  \item \textbf{A1 + A2:} Würden erstmals die Eskalationspipeline der Einzelfallprüfung (V2) mit realen Zahlen füllen. Die gesamte Spalte ``k.D.'' in den Pipeline-Tabellen (Abschnitt~2.4) könnte durch Daten ersetzt werden. Dies wäre der ``größte einzelne Qualitätssprung'' für die Studie.
  \item \textbf{A3:} Würde den Niskanen-Befund (Abschnitt~5.3) empirisch untermauern oder widerlegen. Erstmals wären die Kosten der Prüfinfrastruktur quantifizierbar.
  \item \textbf{A5:} Würde zeigen, ob die 500-Mio.-Behauptung des GKV-SV auf belastbaren Daten basiert oder eine politische Schätzung ist.
  \item \textbf{A6:} Die Antwort ``keine Sonderprüfung seit 2005'' wäre selbst ein Befund über die institutionelle Prüffreiheit des Systems.
  \item \textbf{A9:} Würde erzwingen, dass das BMG schriftlich zu konkreten, dokumentierten Schwachstellen im Normgefüge Stellung nimmt.
  \item \textbf{A11:} Würde den ``unsichtbaren Container'' \S\,48 BMV-Ä erstmals quantifizierbar machen.
\end{itemize}

\textbf{Antwortfrist:} 30 Tage nach Zugang (\S\,7 Abs.\,5 IFG). Erste Empfangsbestätigungen liegen vor (BMG: A1, A7, A9 --- bestätigt 10.04.2026). Sobald substanzielle Antworten eintreffen, werden sie in eine aktualisierte Version dieser Studie eingearbeitet.

\textbf{Transparenzhinweis:} Alle Anfragen sind öffentlich einsehbar. Auch Verweigerungen sind publizistisch verwertbar und werden dokumentiert --- eine Antwort ``diese Daten liegen uns nicht vor'' bestätigt die in dieser Studie diagnostizierte Transparenzlücke.


\newpage
\section*{Änderungshistorie}
\addcontentsline{toc}{section}{Änderungshistorie}

{\small\sloppy
\begin{itemize}[leftmargin=3em, itemsep=0.5em]
  \item \textbf{v0.12} (April 2026, 2026-04-11) --- \textbf{CORE4-Integration + Encoding + BibTeX.} Drei-Schichten-Tabelle des Regressbegriffs (Definitionsdivergenz, Faktor 470--500$\times$). Neuer Abschnitt ``Evidenz in formaler Verkleidung'' (Zwei Ebenen der Formalität, Doppelfunktion A/B/C, CME vs. V2). Einwand~5 neu: ``Ein Drittel des Systems ist rationale Steuerung''. Zertifikatsfehler-Qualifier (``überwiegend nicht messen'', V2b-Fokus vorgezogen). Abstract CORE4-konform umgeschrieben. Bayern SpiFa-Daten (13.332 Regresse, 334~EUR Durchschnitt). Appendix~B: Phase~XIV Multi-Agenten-Experiment (15 Agenten, Prior-Art-Checks). Appendix~C: A1-Framing korrigiert (47.000 bekannt, Aufschlüsselung fehlt). CORE-Referenzpunkte 7$\to$7+2. Encoding: Alle escaped Umlaute $\to$ echtes UTF-8. \texttt{amsmath}-Package ergänzt. \textbf{BibTeX-Literaturverzeichnis} mit 50 Einträgen (\texttt{ST-001\_references.bib}). 87~Seiten.

  \item \textbf{v0.11} (April 2026, 2026-04-11) --- \textbf{CORE3-Integration.} V2a/V2b-Differenzierung in 7 Abschnitten. V1-Schutzstufenmodell eingearbeitet (Anhörung, Beratung, Karenzzeit, Deckel). Zertifikatsfehler nur noch V2b in voller Schärfe. \S\,106/\S\,106b Sachfehler korrigiert. Kommunikationsasymmetrie-Absatz mit Poisson-Modell ($\lambda$=0,85, 57\,\%/Jahr). Bayern/BMG-Zahlen (47.000 V2-Verfahren). Summierung ``Kleines Risiko pro Fall, großes Risiko pro Arzt''. 99,7\,\%-Verwerfungsrate auf V1 beschränkt. Lesarten-Abschnitt komplett umgeschrieben. Executive Summary (ST-001~ES, v1.0) erstellt. 82~Seiten.

  \item \textbf{v0.10} (April 2026, 2026-04-11) --- \textbf{Review-Zyklus + Methodenteil + KNV.} Methodenteil von 12 auf $\sim$60 Zeilen erweitert (Forschungsdesign, Suchstrategie, Multi-Modell-Architektur, Ein-/Ausschlusskriterien, Bias-Management). Ribbat-Operationalisierung (Mapping-Tabelle, Falsifizierbarkeit). Gegeneinwände-Callouts (Vorwärtsverweise auf Einwände~1--3). KNV-Tabelle (14 Maßnahmen, Rang~1--14). Goodhart's/Campbell's Law. Governance-Pharmakovigilanz. Bidirektionale Fehlversorgung. KBV-Berufsmonitoring korrigiert (46,7\,\%/13.900). 3 Review-Agenten: Major Revision $\to$ knapp Minor Revision. 76~Seiten.

  \item \textbf{v0.9} (April 2026, 2026-04-11) --- LG-Sichtprüfung mit 108 PDF-Kommentaren eingearbeitet: ``Die Diagnose'' statt ``Worum es geht'', neue Fragestellung-Überschrift, 50:1-Framing ersetzt durch Ribbat 72\,\%-Befund, Container-Definition, Richtgrößen definiert, Stufenmodell gestrichen, Selektionsbias $\to$ Machtungleichgewicht mit SG Mainz 2022, ``Lesart C'' $\to$ ``System als Formfehler-Erkennung'', Hessendata-Verfassungsrecht (\mbox{1\,BvR\,1547/19}), Doppelprüfungsverbot (LSG BaWü L\,5\,KA\,1505/23), natürliches Experiment 2017 (WIdO-Daten). Appendix~B: Walk of Analysis (Methodenteil). V1/V2 konsequent getrennt. 66~Seiten.

  \item \textbf{v0.8} (April 2026, 2026-04-10) --- 13-Phasen-Review-Zyklus: 26+ Korrekturen (Tonalität, Advocacy-Sprache entschärft, Zertifikatsfehler als ``Interpretationslücke'' reformuliert); normativer Schadensbegriff (BSG B\,6\,KA\,5/09\,R) als neuer Abschnitt; \textbf{neue Synthese ``Drei Scanner, keine Qualitätsprüfung''}: Lesart~C (System als Formfehler-Optimierer), Informationsasymmetrie Prüfkatalog, offene Rechtsfrage automatisiertes Screening vs. \S\,106b-Verdacht. 59~Seiten.

  \item \textbf{v0.7} (April 2026, 2026-04-10) --- Zwei Eskalationspipeline-Tabellen (V1~Auffälligkeitsprüfung + V2~Einzelfallprüfung, je 6~Spalten + Kernaussagen); Vier-Messebenen-Kommentierung; 17-Zeilen-Vergleichstabelle Auffälligkeits- vs. Einzelfallprüfung (3~Gruppen: Verfahrens-, Schutz-, Risikoprofil); Tabelle Vier informelle Ausgleichs\-mechanismen (Optionen A--D); strukturelles Framing V1~=~probabilistisch / V2~=~nahezu deterministisch; Quellenbelege in Fußnoten: KV\,BaWü~2022 (\texttt{kvbawue.de}, 0,097\,\%, 21\,/\,21.500) und KV\,Hessen~2015 (\textit{Dt.\,Ärzteblatt} Art.~113630, 3\,/\,916 = 0,3\,\%\,$|$\,S0); Klarstellung dass die ,,unter 1\,\%``-Zahl für V1+V2 zusammen gilt (keine Verfahrenstrennung veröffentlicht); Version~0.7.

  \item \textbf{v0.6} (April 2026, 2026-04-09) --- Neue Abschnitte: Begriffsklärung der Regressmetriken (Vier Messebenen), Termination without Determination, Doppelter Zertifikatsfehler, Stufenmodell, Selektionsbias, Formfehler als Strategie, Praxisbesonderheiten. 52~Seiten.

  \item \textbf{v0.5} (April 2026) --- Niskanen-Analyse der GPE (6~Kriterien), Spieltheorie (Nash-Gleichgewicht, Schweigeprämium), verfassungsrechtliche Angriffsvektoren (Bestimmtheit, Doppelbeschwerde Arzt + Patient), Modell\-arzt-Hochrechnung mit Sensitivitätstabelle (direkte + anteilige Folgekosten), Verhältnisrechnung Schaden\,:\,Einnahme.

  \item \textbf{v0.4} (April 2026) --- Google-Research-Review: Freidson (Managerialismus, Professional Dominance, Third Logic), Power (Audit Society), Lipsky (Street-Level Bureaucracy), Herd/Moynihan (Administrative Burden), Bernstein (Transparency Paradox als Selbstkritik zu PP-003). Frankreich-Sektion: D\'elit statistique, MSO/MSAP, CNAM-Kampagnen 2023--2025, MG\,France, FMF.

  \item \textbf{v0.3} (April 2026) --- Gemini-Recherchen: GPE-Kosten\-befund (Niskanen), diametrale Manifestation US\,/\,DE (Überbehandlung vs. Unterversorgung).

  \item \textbf{v0.2} (März/April 2026) --- Copilot-Review~1 eingearbeitet: Sensitivitätsbox, Screening-vs.-Verdachtsprüfung, Eigentumsdogmatik, Public Choice, Drei Reformhebel, Vergleichsfolie Defensivmedizin, Glossar. 28~Seiten.

  \item \textbf{v0.1} (März 2026) --- Erstes Skelett mit allen 8~Teilen + Verfassungsrechtskapitel.
\end{itemize}
}

\newpage
\section*{KI-Transparenzhinweis}

\noindent Dieses Dokument wurde unter extensivem Einsatz von KI-Sprachmodellen (u.\,a. Anthropic Claude, Microsoft Copilot, Google Gemini) erarbeitet. Der Einsatz war notwendig, um diese Analyse mit begrenzten Ressourcen zu ermöglichen. Es wurden mehrere Kontroll- und Reviewzyklen durchgeführt~--- sowohl KI-gestützt als auch durch menschliche Redaktion.

\medskip\noindent KI-Sprachmodelle neigen zu Halluzinationen und können sachliche Fehler produzieren. Trotz sorgfältiger Prüfung sind Fehler nicht ausgeschlossen. Wir korrigieren dokumentierte Fehler gerne und veröffentlichen verbesserte Versionen.

\medskip\noindent Dieses Dokument stellt keine Rechtsberatung dar.

\medskip\noindent Hinweise auf Fehler bitte an: \href{mailto:hallo@um-bruch.org}{hallo@um-bruch.org}

\end{document}
